% Môn: Vật lý
% Lớp: 12
% Chương: Chương IV. Vật lí hạt nhân
% Bài: L12C4 Bài 21. Cấu trúc hạt nhân
% Số câu: 189
% App đọc file này trực tiếp — sửa rồi Commit trên GitHub, bấm Đồng bộ GitHub .tex

% L12C4 Bài 21. Cấu trúc hạt nhân

% ===== Câu 1 =====
% ID: L12C4B21-01-DS
% Mức: TH
\dangbt{Bài toán tổng hợp về cấu trúc và đặc trưng của hạt nhân}
\begin{ex}
% ID: L12C4B21-01-DS
% Mức: TH
Xét bài toán tổng hợp về cấu trúc và đặc trưng của hạt nhân (tình huống 1). Hãy xác định tính đúng, sai của bốn phát biểu sau.
\choiceTF
{\True Khi giải bài tổng hợp phải xác định đúng $A,Z,N$, đổi đơn vị khối lượng và dùng nhất quán $E=mc^2$ hoặc $1\,u\,c^2=931,5\,MeV$.}
{Có thể đồng nhất số khối với khối lượng hạt nhân tính theo đơn vị $u$ trong mọi phép tính chính xác.}
{\True Khi giải cần đọc đúng dữ kiện, dùng hệ đơn vị thống nhất và kiểm tra điều kiện áp dụng.}
{Có thể bỏ qua dữ kiện, đơn vị và ý nghĩa vật lí của kết quả.}
\loigiai{
\begin{itemchoice}
\item (Đúng) Khi giải bài tổng hợp phải xác định đúng $A,Z,N$, đổi đơn vị khối lượng và dùng nhất quán $E=mc^2$ hoặc $1\,u\,c^2=931,5\,MeV$. (Vì): Khi giải bài tổng hợp phải xác định đúng $A,Z,N$, đổi đơn vị khối lượng và dùng nhất quán $E=mc^2$ hoặc $1\,u\,c^2=931,5\,MeV$. b) Sai: Có thể đồng nhất số khối với khối lượng hạt nhân tính theo đơn vị $u$ trong mọi phép tính chính xác. c) Đúng: Khi giải cần đọc đúng dữ kiện, dùng hệ đơn vị thống nhất và kiểm tra điều kiện áp dụng. d) Sai: Có thể bỏ qua dữ kiện, đơn vị và ý nghĩa vật lí của kết quả. Kiến thức cốt lõi: Khi giải bài tổng hợp phải xác định đúng $A,Z,N$, đổi đơn vị khối lượng và dùng nhất quán $E=mc^2$ hoặc $1\,u\,c^2=931,5\,MeV$
\item (Sai) Có thể đồng nhất số khối với khối lượng hạt nhân tính theo đơn vị $u$ trong mọi phép tính chính xác.
\item (Đúng) Khi giải cần đọc đúng dữ kiện, dùng hệ đơn vị thống nhất và kiểm tra điều kiện áp dụng.
\item (Sai) Có thể bỏ qua dữ kiện, đơn vị và ý nghĩa vật lí của kết quả.
\end{itemchoice}
}
\end{ex}

% ===== Câu 2 =====
% ID: L12C4B21-01-TL
% Mức: TH
\begin{bt}
% ID: L12C4B21-01-TL
% Mức: TH
Trình bày kiến thức cốt lõi của bài toán tổng hợp về cấu trúc và đặc trưng của hạt nhân và nêu điều kiện áp dụng.
\loigiai{
Cần nêu được: Khi giải bài tổng hợp phải xác định đúng $A,Z,N$, đổi đơn vị khối lượng và dùng nhất quán $E=mc^2$ hoặc $1\,u\,c^2=931,5\,MeV$. Khi vận dụng phải xác định đúng dữ kiện, điều kiện áp dụng, hệ đơn vị và kiểm tra tính hợp lí của kết quả. Nhận định “Có thể đồng nhất số khối với khối lượng hạt nhân tính theo đơn vị $u$ trong mọi phép tính chính xác.” sai vì trái với kiến thức trên.
}
\end{bt}

% ===== Câu 3 =====
% ID: L12C4B21-01-TLN
% Mức: TH
\begin{ex}
% ID: L12C4B21-01-TLN
% Mức: TH
Trong bốn phát biểu sau về bài toán tổng hợp về cấu trúc và đặc trưng của hạt nhân, có bao nhiêu phát biểu đúng? Chỉ ghi một số nguyên. (1) Khi giải bài tổng hợp phải xác định đúng $A,Z,N$, đổi đơn vị khối lượng và dùng nhất quán $E=mc^2$ hoặc $1\,u\,c^2=931,5\,MeV$. (2) Có thể đồng nhất số khối với khối lượng hạt nhân tính theo đơn vị $u$ trong mọi phép tính chính xác. (3) Cần kiểm tra điều kiện áp dụng trước khi tính toán. (4) Có thể bỏ qua đơn vị của kết quả.
\shortans{2}
\loigiai{
Đối chiếu từng phát biểu với kiến thức cốt lõi: Khi giải bài tổng hợp phải xác định đúng $A,Z,N$, đổi đơn vị khối lượng và dùng nhất quán $E=mc^2$ hoặc $1\,u\,c^2=931,5\,MeV$. Có 2 phát biểu đúng.
}
\end{ex}

% ===== Câu 4 =====
% ID: L12C4B21-01-TN
% Mức: NB
\begin{ex}
% ID: L12C4B21-01-TN
% Mức: NB
Phát biểu nào sau đây đúng về bài toán tổng hợp về cấu trúc và đặc trưng của hạt nhân?
\choice
{\True Khi giải bài tổng hợp phải xác định đúng $A,Z,N$, đổi đơn vị khối lượng và dùng nhất quán $E=mc^2$ hoặc $1\,u\,c^2=931,5\,MeV$.}
{Có thể đồng nhất số khối với khối lượng hạt nhân tính theo đơn vị $u$ trong mọi phép tính chính xác.}
{Có thể bỏ qua điều kiện áp dụng và đơn vị trong mọi trường hợp.}
{Kết quả của bài toán không phụ thuộc dữ kiện đã cho.}
\loigiai{
Khi giải bài tổng hợp phải xác định đúng $A,Z,N$, đổi đơn vị khối lượng và dùng nhất quán $E=mc^2$ hoặc $1\,u\,c^2=931,5\,MeV$. Vì vậy phương án $A$ đúng; các phương án còn lại trái với kiến thức hoặc điều kiện áp dụng.
}
\end{ex}

% ===== Câu 5 =====
% ID: L12C4B21-02-DS
% Mức: TH
\begin{ex}
% ID: L12C4B21-02-DS
% Mức: TH
Xét bài toán tổng hợp về cấu trúc và đặc trưng của hạt nhân (tình huống 2). Hãy xác định tính đúng, sai của bốn phát biểu sau.
\choiceTF
{Có thể đồng nhất số khối với khối lượng hạt nhân tính theo đơn vị $u$ trong mọi phép tính chính xác.}
{\True Khi giải bài tổng hợp phải xác định đúng $A,Z,N$, đổi đơn vị khối lượng và dùng nhất quán $E=mc^2$ hoặc $1\,u\,c^2=931,5\,MeV$.}
{Có thể bỏ qua dữ kiện, đơn vị và ý nghĩa vật lí của kết quả.}
{\True Khi giải cần đọc đúng dữ kiện, dùng hệ đơn vị thống nhất và kiểm tra điều kiện áp dụng.}
\loigiai{
\begin{itemchoice}
\item (Sai) Có thể đồng nhất số khối với khối lượng hạt nhân tính theo đơn vị $u$ trong mọi phép tính chính xác. (Vì): Có thể đồng nhất số khối với khối lượng hạt nhân tính theo đơn vị $u$ trong mọi phép tính chính xác. b) Đúng: Khi giải bài tổng hợp phải xác định đúng $A,Z,N$, đổi đơn vị khối lượng và dùng nhất quán $E=mc^2$ hoặc $1\,u\,c^2=931,5\,MeV$. c) Sai: Có thể bỏ qua dữ kiện, đơn vị và ý nghĩa vật lí của kết quả. d) Đúng: Khi giải cần đọc đúng dữ kiện, dùng hệ đơn vị thống nhất và kiểm tra điều kiện áp dụng. Kiến thức cốt lõi: Khi giải bài tổng hợp phải xác định đúng $A,Z,N$, đổi đơn vị khối lượng và dùng nhất quán $E=mc^2$ hoặc $1\,u\,c^2=931,5\,MeV$
\item (Đúng) Khi giải bài tổng hợp phải xác định đúng $A,Z,N$, đổi đơn vị khối lượng và dùng nhất quán $E=mc^2$ hoặc $1\,u\,c^2=931,5\,MeV$.
\item (Sai) Có thể bỏ qua dữ kiện, đơn vị và ý nghĩa vật lí của kết quả.
\item (Đúng) Khi giải cần đọc đúng dữ kiện, dùng hệ đơn vị thống nhất và kiểm tra điều kiện áp dụng.
\end{itemchoice}
}
\end{ex}

% ===== Câu 6 =====
% ID: L12C4B21-02-TL
% Mức: TH
\begin{bt}
% ID: L12C4B21-02-TL
% Mức: TH
Giải thích vì sao nhận định sau đúng: “Khi giải bài tổng hợp phải xác định đúng $A,Z,N$, đổi đơn vị khối lượng và dùng nhất quán $E=mc^2$ hoặc $1\,u\,c^2=931,5\,MeV$.”
\loigiai{
Cần nêu được: Khi giải bài tổng hợp phải xác định đúng $A,Z,N$, đổi đơn vị khối lượng và dùng nhất quán $E=mc^2$ hoặc $1\,u\,c^2=931,5\,MeV$. Khi vận dụng phải xác định đúng dữ kiện, điều kiện áp dụng, hệ đơn vị và kiểm tra tính hợp lí của kết quả. Nhận định “Có thể đồng nhất số khối với khối lượng hạt nhân tính theo đơn vị $u$ trong mọi phép tính chính xác.” sai vì trái với kiến thức trên.
}
\end{bt}

% ===== Câu 7 =====
% ID: L12C4B21-02-TLN
% Mức: TH
\begin{ex}
% ID: L12C4B21-02-TLN
% Mức: TH
Trong bốn phát biểu sau về bài toán tổng hợp về cấu trúc và đặc trưng của hạt nhân, có bao nhiêu phát biểu đúng? Chỉ ghi một số nguyên. (1) Có thể đồng nhất số khối với khối lượng hạt nhân tính theo đơn vị $u$ trong mọi phép tính chính xác. (2) Khi giải bài tổng hợp phải xác định đúng $A,Z,N$, đổi đơn vị khối lượng và dùng nhất quán $E=mc^2$ hoặc $1\,u\,c^2=931,5\,MeV$. (3) Kết quả phải phù hợp ý nghĩa vật lí. (4) Mọi đại lượng đều có cùng bản chất.
\shortans{2}
\loigiai{
Đối chiếu từng phát biểu với kiến thức cốt lõi: Khi giải bài tổng hợp phải xác định đúng $A,Z,N$, đổi đơn vị khối lượng và dùng nhất quán $E=mc^2$ hoặc $1\,u\,c^2=931,5\,MeV$. Có 2 phát biểu đúng.
}
\end{ex}

% ===== Câu 8 =====
% ID: L12C4B21-02-TN
% Mức: NB
\begin{ex}
% ID: L12C4B21-02-TN
% Mức: NB
Nhận định nào phù hợp với kiến thức Vật lí về bài toán tổng hợp về cấu trúc và đặc trưng của hạt nhân?
\choice
{Có thể đồng nhất số khối với khối lượng hạt nhân tính theo đơn vị $u$ trong mọi phép tính chính xác.}
{\True Khi giải bài tổng hợp phải xác định đúng $A,Z,N$, đổi đơn vị khối lượng và dùng nhất quán $E=mc^2$ hoặc $1\,u\,c^2=931,5\,MeV$.}
{Có thể bỏ qua điều kiện áp dụng và đơn vị trong mọi trường hợp.}
{Kết quả của bài toán không phụ thuộc dữ kiện đã cho.}
\loigiai{
Khi giải bài tổng hợp phải xác định đúng $A,Z,N$, đổi đơn vị khối lượng và dùng nhất quán $E=mc^2$ hoặc $1\,u\,c^2=931,5\,MeV$. Vì vậy phương án $B$ đúng; các phương án còn lại trái với kiến thức hoặc điều kiện áp dụng.
}
\end{ex}

% ===== Câu 9 =====
% ID: L12C4B21-03-DS
% Mức: VD
\begin{ex}
% ID: L12C4B21-03-DS
% Mức: VD
Xét bài toán tổng hợp về cấu trúc và đặc trưng của hạt nhân (tình huống 3). Hãy xác định tính đúng, sai của bốn phát biểu sau.
\choiceTF
{\True Khi giải bài tổng hợp phải xác định đúng $A,Z,N$, đổi đơn vị khối lượng và dùng nhất quán $E=mc^2$ hoặc $1\,u\,c^2=931,5\,MeV$.}
{Có thể đồng nhất số khối với khối lượng hạt nhân tính theo đơn vị $u$ trong mọi phép tính chính xác.}
{\True Khi giải cần đọc đúng dữ kiện, dùng hệ đơn vị thống nhất và kiểm tra điều kiện áp dụng.}
{Có thể bỏ qua dữ kiện, đơn vị và ý nghĩa vật lí của kết quả.}
\loigiai{
\begin{itemchoice}
\item (Đúng) Khi giải bài tổng hợp phải xác định đúng $A,Z,N$, đổi đơn vị khối lượng và dùng nhất quán $E=mc^2$ hoặc $1\,u\,c^2=931,5\,MeV$. (Vì): Khi giải bài tổng hợp phải xác định đúng $A,Z,N$, đổi đơn vị khối lượng và dùng nhất quán $E=mc^2$ hoặc $1\,u\,c^2=931,5\,MeV$. b) Sai: Có thể đồng nhất số khối với khối lượng hạt nhân tính theo đơn vị $u$ trong mọi phép tính chính xác. c) Đúng: Khi giải cần đọc đúng dữ kiện, dùng hệ đơn vị thống nhất và kiểm tra điều kiện áp dụng. d) Sai: Có thể bỏ qua dữ kiện, đơn vị và ý nghĩa vật lí của kết quả. Kiến thức cốt lõi: Khi giải bài tổng hợp phải xác định đúng $A,Z,N$, đổi đơn vị khối lượng và dùng nhất quán $E=mc^2$ hoặc $1\,u\,c^2=931,5\,MeV$
\item (Sai) Có thể đồng nhất số khối với khối lượng hạt nhân tính theo đơn vị $u$ trong mọi phép tính chính xác.
\item (Đúng) Khi giải cần đọc đúng dữ kiện, dùng hệ đơn vị thống nhất và kiểm tra điều kiện áp dụng.
\item (Sai) Có thể bỏ qua dữ kiện, đơn vị và ý nghĩa vật lí của kết quả.
\end{itemchoice}
}
\end{ex}

% ===== Câu 10 =====
% ID: L12C4B21-03-TL
% Mức: VD
\begin{bt}
% ID: L12C4B21-03-TL
% Mức: VD
Phân tích sai lầm trong nhận định: “Có thể đồng nhất số khối với khối lượng hạt nhân tính theo đơn vị $u$ trong mọi phép tính chính xác.”
\loigiai{
Cần nêu được: Khi giải bài tổng hợp phải xác định đúng $A,Z,N$, đổi đơn vị khối lượng và dùng nhất quán $E=mc^2$ hoặc $1\,u\,c^2=931,5\,MeV$. Khi vận dụng phải xác định đúng dữ kiện, điều kiện áp dụng, hệ đơn vị và kiểm tra tính hợp lí của kết quả. Nhận định “Có thể đồng nhất số khối với khối lượng hạt nhân tính theo đơn vị $u$ trong mọi phép tính chính xác.” sai vì trái với kiến thức trên.
}
\end{bt}

% ===== Câu 11 =====
% ID: L12C4B21-03-TLN
% Mức: TH
\begin{ex}
% ID: L12C4B21-03-TLN
% Mức: TH
Trong bốn phát biểu sau về bài toán tổng hợp về cấu trúc và đặc trưng của hạt nhân, có bao nhiêu phát biểu đúng? Chỉ ghi một số nguyên. (1) Cần đổi các đại lượng về hệ đơn vị thống nhất. (2) Khi giải bài tổng hợp phải xác định đúng $A,Z,N$, đổi đơn vị khối lượng và dùng nhất quán $E=mc^2$ hoặc $1\,u\,c^2=931,5\,MeV$. (3) Có thể đồng nhất số khối với khối lượng hạt nhân tính theo đơn vị $u$ trong mọi phép tính chính xác. (4) Không cần kiểm tra dữ kiện.
\shortans{2}
\loigiai{
Đối chiếu từng phát biểu với kiến thức cốt lõi: Khi giải bài tổng hợp phải xác định đúng $A,Z,N$, đổi đơn vị khối lượng và dùng nhất quán $E=mc^2$ hoặc $1\,u\,c^2=931,5\,MeV$. Có 2 phát biểu đúng.
}
\end{ex}

% ===== Câu 12 =====
% ID: L12C4B21-03-TN
% Mức: NB
\begin{ex}
% ID: L12C4B21-03-TN
% Mức: NB
Chọn kết luận chính xác về bài toán tổng hợp về cấu trúc và đặc trưng của hạt nhân?
\choice
{Có thể đồng nhất số khối với khối lượng hạt nhân tính theo đơn vị $u$ trong mọi phép tính chính xác.}
{Có thể bỏ qua điều kiện áp dụng và đơn vị trong mọi trường hợp.}
{\True Khi giải bài tổng hợp phải xác định đúng $A,Z,N$, đổi đơn vị khối lượng và dùng nhất quán $E=mc^2$ hoặc $1\,u\,c^2=931,5\,MeV$.}
{Kết quả của bài toán không phụ thuộc dữ kiện đã cho.}
\loigiai{
Khi giải bài tổng hợp phải xác định đúng $A,Z,N$, đổi đơn vị khối lượng và dùng nhất quán $E=mc^2$ hoặc $1\,u\,c^2=931,5\,MeV$. Vì vậy phương án $C$ đúng; các phương án còn lại trái với kiến thức hoặc điều kiện áp dụng.
}
\end{ex}

% ===== Câu 13 =====
% ID: L12C4B21-04-DS
% Mức: VD
\begin{ex}
% ID: L12C4B21-04-DS
% Mức: VD
Xét bài toán tổng hợp về cấu trúc và đặc trưng của hạt nhân (tình huống 4). Hãy xác định tính đúng, sai của bốn phát biểu sau.
\choiceTF
{Có thể đồng nhất số khối với khối lượng hạt nhân tính theo đơn vị $u$ trong mọi phép tính chính xác.}
{\True Khi giải bài tổng hợp phải xác định đúng $A,Z,N$, đổi đơn vị khối lượng và dùng nhất quán $E=mc^2$ hoặc $1\,u\,c^2=931,5\,MeV$.}
{Có thể bỏ qua dữ kiện, đơn vị và ý nghĩa vật lí của kết quả.}
{\True Khi giải cần đọc đúng dữ kiện, dùng hệ đơn vị thống nhất và kiểm tra điều kiện áp dụng.}
\loigiai{
\begin{itemchoice}
\item (Sai) Có thể đồng nhất số khối với khối lượng hạt nhân tính theo đơn vị $u$ trong mọi phép tính chính xác. (Vì): Có thể đồng nhất số khối với khối lượng hạt nhân tính theo đơn vị $u$ trong mọi phép tính chính xác. b) Đúng: Khi giải bài tổng hợp phải xác định đúng $A,Z,N$, đổi đơn vị khối lượng và dùng nhất quán $E=mc^2$ hoặc $1\,u\,c^2=931,5\,MeV$. c) Sai: Có thể bỏ qua dữ kiện, đơn vị và ý nghĩa vật lí của kết quả. d) Đúng: Khi giải cần đọc đúng dữ kiện, dùng hệ đơn vị thống nhất và kiểm tra điều kiện áp dụng. Kiến thức cốt lõi: Khi giải bài tổng hợp phải xác định đúng $A,Z,N$, đổi đơn vị khối lượng và dùng nhất quán $E=mc^2$ hoặc $1\,u\,c^2=931,5\,MeV$
\item (Đúng) Khi giải bài tổng hợp phải xác định đúng $A,Z,N$, đổi đơn vị khối lượng và dùng nhất quán $E=mc^2$ hoặc $1\,u\,c^2=931,5\,MeV$.
\item (Sai) Có thể bỏ qua dữ kiện, đơn vị và ý nghĩa vật lí của kết quả.
\item (Đúng) Khi giải cần đọc đúng dữ kiện, dùng hệ đơn vị thống nhất và kiểm tra điều kiện áp dụng.
\end{itemchoice}
}
\end{ex}

% ===== Câu 14 =====
% ID: L12C4B21-04-TL
% Mức: VD
\begin{bt}
% ID: L12C4B21-04-TL
% Mức: VD
Đề xuất các bước giải một bài toán thuộc dạng bài toán tổng hợp về cấu trúc và đặc trưng của hạt nhân.
\loigiai{
Cần nêu được: Khi giải bài tổng hợp phải xác định đúng $A,Z,N$, đổi đơn vị khối lượng và dùng nhất quán $E=mc^2$ hoặc $1\,u\,c^2=931,5\,MeV$. Khi vận dụng phải xác định đúng dữ kiện, điều kiện áp dụng, hệ đơn vị và kiểm tra tính hợp lí của kết quả. Nhận định “Có thể đồng nhất số khối với khối lượng hạt nhân tính theo đơn vị $u$ trong mọi phép tính chính xác.” sai vì trái với kiến thức trên.
}
\end{bt}

% ===== Câu 15 =====
% ID: L12C4B21-04-TLN
% Mức: VD
\begin{ex}
% ID: L12C4B21-04-TLN
% Mức: VD
Trong bốn phát biểu sau về bài toán tổng hợp về cấu trúc và đặc trưng của hạt nhân, có bao nhiêu phát biểu đúng? Chỉ ghi một số nguyên. (1) Khi giải bài tổng hợp phải xác định đúng $A,Z,N$, đổi đơn vị khối lượng và dùng nhất quán $E=mc^2$ hoặc $1\,u\,c^2=931,5\,MeV$. (2) Cần đối chiếu kết quả với điều kiện bài toán. (3) Có thể đồng nhất số khối với khối lượng hạt nhân tính theo đơn vị $u$ trong mọi phép tính chính xác. (4) Mọi trường hợp đều cho cùng một kết quả.
\shortans{2}
\loigiai{
Đối chiếu từng phát biểu với kiến thức cốt lõi: Khi giải bài tổng hợp phải xác định đúng $A,Z,N$, đổi đơn vị khối lượng và dùng nhất quán $E=mc^2$ hoặc $1\,u\,c^2=931,5\,MeV$. Có 2 phát biểu đúng.
}
\end{ex}

% ===== Câu 16 =====
% ID: L12C4B21-04-TN
% Mức: TH
\begin{ex}
% ID: L12C4B21-04-TN
% Mức: TH
Nội dung nào mô tả đúng nhất về bài toán tổng hợp về cấu trúc và đặc trưng của hạt nhân?
\choice
{Có thể đồng nhất số khối với khối lượng hạt nhân tính theo đơn vị $u$ trong mọi phép tính chính xác.}
{Có thể bỏ qua điều kiện áp dụng và đơn vị trong mọi trường hợp.}
{Kết quả của bài toán không phụ thuộc dữ kiện đã cho.}
{\True Khi giải bài tổng hợp phải xác định đúng $A,Z,N$, đổi đơn vị khối lượng và dùng nhất quán $E=mc^2$ hoặc $1\,u\,c^2=931,5\,MeV$.}
\loigiai{
Khi giải bài tổng hợp phải xác định đúng $A,Z,N$, đổi đơn vị khối lượng và dùng nhất quán $E=mc^2$ hoặc $1\,u\,c^2=931,5\,MeV$. Vì vậy phương án $D$ đúng; các phương án còn lại trái với kiến thức hoặc điều kiện áp dụng.
}
\end{ex}

% ===== Câu 17 =====
% ID: L12C4B21-05-DS
% Mức: VD
\begin{ex}
% ID: L12C4B21-05-DS
% Mức: VD
Xét bài toán tổng hợp về cấu trúc và đặc trưng của hạt nhân (tình huống 5). Hãy xác định tính đúng, sai của bốn phát biểu sau.
\choiceTF
{\True Khi giải bài tổng hợp phải xác định đúng $A,Z,N$, đổi đơn vị khối lượng và dùng nhất quán $E=mc^2$ hoặc $1\,u\,c^2=931,5\,MeV$.}
{Có thể bỏ qua dữ kiện, đơn vị và ý nghĩa vật lí của kết quả.}
{Có thể đồng nhất số khối với khối lượng hạt nhân tính theo đơn vị $u$ trong mọi phép tính chính xác.}
{Có thể bỏ qua dữ kiện, đơn vị và ý nghĩa vật lí của kết quả.}
\loigiai{
\begin{itemchoice}
\item (Đúng) Khi giải bài tổng hợp phải xác định đúng $A,Z,N$, đổi đơn vị khối lượng và dùng nhất quán $E=mc^2$ hoặc $1\,u\,c^2=931,5\,MeV$. (Vì): Khi giải bài tổng hợp phải xác định đúng $A,Z,N$, đổi đơn vị khối lượng và dùng nhất quán $E=mc^2$ hoặc $1\,u\,c^2=931,5\,MeV$. b) Sai: Có thể bỏ qua dữ kiện, đơn vị và ý nghĩa vật lí của kết quả. c) Sai: Có thể đồng nhất số khối với khối lượng hạt nhân tính theo đơn vị $u$ trong mọi phép tính chính xác. d) Sai: Có thể bỏ qua dữ kiện, đơn vị và ý nghĩa vật lí của kết quả. Kiến thức cốt lõi: Khi giải bài tổng hợp phải xác định đúng $A,Z,N$, đổi đơn vị khối lượng và dùng nhất quán $E=mc^2$ hoặc $1\,u\,c^2=931,5\,MeV$
\item (Sai) Có thể bỏ qua dữ kiện, đơn vị và ý nghĩa vật lí của kết quả.
\item (Sai) Có thể đồng nhất số khối với khối lượng hạt nhân tính theo đơn vị $u$ trong mọi phép tính chính xác.
\item (Sai) Có thể bỏ qua dữ kiện, đơn vị và ý nghĩa vật lí của kết quả.
\end{itemchoice}
}
\end{ex}

% ===== Câu 18 =====
% ID: L12C4B21-05-TL
% Mức: VD
\begin{bt}
% ID: L12C4B21-05-TL
% Mức: VD
Nêu một tình huống thực tiễn liên quan đến bài toán tổng hợp về cấu trúc và đặc trưng của hạt nhân và giải thích bằng kiến thức Vật lí.
\loigiai{
Cần nêu được: Khi giải bài tổng hợp phải xác định đúng $A,Z,N$, đổi đơn vị khối lượng và dùng nhất quán $E=mc^2$ hoặc $1\,u\,c^2=931,5\,MeV$. Khi vận dụng phải xác định đúng dữ kiện, điều kiện áp dụng, hệ đơn vị và kiểm tra tính hợp lí của kết quả. Nhận định “Có thể đồng nhất số khối với khối lượng hạt nhân tính theo đơn vị $u$ trong mọi phép tính chính xác.” sai vì trái với kiến thức trên.
}
\end{bt}

% ===== Câu 19 =====
% ID: L12C4B21-05-TLN
% Mức: VD
\begin{ex}
% ID: L12C4B21-05-TLN
% Mức: VD
Trong bốn phát biểu sau về bài toán tổng hợp về cấu trúc và đặc trưng của hạt nhân, có bao nhiêu phát biểu đúng? Chỉ ghi một số nguyên. (1) Có thể đồng nhất số khối với khối lượng hạt nhân tính theo đơn vị $u$ trong mọi phép tính chính xác. (2) Cần đọc đúng dữ kiện và giả thiết. (3) Khi giải bài tổng hợp phải xác định đúng $A,Z,N$, đổi đơn vị khối lượng và dùng nhất quán $E=mc^2$ hoặc $1\,u\,c^2=931,5\,MeV$. (4) Có thể thay số trước khi chọn công thức.
\shortans{2}
\loigiai{
Đối chiếu từng phát biểu với kiến thức cốt lõi: Khi giải bài tổng hợp phải xác định đúng $A,Z,N$, đổi đơn vị khối lượng và dùng nhất quán $E=mc^2$ hoặc $1\,u\,c^2=931,5\,MeV$. Có 2 phát biểu đúng.
}
\end{ex}

% ===== Câu 20 =====
% ID: L12C4B21-05-TN
% Mức: TH
\begin{ex}
% ID: L12C4B21-05-TN
% Mức: TH
Khi phân tích, phát biểu nào đúng về bài toán tổng hợp về cấu trúc và đặc trưng của hạt nhân?
\choice
{\True Khi giải bài tổng hợp phải xác định đúng $A,Z,N$, đổi đơn vị khối lượng và dùng nhất quán $E=mc^2$ hoặc $1\,u\,c^2=931,5\,MeV$.}
{Có thể đồng nhất số khối với khối lượng hạt nhân tính theo đơn vị $u$ trong mọi phép tính chính xác.}
{Có thể bỏ qua điều kiện áp dụng và đơn vị trong mọi trường hợp.}
{Kết quả của bài toán không phụ thuộc dữ kiện đã cho.}
\loigiai{
Khi giải bài tổng hợp phải xác định đúng $A,Z,N$, đổi đơn vị khối lượng và dùng nhất quán $E=mc^2$ hoặc $1\,u\,c^2=931,5\,MeV$. Vì vậy phương án $A$ đúng; các phương án còn lại trái với kiến thức hoặc điều kiện áp dụng.
}
\end{ex}

% ===== Câu 21 =====
% ID: L12C4B21-06-DS
% Mức: TH
\dangbt{Nhận biết cấu tạo hạt nhân, nuclon, số khối và điện tích hạt nhân}
\begin{ex}
% ID: L12C4B21-06-DS
% Mức: TH
Xét nhận biết cấu tạo hạt nhân, nuclon, số khối và điện tích hạt nhân (tình huống 1). Hãy xác định tính đúng, sai của bốn phát biểu sau.
\choiceTF
{\True Hạt nhân gồm proton và neutron; số khối $A=Z+N$, điện tích hạt nhân bằng $+Ze$.}
{Số khối $A$ luôn bằng số neutron $N$.}
{\True Khi giải cần đọc đúng dữ kiện, dùng hệ đơn vị thống nhất và kiểm tra điều kiện áp dụng.}
{Có thể bỏ qua dữ kiện, đơn vị và ý nghĩa vật lí của kết quả.}
\loigiai{
\begin{itemchoice}
\item (Đúng) Hạt nhân gồm proton và neutron; số khối $A=Z+N$, điện tích hạt nhân bằng $+Ze$. (Vì): Hạt nhân gồm proton và neutron; số khối $A=Z+N$, điện tích hạt nhân bằng $+Ze$. b) Sai: Số khối $A$ luôn bằng số neutron $N$. c) Đúng: Khi giải cần đọc đúng dữ kiện, dùng hệ đơn vị thống nhất và kiểm tra điều kiện áp dụng. d) Sai: Có thể bỏ qua dữ kiện, đơn vị và ý nghĩa vật lí của kết quả. Kiến thức cốt lõi: Hạt nhân gồm proton và neutron; số khối $A=Z+N$, điện tích hạt nhân bằng $+Ze$
\item (Sai) Số khối $A$ luôn bằng số neutron $N$.
\item (Đúng) Khi giải cần đọc đúng dữ kiện, dùng hệ đơn vị thống nhất và kiểm tra điều kiện áp dụng.
\item (Sai) Có thể bỏ qua dữ kiện, đơn vị và ý nghĩa vật lí của kết quả.
\end{itemchoice}
}
\end{ex}

% ===== Câu 22 =====
% ID: L12C4B21-06-TL
% Mức: VDC
\dangbt{Bài toán tổng hợp về cấu trúc và đặc trưng của hạt nhân}
\begin{bt}
% ID: L12C4B21-06-TL
% Mức: VDC
So sánh nhận định đúng “Khi giải bài tổng hợp phải xác định đúng $A,Z,N$, đổi đơn vị khối lượng và dùng nhất quán $E=mc^2$ hoặc $1\,u\,c^2=931,5\,MeV$.” với nhận định sai “Có thể đồng nhất số khối với khối lượng hạt nhân tính theo đơn vị $u$ trong mọi phép tính chính xác.”; chỉ rõ căn cứ Vật lí.
\loigiai{
Cần nêu được: Khi giải bài tổng hợp phải xác định đúng $A,Z,N$, đổi đơn vị khối lượng và dùng nhất quán $E=mc^2$ hoặc $1\,u\,c^2=931,5\,MeV$. Khi vận dụng phải xác định đúng dữ kiện, điều kiện áp dụng, hệ đơn vị và kiểm tra tính hợp lí của kết quả. Nhận định “Có thể đồng nhất số khối với khối lượng hạt nhân tính theo đơn vị $u$ trong mọi phép tính chính xác.” sai vì trái với kiến thức trên.
}
\end{bt}

% ===== Câu 23 =====
% ID: L12C4B21-06-TLN
% Mức: VDC
\begin{ex}
% ID: L12C4B21-06-TLN
% Mức: VDC
Trong bốn phát biểu sau về bài toán tổng hợp về cấu trúc và đặc trưng của hạt nhân, có bao nhiêu phát biểu đúng? Chỉ ghi một số nguyên. (1) Kết quả cần có ý nghĩa vật lí. (2) Có thể đồng nhất số khối với khối lượng hạt nhân tính theo đơn vị $u$ trong mọi phép tính chính xác. (3) Phải dùng đúng định luật cho tình huống. (4) Khi giải bài tổng hợp phải xác định đúng $A,Z,N$, đổi đơn vị khối lượng và dùng nhất quán $E=mc^2$ hoặc $1\,u\,c^2=931,5\,MeV$.
\shortans{3}
\loigiai{
Đối chiếu từng phát biểu với kiến thức cốt lõi: Khi giải bài tổng hợp phải xác định đúng $A,Z,N$, đổi đơn vị khối lượng và dùng nhất quán $E=mc^2$ hoặc $1\,u\,c^2=931,5\,MeV$. Có 3 phát biểu đúng.
}
\end{ex}

% ===== Câu 24 =====
% ID: L12C4B21-06-TN
% Mức: TH
\begin{ex}
% ID: L12C4B21-06-TN
% Mức: TH
Kết luận nào của học sinh là đúng về bài toán tổng hợp về cấu trúc và đặc trưng của hạt nhân?
\choice
{Có thể đồng nhất số khối với khối lượng hạt nhân tính theo đơn vị $u$ trong mọi phép tính chính xác.}
{\True Khi giải bài tổng hợp phải xác định đúng $A,Z,N$, đổi đơn vị khối lượng và dùng nhất quán $E=mc^2$ hoặc $1\,u\,c^2=931,5\,MeV$.}
{Có thể bỏ qua điều kiện áp dụng và đơn vị trong mọi trường hợp.}
{Kết quả của bài toán không phụ thuộc dữ kiện đã cho.}
\loigiai{
Khi giải bài tổng hợp phải xác định đúng $A,Z,N$, đổi đơn vị khối lượng và dùng nhất quán $E=mc^2$ hoặc $1\,u\,c^2=931,5\,MeV$. Vì vậy phương án $B$ đúng; các phương án còn lại trái với kiến thức hoặc điều kiện áp dụng.
}
\end{ex}

% ===== Câu 25 =====
% ID: L12C4B21-07-DS
% Mức: TH
\dangbt{Nhận biết cấu tạo hạt nhân, nuclon, số khối và điện tích hạt nhân}
\begin{ex}
% ID: L12C4B21-07-DS
% Mức: TH
Xét nhận biết cấu tạo hạt nhân, nuclon, số khối và điện tích hạt nhân (tình huống 2). Hãy xác định tính đúng, sai của bốn phát biểu sau.
\choiceTF
{Số khối $A$ luôn bằng số neutron $N$.}
{\True Hạt nhân gồm proton và neutron; số khối $A=Z+N$, điện tích hạt nhân bằng $+Ze$.}
{Có thể bỏ qua dữ kiện, đơn vị và ý nghĩa vật lí của kết quả.}
{\True Khi giải cần đọc đúng dữ kiện, dùng hệ đơn vị thống nhất và kiểm tra điều kiện áp dụng.}
\loigiai{
\begin{itemchoice}
\item (Sai) Số khối $A$ luôn bằng số neutron $N$. (Vì): ố khối $A$ luôn bằng số neutron $N$. b) Đúng: Hạt nhân gồm proton và neutron; số khối $A=Z+N$, điện tích hạt nhân bằng $+Ze$. c) Sai: Có thể bỏ qua dữ kiện, đơn vị và ý nghĩa vật lí của kết quả. d) Đúng: Khi giải cần đọc đúng dữ kiện, dùng hệ đơn vị thống nhất và kiểm tra điều kiện áp dụng. Kiến thức cốt lõi: Hạt nhân gồm proton và neutron; số khối $A=Z+N$, điện tích hạt nhân bằng $+Ze$
\item (Đúng) Hạt nhân gồm proton và neutron; số khối $A=Z+N$, điện tích hạt nhân bằng $+Ze$.
\item (Sai) Có thể bỏ qua dữ kiện, đơn vị và ý nghĩa vật lí của kết quả.
\item (Đúng) Khi giải cần đọc đúng dữ kiện, dùng hệ đơn vị thống nhất và kiểm tra điều kiện áp dụng.
\end{itemchoice}
}
\end{ex}

% ===== Câu 26 =====
% ID: L12C4B21-07-TL
% Mức: TH
\begin{bt}
% ID: L12C4B21-07-TL
% Mức: TH
Trình bày kiến thức cốt lõi của nhận biết cấu tạo hạt nhân, nuclon, số khối và điện tích hạt nhân và nêu điều kiện áp dụng.
\loigiai{
Cần nêu được: Hạt nhân gồm proton và neutron; số khối $A=Z+N$, điện tích hạt nhân bằng $+Ze$. Khi vận dụng phải xác định đúng dữ kiện, điều kiện áp dụng, hệ đơn vị và kiểm tra tính hợp lí của kết quả. Nhận định “Số khối $A$ luôn bằng số neutron $N$.” sai vì trái với kiến thức trên.
}
\end{bt}

% ===== Câu 27 =====
% ID: L12C4B21-07-TLN
% Mức: TH
\begin{ex}
% ID: L12C4B21-07-TLN
% Mức: TH
Trong bốn phát biểu sau về nhận biết cấu tạo hạt nhân, nuclon, số khối và điện tích hạt nhân, có bao nhiêu phát biểu đúng? Chỉ ghi một số nguyên. (1) Hạt nhân gồm proton và neutron; số khối $A=Z+N$, điện tích hạt nhân bằng $+Ze$. (2) Số khối $A$ luôn bằng số neutron $N$. (3) Cần kiểm tra điều kiện áp dụng trước khi tính toán. (4) Có thể bỏ qua đơn vị của kết quả.
\shortans{2}
\loigiai{
Đối chiếu từng phát biểu với kiến thức cốt lõi: Hạt nhân gồm proton và neutron; số khối $A=Z+N$, điện tích hạt nhân bằng $+Ze$. Có 2 phát biểu đúng.
}
\end{ex}

% ===== Câu 28 =====
% ID: L12C4B21-07-TN
% Mức: TH
\dangbt{Bài toán tổng hợp về cấu trúc và đặc trưng của hạt nhân}
\begin{ex}
% ID: L12C4B21-07-TN
% Mức: TH
Chọn phát biểu không trái với kiến thức về bài toán tổng hợp về cấu trúc và đặc trưng của hạt nhân?
\choice
{Có thể đồng nhất số khối với khối lượng hạt nhân tính theo đơn vị $u$ trong mọi phép tính chính xác.}
{Có thể bỏ qua điều kiện áp dụng và đơn vị trong mọi trường hợp.}
{\True Khi giải bài tổng hợp phải xác định đúng $A,Z,N$, đổi đơn vị khối lượng và dùng nhất quán $E=mc^2$ hoặc $1\,u\,c^2=931,5\,MeV$.}
{Kết quả của bài toán không phụ thuộc dữ kiện đã cho.}
\loigiai{
Khi giải bài tổng hợp phải xác định đúng $A,Z,N$, đổi đơn vị khối lượng và dùng nhất quán $E=mc^2$ hoặc $1\,u\,c^2=931,5\,MeV$. Vì vậy phương án $C$ đúng; các phương án còn lại trái với kiến thức hoặc điều kiện áp dụng.
}
\end{ex}

% ===== Câu 29 =====
% ID: L12C4B21-08-DS
% Mức: VD
\dangbt{Nhận biết cấu tạo hạt nhân, nuclon, số khối và điện tích hạt nhân}
\begin{ex}
% ID: L12C4B21-08-DS
% Mức: VD
Xét nhận biết cấu tạo hạt nhân, nuclon, số khối và điện tích hạt nhân (tình huống 3). Hãy xác định tính đúng, sai của bốn phát biểu sau.
\choiceTF
{\True Hạt nhân gồm proton và neutron; số khối $A=Z+N$, điện tích hạt nhân bằng $+Ze$.}
{Số khối $A$ luôn bằng số neutron $N$.}
{\True Khi giải cần đọc đúng dữ kiện, dùng hệ đơn vị thống nhất và kiểm tra điều kiện áp dụng.}
{Có thể bỏ qua dữ kiện, đơn vị và ý nghĩa vật lí của kết quả.}
\loigiai{
\begin{itemchoice}
\item (Đúng) Hạt nhân gồm proton và neutron; số khối $A=Z+N$, điện tích hạt nhân bằng $+Ze$. (Vì): Hạt nhân gồm proton và neutron; số khối $A=Z+N$, điện tích hạt nhân bằng $+Ze$. b) Sai: Số khối $A$ luôn bằng số neutron $N$. c) Đúng: Khi giải cần đọc đúng dữ kiện, dùng hệ đơn vị thống nhất và kiểm tra điều kiện áp dụng. d) Sai: Có thể bỏ qua dữ kiện, đơn vị và ý nghĩa vật lí của kết quả. Kiến thức cốt lõi: Hạt nhân gồm proton và neutron; số khối $A=Z+N$, điện tích hạt nhân bằng $+Ze$
\item (Sai) Số khối $A$ luôn bằng số neutron $N$.
\item (Đúng) Khi giải cần đọc đúng dữ kiện, dùng hệ đơn vị thống nhất và kiểm tra điều kiện áp dụng.
\item (Sai) Có thể bỏ qua dữ kiện, đơn vị và ý nghĩa vật lí của kết quả.
\end{itemchoice}
}
\end{ex}

% ===== Câu 30 =====
% ID: L12C4B21-08-TL
% Mức: TH
\begin{bt}
% ID: L12C4B21-08-TL
% Mức: TH
Giải thích vì sao nhận định sau đúng: “Hạt nhân gồm proton và neutron; số khối $A=Z+N$, điện tích hạt nhân bằng $+Ze$.”
\loigiai{
Cần nêu được: Hạt nhân gồm proton và neutron; số khối $A=Z+N$, điện tích hạt nhân bằng $+Ze$. Khi vận dụng phải xác định đúng dữ kiện, điều kiện áp dụng, hệ đơn vị và kiểm tra tính hợp lí của kết quả. Nhận định “Số khối $A$ luôn bằng số neutron $N$.” sai vì trái với kiến thức trên.
}
\end{bt}

% ===== Câu 31 =====
% ID: L12C4B21-08-TLN
% Mức: TH
\begin{ex}
% ID: L12C4B21-08-TLN
% Mức: TH
Trong bốn phát biểu sau về nhận biết cấu tạo hạt nhân, nuclon, số khối và điện tích hạt nhân, có bao nhiêu phát biểu đúng? Chỉ ghi một số nguyên. (1) Số khối $A$ luôn bằng số neutron $N$. (2) Hạt nhân gồm proton và neutron; số khối $A=Z+N$, điện tích hạt nhân bằng $+Ze$. (3) Kết quả phải phù hợp ý nghĩa vật lí. (4) Mọi đại lượng đều có cùng bản chất.
\shortans{2}
\loigiai{
Đối chiếu từng phát biểu với kiến thức cốt lõi: Hạt nhân gồm proton và neutron; số khối $A=Z+N$, điện tích hạt nhân bằng $+Ze$. Có 2 phát biểu đúng.
}
\end{ex}

% ===== Câu 32 =====
% ID: L12C4B21-08-TN
% Mức: VD
\dangbt{Bài toán tổng hợp về cấu trúc và đặc trưng của hạt nhân}
\begin{ex}
% ID: L12C4B21-08-TN
% Mức: VD
Cơ sở Vật lí đúng để giải bài là gì về bài toán tổng hợp về cấu trúc và đặc trưng của hạt nhân?
\choice
{Có thể đồng nhất số khối với khối lượng hạt nhân tính theo đơn vị $u$ trong mọi phép tính chính xác.}
{Có thể bỏ qua điều kiện áp dụng và đơn vị trong mọi trường hợp.}
{Kết quả của bài toán không phụ thuộc dữ kiện đã cho.}
{\True Khi giải bài tổng hợp phải xác định đúng $A,Z,N$, đổi đơn vị khối lượng và dùng nhất quán $E=mc^2$ hoặc $1\,u\,c^2=931,5\,MeV$.}
\loigiai{
Khi giải bài tổng hợp phải xác định đúng $A,Z,N$, đổi đơn vị khối lượng và dùng nhất quán $E=mc^2$ hoặc $1\,u\,c^2=931,5\,MeV$. Vì vậy phương án $D$ đúng; các phương án còn lại trái với kiến thức hoặc điều kiện áp dụng.
}
\end{ex}

% ===== Câu 33 =====
% ID: L12C4B21-09-DS
% Mức: VD
\dangbt{Nhận biết cấu tạo hạt nhân, nuclon, số khối và điện tích hạt nhân}
\begin{ex}
% ID: L12C4B21-09-DS
% Mức: VD
Xét nhận biết cấu tạo hạt nhân, nuclon, số khối và điện tích hạt nhân (tình huống 4). Hãy xác định tính đúng, sai của bốn phát biểu sau.
\choiceTF
{Số khối $A$ luôn bằng số neutron $N$.}
{\True Hạt nhân gồm proton và neutron; số khối $A=Z+N$, điện tích hạt nhân bằng $+Ze$.}
{Có thể bỏ qua dữ kiện, đơn vị và ý nghĩa vật lí của kết quả.}
{\True Khi giải cần đọc đúng dữ kiện, dùng hệ đơn vị thống nhất và kiểm tra điều kiện áp dụng.}
\loigiai{
\begin{itemchoice}
\item (Sai) Số khối $A$ luôn bằng số neutron $N$. (Vì): ố khối $A$ luôn bằng số neutron $N$. b) Đúng: Hạt nhân gồm proton và neutron; số khối $A=Z+N$, điện tích hạt nhân bằng $+Ze$. c) Sai: Có thể bỏ qua dữ kiện, đơn vị và ý nghĩa vật lí của kết quả. d) Đúng: Khi giải cần đọc đúng dữ kiện, dùng hệ đơn vị thống nhất và kiểm tra điều kiện áp dụng. Kiến thức cốt lõi: Hạt nhân gồm proton và neutron; số khối $A=Z+N$, điện tích hạt nhân bằng $+Ze$
\item (Đúng) Hạt nhân gồm proton và neutron; số khối $A=Z+N$, điện tích hạt nhân bằng $+Ze$.
\item (Sai) Có thể bỏ qua dữ kiện, đơn vị và ý nghĩa vật lí của kết quả.
\item (Đúng) Khi giải cần đọc đúng dữ kiện, dùng hệ đơn vị thống nhất và kiểm tra điều kiện áp dụng.
\end{itemchoice}
}
\end{ex}

% ===== Câu 34 =====
% ID: L12C4B21-09-TL
% Mức: VD
\begin{bt}
% ID: L12C4B21-09-TL
% Mức: VD
Phân tích sai lầm trong nhận định: “Số khối $A$ luôn bằng số neutron $N$.”
\loigiai{
Cần nêu được: Hạt nhân gồm proton và neutron; số khối $A=Z+N$, điện tích hạt nhân bằng $+Ze$. Khi vận dụng phải xác định đúng dữ kiện, điều kiện áp dụng, hệ đơn vị và kiểm tra tính hợp lí của kết quả. Nhận định “Số khối $A$ luôn bằng số neutron $N$.” sai vì trái với kiến thức trên.
}
\end{bt}

% ===== Câu 35 =====
% ID: L12C4B21-09-TLN
% Mức: TH
\begin{ex}
% ID: L12C4B21-09-TLN
% Mức: TH
Trong bốn phát biểu sau về nhận biết cấu tạo hạt nhân, nuclon, số khối và điện tích hạt nhân, có bao nhiêu phát biểu đúng? Chỉ ghi một số nguyên. (1) Cần đổi các đại lượng về hệ đơn vị thống nhất. (2) Hạt nhân gồm proton và neutron; số khối $A=Z+N$, điện tích hạt nhân bằng $+Ze$. (3) Số khối $A$ luôn bằng số neutron $N$. (4) Không cần kiểm tra dữ kiện.
\shortans{2}
\loigiai{
Đối chiếu từng phát biểu với kiến thức cốt lõi: Hạt nhân gồm proton và neutron; số khối $A=Z+N$, điện tích hạt nhân bằng $+Ze$. Có 2 phát biểu đúng.
}
\end{ex}

% ===== Câu 36 =====
% ID: L12C4B21-09-TN
% Mức: VD
\dangbt{Bài toán tổng hợp về cấu trúc và đặc trưng của hạt nhân}
\begin{ex}
% ID: L12C4B21-09-TN
% Mức: VD
Trong một bài thực hành, nhận xét nào đúng về bài toán tổng hợp về cấu trúc và đặc trưng của hạt nhân?
\choice
{\True Khi giải bài tổng hợp phải xác định đúng $A,Z,N$, đổi đơn vị khối lượng và dùng nhất quán $E=mc^2$ hoặc $1\,u\,c^2=931,5\,MeV$.}
{Có thể đồng nhất số khối với khối lượng hạt nhân tính theo đơn vị $u$ trong mọi phép tính chính xác.}
{Có thể bỏ qua điều kiện áp dụng và đơn vị trong mọi trường hợp.}
{Kết quả của bài toán không phụ thuộc dữ kiện đã cho.}
\loigiai{
Khi giải bài tổng hợp phải xác định đúng $A,Z,N$, đổi đơn vị khối lượng và dùng nhất quán $E=mc^2$ hoặc $1\,u\,c^2=931,5\,MeV$. Vì vậy phương án $A$ đúng; các phương án còn lại trái với kiến thức hoặc điều kiện áp dụng.
}
\end{ex}

% ===== Câu 37 =====
% ID: L12C4B21-10-DS
% Mức: VD
\dangbt{Nhận biết cấu tạo hạt nhân, nuclon, số khối và điện tích hạt nhân}
\begin{ex}
% ID: L12C4B21-10-DS
% Mức: VD
Xét nhận biết cấu tạo hạt nhân, nuclon, số khối và điện tích hạt nhân (tình huống 5). Hãy xác định tính đúng, sai của bốn phát biểu sau.
\choiceTF
{\True Hạt nhân gồm proton và neutron; số khối $A=Z+N$, điện tích hạt nhân bằng $+Ze$.}
{Có thể bỏ qua dữ kiện, đơn vị và ý nghĩa vật lí của kết quả.}
{Số khối $A$ luôn bằng số neutron $N$.}
{Có thể bỏ qua dữ kiện, đơn vị và ý nghĩa vật lí của kết quả.}
\loigiai{
\begin{itemchoice}
\item (Đúng) Hạt nhân gồm proton và neutron; số khối $A=Z+N$, điện tích hạt nhân bằng $+Ze$. (Vì): Hạt nhân gồm proton và neutron; số khối $A=Z+N$, điện tích hạt nhân bằng $+Ze$. b) Sai: Có thể bỏ qua dữ kiện, đơn vị và ý nghĩa vật lí của kết quả. c) Sai: Số khối $A$ luôn bằng số neutron $N$. d) Sai: Có thể bỏ qua dữ kiện, đơn vị và ý nghĩa vật lí của kết quả. Kiến thức cốt lõi: Hạt nhân gồm proton và neutron; số khối $A=Z+N$, điện tích hạt nhân bằng $+Ze$
\item (Sai) Có thể bỏ qua dữ kiện, đơn vị và ý nghĩa vật lí của kết quả.
\item (Sai) Số khối $A$ luôn bằng số neutron $N$.
\item (Sai) Có thể bỏ qua dữ kiện, đơn vị và ý nghĩa vật lí của kết quả.
\end{itemchoice}
}
\end{ex}

% ===== Câu 38 =====
% ID: L12C4B21-10-TL
% Mức: VD
\begin{bt}
% ID: L12C4B21-10-TL
% Mức: VD
Đề xuất các bước giải một bài toán thuộc dạng nhận biết cấu tạo hạt nhân, nuclon, số khối và điện tích hạt nhân.
\loigiai{
Cần nêu được: Hạt nhân gồm proton và neutron; số khối $A=Z+N$, điện tích hạt nhân bằng $+Ze$. Khi vận dụng phải xác định đúng dữ kiện, điều kiện áp dụng, hệ đơn vị và kiểm tra tính hợp lí của kết quả. Nhận định “Số khối $A$ luôn bằng số neutron $N$.” sai vì trái với kiến thức trên.
}
\end{bt}

% ===== Câu 39 =====
% ID: L12C4B21-10-TLN
% Mức: VD
\begin{ex}
% ID: L12C4B21-10-TLN
% Mức: VD
Trong bốn phát biểu sau về nhận biết cấu tạo hạt nhân, nuclon, số khối và điện tích hạt nhân, có bao nhiêu phát biểu đúng? Chỉ ghi một số nguyên. (1) Hạt nhân gồm proton và neutron; số khối $A=Z+N$, điện tích hạt nhân bằng $+Ze$. (2) Cần đối chiếu kết quả với điều kiện bài toán. (3) Số khối $A$ luôn bằng số neutron $N$. (4) Mọi trường hợp đều cho cùng một kết quả.
\shortans{2}
\loigiai{
Đối chiếu từng phát biểu với kiến thức cốt lõi: Hạt nhân gồm proton và neutron; số khối $A=Z+N$, điện tích hạt nhân bằng $+Ze$. Có 2 phát biểu đúng.
}
\end{ex}

% ===== Câu 40 =====
% ID: L12C4B21-10-TN
% Mức: VD
\dangbt{Bài toán tổng hợp về cấu trúc và đặc trưng của hạt nhân}
\begin{ex}
% ID: L12C4B21-10-TN
% Mức: VD
Kết luận đúng khi vận dụng là về bài toán tổng hợp về cấu trúc và đặc trưng của hạt nhân?
\choice
{Có thể đồng nhất số khối với khối lượng hạt nhân tính theo đơn vị $u$ trong mọi phép tính chính xác.}
{\True Khi giải bài tổng hợp phải xác định đúng $A,Z,N$, đổi đơn vị khối lượng và dùng nhất quán $E=mc^2$ hoặc $1\,u\,c^2=931,5\,MeV$.}
{Có thể bỏ qua điều kiện áp dụng và đơn vị trong mọi trường hợp.}
{Kết quả của bài toán không phụ thuộc dữ kiện đã cho.}
\loigiai{
Khi giải bài tổng hợp phải xác định đúng $A,Z,N$, đổi đơn vị khối lượng và dùng nhất quán $E=mc^2$ hoặc $1\,u\,c^2=931,5\,MeV$. Vì vậy phương án $B$ đúng; các phương án còn lại trái với kiến thức hoặc điều kiện áp dụng.
}
\end{ex}

% ===== Câu 41 =====
% ID: L12C4B21-11-DS
% Mức: TH
\dangbt{Viết kí hiệu hạt nhân và xác định số proton, neutron, electron}
\begin{ex}
% ID: L12C4B21-11-DS
% Mức: TH
Xét viết kí hiệu hạt nhân và xác định số proton, neutron, electron (tình huống 1). Hãy xác định tính đúng, sai của bốn phát biểu sau.
\choiceTF
{\True Với hạt nhân $^{A}_{Z}X$, số proton là $Z$, số neutron là $A-Z$; nguyên tử trung hòa có $Z$ electron.}
{Trong kí hiệu $^{A}_{Z}X$, số neutron bằng $Z$.}
{\True Khi giải cần đọc đúng dữ kiện, dùng hệ đơn vị thống nhất và kiểm tra điều kiện áp dụng.}
{Có thể bỏ qua dữ kiện, đơn vị và ý nghĩa vật lí của kết quả.}
\loigiai{
\begin{itemchoice}
\item (Đúng) Với hạt nhân $^{A}_{Z}X$, số proton là $Z$, số neutron là $A-Z$; nguyên tử trung hòa có $Z$ electron. (Vì): Với hạt nhân $^{A}_{Z}X$, số proton là $Z$, số neutron là $A-Z$; nguyên tử trung hòa có $Z$ electron. b) Sai: Trong kí hiệu $^{A}_{Z}X$, số neutron bằng $Z$. c) Đúng: Khi giải cần đọc đúng dữ kiện, dùng hệ đơn vị thống nhất và kiểm tra điều kiện áp dụng. d) Sai: Có thể bỏ qua dữ kiện, đơn vị và ý nghĩa vật lí của kết quả. Kiến thức cốt lõi: Với hạt nhân $^{A}_{Z}X$, số proton là $Z$, số neutron là $A-Z$; nguyên tử trung hòa có $Z$ electron
\item (Sai) Trong kí hiệu $^{A}_{Z}X$, số neutron bằng $Z$.
\item (Đúng) Khi giải cần đọc đúng dữ kiện, dùng hệ đơn vị thống nhất và kiểm tra điều kiện áp dụng.
\item (Sai) Có thể bỏ qua dữ kiện, đơn vị và ý nghĩa vật lí của kết quả.
\end{itemchoice}
}
\end{ex}

% ===== Câu 42 =====
% ID: L12C4B21-11-TL
% Mức: VD
\dangbt{Nhận biết cấu tạo hạt nhân, nuclon, số khối và điện tích hạt nhân}
\begin{bt}
% ID: L12C4B21-11-TL
% Mức: VD
Nêu một tình huống thực tiễn liên quan đến nhận biết cấu tạo hạt nhân, nuclon, số khối và điện tích hạt nhân và giải thích bằng kiến thức Vật lí.
\loigiai{
Cần nêu được: Hạt nhân gồm proton và neutron; số khối $A=Z+N$, điện tích hạt nhân bằng $+Ze$. Khi vận dụng phải xác định đúng dữ kiện, điều kiện áp dụng, hệ đơn vị và kiểm tra tính hợp lí của kết quả. Nhận định “Số khối $A$ luôn bằng số neutron $N$.” sai vì trái với kiến thức trên.
}
\end{bt}

% ===== Câu 43 =====
% ID: L12C4B21-11-TLN
% Mức: VD
\begin{ex}
% ID: L12C4B21-11-TLN
% Mức: VD
Trong bốn phát biểu sau về nhận biết cấu tạo hạt nhân, nuclon, số khối và điện tích hạt nhân, có bao nhiêu phát biểu đúng? Chỉ ghi một số nguyên. (1) Số khối $A$ luôn bằng số neutron $N$. (2) Cần đọc đúng dữ kiện và giả thiết. (3) Hạt nhân gồm proton và neutron; số khối $A=Z+N$, điện tích hạt nhân bằng $+Ze$. (4) Có thể thay số trước khi chọn công thức.
\shortans{2}
\loigiai{
Đối chiếu từng phát biểu với kiến thức cốt lõi: Hạt nhân gồm proton và neutron; số khối $A=Z+N$, điện tích hạt nhân bằng $+Ze$. Có 2 phát biểu đúng.
}
\end{ex}

% ===== Câu 44 =====
% ID: L12C4B21-11-TN
% Mức: NB
\begin{ex}
% ID: L12C4B21-11-TN
% Mức: NB
Phát biểu nào sau đây đúng về nhận biết cấu tạo hạt nhân, nuclon, số khối và điện tích hạt nhân?
\choice
{\True Hạt nhân gồm proton và neutron; số khối $A=Z+N$, điện tích hạt nhân bằng $+Ze$.}
{Số khối $A$ luôn bằng số neutron $N$.}
{Có thể bỏ qua điều kiện áp dụng và đơn vị trong mọi trường hợp.}
{Kết quả của bài toán không phụ thuộc dữ kiện đã cho.}
\loigiai{
Hạt nhân gồm proton và neutron; số khối $A=Z+N$, điện tích hạt nhân bằng $+Ze$. Vì vậy phương án $A$ đúng; các phương án còn lại trái với kiến thức hoặc điều kiện áp dụng.
}
\end{ex}

% ===== Câu 45 =====
% ID: L12C4B21-12-DS
% Mức: TH
\dangbt{Viết kí hiệu hạt nhân và xác định số proton, neutron, electron}
\begin{ex}
% ID: L12C4B21-12-DS
% Mức: TH
Xét viết kí hiệu hạt nhân và xác định số proton, neutron, electron (tình huống 2). Hãy xác định tính đúng, sai của bốn phát biểu sau.
\choiceTF
{Trong kí hiệu $^{A}_{Z}X$, số neutron bằng $Z$.}
{\True Với hạt nhân $^{A}_{Z}X$, số proton là $Z$, số neutron là $A-Z$; nguyên tử trung hòa có $Z$ electron.}
{Có thể bỏ qua dữ kiện, đơn vị và ý nghĩa vật lí của kết quả.}
{\True Khi giải cần đọc đúng dữ kiện, dùng hệ đơn vị thống nhất và kiểm tra điều kiện áp dụng.}
\loigiai{
\begin{itemchoice}
\item (Sai) Trong kí hiệu $^{A}_{Z}X$, số neutron bằng $Z$. (Vì): Trong kí hiệu $^{A}_{Z}X$, số neutron bằng $Z$. b) Đúng: Với hạt nhân $^{A}_{Z}X$, số proton là $Z$, số neutron là $A-Z$; nguyên tử trung hòa có $Z$ electron. c) Sai: Có thể bỏ qua dữ kiện, đơn vị và ý nghĩa vật lí của kết quả. d) Đúng: Khi giải cần đọc đúng dữ kiện, dùng hệ đơn vị thống nhất và kiểm tra điều kiện áp dụng. Kiến thức cốt lõi: Với hạt nhân $^{A}_{Z}X$, số proton là $Z$, số neutron là $A-Z$; nguyên tử trung hòa có $Z$ electron
\item (Đúng) Với hạt nhân $^{A}_{Z}X$, số proton là $Z$, số neutron là $A-Z$; nguyên tử trung hòa có $Z$ electron.
\item (Sai) Có thể bỏ qua dữ kiện, đơn vị và ý nghĩa vật lí của kết quả.
\item (Đúng) Khi giải cần đọc đúng dữ kiện, dùng hệ đơn vị thống nhất và kiểm tra điều kiện áp dụng.
\end{itemchoice}
}
\end{ex}

% ===== Câu 46 =====
% ID: L12C4B21-12-TL
% Mức: VDC
\dangbt{Nhận biết cấu tạo hạt nhân, nuclon, số khối và điện tích hạt nhân}
\begin{bt}
% ID: L12C4B21-12-TL
% Mức: VDC
So sánh nhận định đúng “Hạt nhân gồm proton và neutron; số khối $A=Z+N$, điện tích hạt nhân bằng $+Ze$.” với nhận định sai “Số khối $A$ luôn bằng số neutron $N$.”; chỉ rõ căn cứ Vật lí.
\loigiai{
Cần nêu được: Hạt nhân gồm proton và neutron; số khối $A=Z+N$, điện tích hạt nhân bằng $+Ze$. Khi vận dụng phải xác định đúng dữ kiện, điều kiện áp dụng, hệ đơn vị và kiểm tra tính hợp lí của kết quả. Nhận định “Số khối $A$ luôn bằng số neutron $N$.” sai vì trái với kiến thức trên.
}
\end{bt}

% ===== Câu 47 =====
% ID: L12C4B21-12-TLN
% Mức: VDC
\begin{ex}
% ID: L12C4B21-12-TLN
% Mức: VDC
Trong bốn phát biểu sau về nhận biết cấu tạo hạt nhân, nuclon, số khối và điện tích hạt nhân, có bao nhiêu phát biểu đúng? Chỉ ghi một số nguyên. (1) Kết quả cần có ý nghĩa vật lí. (2) Số khối $A$ luôn bằng số neutron $N$. (3) Phải dùng đúng định luật cho tình huống. (4) Hạt nhân gồm proton và neutron; số khối $A=Z+N$, điện tích hạt nhân bằng $+Ze$.
\shortans{3}
\loigiai{
Đối chiếu từng phát biểu với kiến thức cốt lõi: Hạt nhân gồm proton và neutron; số khối $A=Z+N$, điện tích hạt nhân bằng $+Ze$. Có 3 phát biểu đúng.
}
\end{ex}

% ===== Câu 48 =====
% ID: L12C4B21-12-TN
% Mức: NB
\begin{ex}
% ID: L12C4B21-12-TN
% Mức: NB
Nhận định nào phù hợp với kiến thức Vật lí về nhận biết cấu tạo hạt nhân, nuclon, số khối và điện tích hạt nhân?
\choice
{Số khối $A$ luôn bằng số neutron $N$.}
{\True Hạt nhân gồm proton và neutron; số khối $A=Z+N$, điện tích hạt nhân bằng $+Ze$.}
{Có thể bỏ qua điều kiện áp dụng và đơn vị trong mọi trường hợp.}
{Kết quả của bài toán không phụ thuộc dữ kiện đã cho.}
\loigiai{
Hạt nhân gồm proton và neutron; số khối $A=Z+N$, điện tích hạt nhân bằng $+Ze$. Vì vậy phương án $B$ đúng; các phương án còn lại trái với kiến thức hoặc điều kiện áp dụng.
}
\end{ex}

% ===== Câu 49 =====
% ID: L12C4B21-13-DS
% Mức: VD
\dangbt{Viết kí hiệu hạt nhân và xác định số proton, neutron, electron}
\begin{ex}
% ID: L12C4B21-13-DS
% Mức: VD
Xét viết kí hiệu hạt nhân và xác định số proton, neutron, electron (tình huống 3). Hãy xác định tính đúng, sai của bốn phát biểu sau.
\choiceTF
{\True Với hạt nhân $^{A}_{Z}X$, số proton là $Z$, số neutron là $A-Z$; nguyên tử trung hòa có $Z$ electron.}
{Trong kí hiệu $^{A}_{Z}X$, số neutron bằng $Z$.}
{\True Khi giải cần đọc đúng dữ kiện, dùng hệ đơn vị thống nhất và kiểm tra điều kiện áp dụng.}
{Có thể bỏ qua dữ kiện, đơn vị và ý nghĩa vật lí của kết quả.}
\loigiai{
\begin{itemchoice}
\item (Đúng) Với hạt nhân $^{A}_{Z}X$, số proton là $Z$, số neutron là $A-Z$; nguyên tử trung hòa có $Z$ electron. (Vì): Với hạt nhân $^{A}_{Z}X$, số proton là $Z$, số neutron là $A-Z$; nguyên tử trung hòa có $Z$ electron. b) Sai: Trong kí hiệu $^{A}_{Z}X$, số neutron bằng $Z$. c) Đúng: Khi giải cần đọc đúng dữ kiện, dùng hệ đơn vị thống nhất và kiểm tra điều kiện áp dụng. d) Sai: Có thể bỏ qua dữ kiện, đơn vị và ý nghĩa vật lí của kết quả. Kiến thức cốt lõi: Với hạt nhân $^{A}_{Z}X$, số proton là $Z$, số neutron là $A-Z$; nguyên tử trung hòa có $Z$ electron
\item (Sai) Trong kí hiệu $^{A}_{Z}X$, số neutron bằng $Z$.
\item (Đúng) Khi giải cần đọc đúng dữ kiện, dùng hệ đơn vị thống nhất và kiểm tra điều kiện áp dụng.
\item (Sai) Có thể bỏ qua dữ kiện, đơn vị và ý nghĩa vật lí của kết quả.
\end{itemchoice}
}
\end{ex}

% ===== Câu 50 =====
% ID: L12C4B21-13-TL
% Mức: TH
\begin{bt}
% ID: L12C4B21-13-TL
% Mức: TH
Trình bày kiến thức cốt lõi của viết kí hiệu hạt nhân và xác định số proton, neutron, electron và nêu điều kiện áp dụng.
\loigiai{
Cần nêu được: Với hạt nhân $^{A}_{Z}X$, số proton là $Z$, số neutron là $A-Z$; nguyên tử trung hòa có $Z$ electron. Khi vận dụng phải xác định đúng dữ kiện, điều kiện áp dụng, hệ đơn vị và kiểm tra tính hợp lí của kết quả. Nhận định “Trong kí hiệu $^{A}_{Z}X$, số neutron bằng $Z$.” sai vì trái với kiến thức trên.
}
\end{bt}

% ===== Câu 51 =====
% ID: L12C4B21-13-TLN
% Mức: TH
\begin{ex}
% ID: L12C4B21-13-TLN
% Mức: TH
Trong bốn phát biểu sau về viết kí hiệu hạt nhân và xác định số proton, neutron, electron, có bao nhiêu phát biểu đúng? Chỉ ghi một số nguyên. (1) Với hạt nhân $^{A}_{Z}X$, số proton là $Z$, số neutron là $A-Z$; nguyên tử trung hòa có $Z$ electron. (2) Trong kí hiệu $^{A}_{Z}X$, số neutron bằng $Z$. (3) Cần kiểm tra điều kiện áp dụng trước khi tính toán. (4) Có thể bỏ qua đơn vị của kết quả.
\shortans{2}
\loigiai{
Đối chiếu từng phát biểu với kiến thức cốt lõi: Với hạt nhân $^{A}_{Z}X$, số proton là $Z$, số neutron là $A-Z$; nguyên tử trung hòa có $Z$ electron. Có 2 phát biểu đúng.
}
\end{ex}

% ===== Câu 52 =====
% ID: L12C4B21-13-TN
% Mức: NB
\dangbt{Nhận biết cấu tạo hạt nhân, nuclon, số khối và điện tích hạt nhân}
\begin{ex}
% ID: L12C4B21-13-TN
% Mức: NB
Chọn kết luận chính xác về nhận biết cấu tạo hạt nhân, nuclon, số khối và điện tích hạt nhân?
\choice
{Số khối $A$ luôn bằng số neutron $N$.}
{Có thể bỏ qua điều kiện áp dụng và đơn vị trong mọi trường hợp.}
{\True Hạt nhân gồm proton và neutron; số khối $A=Z+N$, điện tích hạt nhân bằng $+Ze$.}
{Kết quả của bài toán không phụ thuộc dữ kiện đã cho.}
\loigiai{
Hạt nhân gồm proton và neutron; số khối $A=Z+N$, điện tích hạt nhân bằng $+Ze$. Vì vậy phương án $C$ đúng; các phương án còn lại trái với kiến thức hoặc điều kiện áp dụng.
}
\end{ex}

% ===== Câu 53 =====
% ID: L12C4B21-14-DS
% Mức: VD
\dangbt{Viết kí hiệu hạt nhân và xác định số proton, neutron, electron}
\begin{ex}
% ID: L12C4B21-14-DS
% Mức: VD
Xét viết kí hiệu hạt nhân và xác định số proton, neutron, electron (tình huống 4). Hãy xác định tính đúng, sai của bốn phát biểu sau.
\choiceTF
{Trong kí hiệu $^{A}_{Z}X$, số neutron bằng $Z$.}
{\True Với hạt nhân $^{A}_{Z}X$, số proton là $Z$, số neutron là $A-Z$; nguyên tử trung hòa có $Z$ electron.}
{Có thể bỏ qua dữ kiện, đơn vị và ý nghĩa vật lí của kết quả.}
{\True Khi giải cần đọc đúng dữ kiện, dùng hệ đơn vị thống nhất và kiểm tra điều kiện áp dụng.}
\loigiai{
\begin{itemchoice}
\item (Sai) Trong kí hiệu $^{A}_{Z}X$, số neutron bằng $Z$. (Vì): Trong kí hiệu $^{A}_{Z}X$, số neutron bằng $Z$. b) Đúng: Với hạt nhân $^{A}_{Z}X$, số proton là $Z$, số neutron là $A-Z$; nguyên tử trung hòa có $Z$ electron. c) Sai: Có thể bỏ qua dữ kiện, đơn vị và ý nghĩa vật lí của kết quả. d) Đúng: Khi giải cần đọc đúng dữ kiện, dùng hệ đơn vị thống nhất và kiểm tra điều kiện áp dụng. Kiến thức cốt lõi: Với hạt nhân $^{A}_{Z}X$, số proton là $Z$, số neutron là $A-Z$; nguyên tử trung hòa có $Z$ electron
\item (Đúng) Với hạt nhân $^{A}_{Z}X$, số proton là $Z$, số neutron là $A-Z$; nguyên tử trung hòa có $Z$ electron.
\item (Sai) Có thể bỏ qua dữ kiện, đơn vị và ý nghĩa vật lí của kết quả.
\item (Đúng) Khi giải cần đọc đúng dữ kiện, dùng hệ đơn vị thống nhất và kiểm tra điều kiện áp dụng.
\end{itemchoice}
}
\end{ex}

% ===== Câu 54 =====
% ID: L12C4B21-14-TL
% Mức: TH
\begin{bt}
% ID: L12C4B21-14-TL
% Mức: TH
Giải thích vì sao nhận định sau đúng: “Với hạt nhân $^{A}_{Z}X$, số proton là $Z$, số neutron là $A-Z$; nguyên tử trung hòa có $Z$ electron.”
\loigiai{
Cần nêu được: Với hạt nhân $^{A}_{Z}X$, số proton là $Z$, số neutron là $A-Z$; nguyên tử trung hòa có $Z$ electron. Khi vận dụng phải xác định đúng dữ kiện, điều kiện áp dụng, hệ đơn vị và kiểm tra tính hợp lí của kết quả. Nhận định “Trong kí hiệu $^{A}_{Z}X$, số neutron bằng $Z$.” sai vì trái với kiến thức trên.
}
\end{bt}

% ===== Câu 55 =====
% ID: L12C4B21-14-TLN
% Mức: TH
\begin{ex}
% ID: L12C4B21-14-TLN
% Mức: TH
Trong bốn phát biểu sau về viết kí hiệu hạt nhân và xác định số proton, neutron, electron, có bao nhiêu phát biểu đúng? Chỉ ghi một số nguyên. (1) Trong kí hiệu $^{A}_{Z}X$, số neutron bằng $Z$. (2) Với hạt nhân $^{A}_{Z}X$, số proton là $Z$, số neutron là $A-Z$; nguyên tử trung hòa có $Z$ electron. (3) Kết quả phải phù hợp ý nghĩa vật lí. (4) Mọi đại lượng đều có cùng bản chất.
\shortans{2}
\loigiai{
Đối chiếu từng phát biểu với kiến thức cốt lõi: Với hạt nhân $^{A}_{Z}X$, số proton là $Z$, số neutron là $A-Z$; nguyên tử trung hòa có $Z$ electron. Có 2 phát biểu đúng.
}
\end{ex}

% ===== Câu 56 =====
% ID: L12C4B21-14-TN
% Mức: TH
\dangbt{Nhận biết cấu tạo hạt nhân, nuclon, số khối và điện tích hạt nhân}
\begin{ex}
% ID: L12C4B21-14-TN
% Mức: TH
Nội dung nào mô tả đúng nhất về nhận biết cấu tạo hạt nhân, nuclon, số khối và điện tích hạt nhân?
\choice
{Số khối $A$ luôn bằng số neutron $N$.}
{Có thể bỏ qua điều kiện áp dụng và đơn vị trong mọi trường hợp.}
{Kết quả của bài toán không phụ thuộc dữ kiện đã cho.}
{\True Hạt nhân gồm proton và neutron; số khối $A=Z+N$, điện tích hạt nhân bằng $+Ze$.}
\loigiai{
Hạt nhân gồm proton và neutron; số khối $A=Z+N$, điện tích hạt nhân bằng $+Ze$. Vì vậy phương án $D$ đúng; các phương án còn lại trái với kiến thức hoặc điều kiện áp dụng.
}
\end{ex}

% ===== Câu 57 =====
% ID: L12C4B21-15-DS
% Mức: VD
\dangbt{Viết kí hiệu hạt nhân và xác định số proton, neutron, electron}
\begin{ex}
% ID: L12C4B21-15-DS
% Mức: VD
Xét viết kí hiệu hạt nhân và xác định số proton, neutron, electron (tình huống 5). Hãy xác định tính đúng, sai của bốn phát biểu sau.
\choiceTF
{\True Với hạt nhân $^{A}_{Z}X$, số proton là $Z$, số neutron là $A-Z$; nguyên tử trung hòa có $Z$ electron.}
{Có thể bỏ qua dữ kiện, đơn vị và ý nghĩa vật lí của kết quả.}
{Trong kí hiệu $^{A}_{Z}X$, số neutron bằng $Z$.}
{Có thể bỏ qua dữ kiện, đơn vị và ý nghĩa vật lí của kết quả.}
\loigiai{
\begin{itemchoice}
\item (Đúng) Với hạt nhân $^{A}_{Z}X$, số proton là $Z$, số neutron là $A-Z$; nguyên tử trung hòa có $Z$ electron. (Vì): Với hạt nhân $^{A}_{Z}X$, số proton là $Z$, số neutron là $A-Z$; nguyên tử trung hòa có $Z$ electron. b) Sai: Có thể bỏ qua dữ kiện, đơn vị và ý nghĩa vật lí của kết quả. c) Sai: Trong kí hiệu $^{A}_{Z}X$, số neutron bằng $Z$. d) Sai: Có thể bỏ qua dữ kiện, đơn vị và ý nghĩa vật lí của kết quả. Kiến thức cốt lõi: Với hạt nhân $^{A}_{Z}X$, số proton là $Z$, số neutron là $A-Z$; nguyên tử trung hòa có $Z$ electron
\item (Sai) Có thể bỏ qua dữ kiện, đơn vị và ý nghĩa vật lí của kết quả.
\item (Sai) Trong kí hiệu $^{A}_{Z}X$, số neutron bằng $Z$.
\item (Sai) Có thể bỏ qua dữ kiện, đơn vị và ý nghĩa vật lí của kết quả.
\end{itemchoice}
}
\end{ex}

% ===== Câu 58 =====
% ID: L12C4B21-15-TL
% Mức: VD
\begin{bt}
% ID: L12C4B21-15-TL
% Mức: VD
Phân tích sai lầm trong nhận định: “Trong kí hiệu $^{A}_{Z}X$, số neutron bằng $Z$.”
\loigiai{
Cần nêu được: Với hạt nhân $^{A}_{Z}X$, số proton là $Z$, số neutron là $A-Z$; nguyên tử trung hòa có $Z$ electron. Khi vận dụng phải xác định đúng dữ kiện, điều kiện áp dụng, hệ đơn vị và kiểm tra tính hợp lí của kết quả. Nhận định “Trong kí hiệu $^{A}_{Z}X$, số neutron bằng $Z$.” sai vì trái với kiến thức trên.
}
\end{bt}

% ===== Câu 59 =====
% ID: L12C4B21-15-TLN
% Mức: TH
\begin{ex}
% ID: L12C4B21-15-TLN
% Mức: TH
Trong bốn phát biểu sau về viết kí hiệu hạt nhân và xác định số proton, neutron, electron, có bao nhiêu phát biểu đúng? Chỉ ghi một số nguyên. (1) Cần đổi các đại lượng về hệ đơn vị thống nhất. (2) Với hạt nhân $^{A}_{Z}X$, số proton là $Z$, số neutron là $A-Z$; nguyên tử trung hòa có $Z$ electron. (3) Trong kí hiệu $^{A}_{Z}X$, số neutron bằng $Z$. (4) Không cần kiểm tra dữ kiện.
\shortans{2}
\loigiai{
Đối chiếu từng phát biểu với kiến thức cốt lõi: Với hạt nhân $^{A}_{Z}X$, số proton là $Z$, số neutron là $A-Z$; nguyên tử trung hòa có $Z$ electron. Có 2 phát biểu đúng.
}
\end{ex}

% ===== Câu 60 =====
% ID: L12C4B21-15-TN
% Mức: TH
\dangbt{Nhận biết cấu tạo hạt nhân, nuclon, số khối và điện tích hạt nhân}
\begin{ex}
% ID: L12C4B21-15-TN
% Mức: TH
Khi phân tích, phát biểu nào đúng về nhận biết cấu tạo hạt nhân, nuclon, số khối và điện tích hạt nhân?
\choice
{\True Hạt nhân gồm proton và neutron; số khối $A=Z+N$, điện tích hạt nhân bằng $+Ze$.}
{Số khối $A$ luôn bằng số neutron $N$.}
{Có thể bỏ qua điều kiện áp dụng và đơn vị trong mọi trường hợp.}
{Kết quả của bài toán không phụ thuộc dữ kiện đã cho.}
\loigiai{
Hạt nhân gồm proton và neutron; số khối $A=Z+N$, điện tích hạt nhân bằng $+Ze$. Vì vậy phương án $A$ đúng; các phương án còn lại trái với kiến thức hoặc điều kiện áp dụng.
}
\end{ex}

% ===== Câu 61 =====
% ID: L12C4B21-16-DS
% Mức: TH
\dangbt{Nhận biết đồng vị và tính nguyên tử khối trung bình}
\begin{ex}
% ID: L12C4B21-16-DS
% Mức: TH
Xét nhận biết đồng vị và tính nguyên tử khối trung bình (tình huống 1). Hãy xác định tính đúng, sai của bốn phát biểu sau.
\choiceTF
{\True Các đồng vị có cùng số proton nhưng khác số neutron; nguyên tử khối trung bình là trung bình có trọng số theo tỉ lệ đồng vị.}
{Các đồng vị của một nguyên tố có số proton khác nhau.}
{\True Khi giải cần đọc đúng dữ kiện, dùng hệ đơn vị thống nhất và kiểm tra điều kiện áp dụng.}
{Có thể bỏ qua dữ kiện, đơn vị và ý nghĩa vật lí của kết quả.}
\loigiai{
\begin{itemchoice}
\item (Đúng) Các đồng vị có cùng số proton nhưng khác số neutron; nguyên tử khối trung bình là trung bình có trọng số theo tỉ lệ đồng vị. (Vì): Các đồng vị có cùng số proton nhưng khác số neutron; nguyên tử khối trung bình là trung bình có trọng số theo tỉ lệ đồng vị. b) Sai: Các đồng vị của một nguyên tố có số proton khác nhau. c) Đúng: Khi giải cần đọc đúng dữ kiện, dùng hệ đơn vị thống nhất và kiểm tra điều kiện áp dụng. d) Sai: Có thể bỏ qua dữ kiện, đơn vị và ý nghĩa vật lí của kết quả. Kiến thức cốt lõi: Các đồng vị có cùng số proton nhưng khác số neutron; nguyên tử khối trung bình là trung bình có trọng số theo tỉ lệ đồng vị
\item (Sai) Các đồng vị của một nguyên tố có số proton khác nhau.
\item (Đúng) Khi giải cần đọc đúng dữ kiện, dùng hệ đơn vị thống nhất và kiểm tra điều kiện áp dụng.
\item (Sai) Có thể bỏ qua dữ kiện, đơn vị và ý nghĩa vật lí của kết quả.
\end{itemchoice}
}
\end{ex}

% ===== Câu 62 =====
% ID: L12C4B21-16-TL
% Mức: VD
\dangbt{Viết kí hiệu hạt nhân và xác định số proton, neutron, electron}
\begin{bt}
% ID: L12C4B21-16-TL
% Mức: VD
Đề xuất các bước giải một bài toán thuộc dạng viết kí hiệu hạt nhân và xác định số proton, neutron, electron.
\loigiai{
Cần nêu được: Với hạt nhân $^{A}_{Z}X$, số proton là $Z$, số neutron là $A-Z$; nguyên tử trung hòa có $Z$ electron. Khi vận dụng phải xác định đúng dữ kiện, điều kiện áp dụng, hệ đơn vị và kiểm tra tính hợp lí của kết quả. Nhận định “Trong kí hiệu $^{A}_{Z}X$, số neutron bằng $Z$.” sai vì trái với kiến thức trên.
}
\end{bt}

% ===== Câu 63 =====
% ID: L12C4B21-16-TLN
% Mức: VD
\begin{ex}
% ID: L12C4B21-16-TLN
% Mức: VD
Trong bốn phát biểu sau về viết kí hiệu hạt nhân và xác định số proton, neutron, electron, có bao nhiêu phát biểu đúng? Chỉ ghi một số nguyên. (1) Với hạt nhân $^{A}_{Z}X$, số proton là $Z$, số neutron là $A-Z$; nguyên tử trung hòa có $Z$ electron. (2) Cần đối chiếu kết quả với điều kiện bài toán. (3) Trong kí hiệu $^{A}_{Z}X$, số neutron bằng $Z$. (4) Mọi trường hợp đều cho cùng một kết quả.
\shortans{2}
\loigiai{
Đối chiếu từng phát biểu với kiến thức cốt lõi: Với hạt nhân $^{A}_{Z}X$, số proton là $Z$, số neutron là $A-Z$; nguyên tử trung hòa có $Z$ electron. Có 2 phát biểu đúng.
}
\end{ex}

% ===== Câu 64 =====
% ID: L12C4B21-16-TN
% Mức: TH
\dangbt{Nhận biết cấu tạo hạt nhân, nuclon, số khối và điện tích hạt nhân}
\begin{ex}
% ID: L12C4B21-16-TN
% Mức: TH
Kết luận nào của học sinh là đúng về nhận biết cấu tạo hạt nhân, nuclon, số khối và điện tích hạt nhân?
\choice
{Số khối $A$ luôn bằng số neutron $N$.}
{\True Hạt nhân gồm proton và neutron; số khối $A=Z+N$, điện tích hạt nhân bằng $+Ze$.}
{Có thể bỏ qua điều kiện áp dụng và đơn vị trong mọi trường hợp.}
{Kết quả của bài toán không phụ thuộc dữ kiện đã cho.}
\loigiai{
Hạt nhân gồm proton và neutron; số khối $A=Z+N$, điện tích hạt nhân bằng $+Ze$. Vì vậy phương án $B$ đúng; các phương án còn lại trái với kiến thức hoặc điều kiện áp dụng.
}
\end{ex}

% ===== Câu 65 =====
% ID: L12C4B21-17-DS
% Mức: TH
\dangbt{Nhận biết đồng vị và tính nguyên tử khối trung bình}
\begin{ex}
% ID: L12C4B21-17-DS
% Mức: TH
Xét nhận biết đồng vị và tính nguyên tử khối trung bình (tình huống 2). Hãy xác định tính đúng, sai của bốn phát biểu sau.
\choiceTF
{Các đồng vị của một nguyên tố có số proton khác nhau.}
{\True Các đồng vị có cùng số proton nhưng khác số neutron; nguyên tử khối trung bình là trung bình có trọng số theo tỉ lệ đồng vị.}
{Có thể bỏ qua dữ kiện, đơn vị và ý nghĩa vật lí của kết quả.}
{\True Khi giải cần đọc đúng dữ kiện, dùng hệ đơn vị thống nhất và kiểm tra điều kiện áp dụng.}
\loigiai{
\begin{itemchoice}
\item (Sai) Các đồng vị của một nguyên tố có số proton khác nhau. (Vì): Các đồng vị của một nguyên tố có số proton khác nhau. b) Đúng: Các đồng vị có cùng số proton nhưng khác số neutron; nguyên tử khối trung bình là trung bình có trọng số theo tỉ lệ đồng vị. c) Sai: Có thể bỏ qua dữ kiện, đơn vị và ý nghĩa vật lí của kết quả. d) Đúng: Khi giải cần đọc đúng dữ kiện, dùng hệ đơn vị thống nhất và kiểm tra điều kiện áp dụng. Kiến thức cốt lõi: Các đồng vị có cùng số proton nhưng khác số neutron; nguyên tử khối trung bình là trung bình có trọng số theo tỉ lệ đồng vị
\item (Đúng) Các đồng vị có cùng số proton nhưng khác số neutron; nguyên tử khối trung bình là trung bình có trọng số theo tỉ lệ đồng vị.
\item (Sai) Có thể bỏ qua dữ kiện, đơn vị và ý nghĩa vật lí của kết quả.
\item (Đúng) Khi giải cần đọc đúng dữ kiện, dùng hệ đơn vị thống nhất và kiểm tra điều kiện áp dụng.
\end{itemchoice}
}
\end{ex}

% ===== Câu 66 =====
% ID: L12C4B21-17-TL
% Mức: VD
\dangbt{Viết kí hiệu hạt nhân và xác định số proton, neutron, electron}
\begin{bt}
% ID: L12C4B21-17-TL
% Mức: VD
Nêu một tình huống thực tiễn liên quan đến viết kí hiệu hạt nhân và xác định số proton, neutron, electron và giải thích bằng kiến thức Vật lí.
\loigiai{
Cần nêu được: Với hạt nhân $^{A}_{Z}X$, số proton là $Z$, số neutron là $A-Z$; nguyên tử trung hòa có $Z$ electron. Khi vận dụng phải xác định đúng dữ kiện, điều kiện áp dụng, hệ đơn vị và kiểm tra tính hợp lí của kết quả. Nhận định “Trong kí hiệu $^{A}_{Z}X$, số neutron bằng $Z$.” sai vì trái với kiến thức trên.
}
\end{bt}

% ===== Câu 67 =====
% ID: L12C4B21-17-TLN
% Mức: VD
\begin{ex}
% ID: L12C4B21-17-TLN
% Mức: VD
Trong bốn phát biểu sau về viết kí hiệu hạt nhân và xác định số proton, neutron, electron, có bao nhiêu phát biểu đúng? Chỉ ghi một số nguyên. (1) Trong kí hiệu $^{A}_{Z}X$, số neutron bằng $Z$. (2) Cần đọc đúng dữ kiện và giả thiết. (3) Với hạt nhân $^{A}_{Z}X$, số proton là $Z$, số neutron là $A-Z$; nguyên tử trung hòa có $Z$ electron. (4) Có thể thay số trước khi chọn công thức.
\shortans{2}
\loigiai{
Đối chiếu từng phát biểu với kiến thức cốt lõi: Với hạt nhân $^{A}_{Z}X$, số proton là $Z$, số neutron là $A-Z$; nguyên tử trung hòa có $Z$ electron. Có 2 phát biểu đúng.
}
\end{ex}

% ===== Câu 68 =====
% ID: L12C4B21-17-TN
% Mức: TH
\dangbt{Nhận biết cấu tạo hạt nhân, nuclon, số khối và điện tích hạt nhân}
\begin{ex}
% ID: L12C4B21-17-TN
% Mức: TH
Chọn phát biểu không trái với kiến thức về nhận biết cấu tạo hạt nhân, nuclon, số khối và điện tích hạt nhân?
\choice
{Số khối $A$ luôn bằng số neutron $N$.}
{Có thể bỏ qua điều kiện áp dụng và đơn vị trong mọi trường hợp.}
{\True Hạt nhân gồm proton và neutron; số khối $A=Z+N$, điện tích hạt nhân bằng $+Ze$.}
{Kết quả của bài toán không phụ thuộc dữ kiện đã cho.}
\loigiai{
Hạt nhân gồm proton và neutron; số khối $A=Z+N$, điện tích hạt nhân bằng $+Ze$. Vì vậy phương án $C$ đúng; các phương án còn lại trái với kiến thức hoặc điều kiện áp dụng.
}
\end{ex}

% ===== Câu 69 =====
% ID: L12C4B21-18-DS
% Mức: VD
\dangbt{Nhận biết đồng vị và tính nguyên tử khối trung bình}
\begin{ex}
% ID: L12C4B21-18-DS
% Mức: VD
Xét nhận biết đồng vị và tính nguyên tử khối trung bình (tình huống 3). Hãy xác định tính đúng, sai của bốn phát biểu sau.
\choiceTF
{\True Các đồng vị có cùng số proton nhưng khác số neutron; nguyên tử khối trung bình là trung bình có trọng số theo tỉ lệ đồng vị.}
{Các đồng vị của một nguyên tố có số proton khác nhau.}
{\True Khi giải cần đọc đúng dữ kiện, dùng hệ đơn vị thống nhất và kiểm tra điều kiện áp dụng.}
{Có thể bỏ qua dữ kiện, đơn vị và ý nghĩa vật lí của kết quả.}
\loigiai{
\begin{itemchoice}
\item (Đúng) Các đồng vị có cùng số proton nhưng khác số neutron; nguyên tử khối trung bình là trung bình có trọng số theo tỉ lệ đồng vị. (Vì): Các đồng vị có cùng số proton nhưng khác số neutron; nguyên tử khối trung bình là trung bình có trọng số theo tỉ lệ đồng vị. b) Sai: Các đồng vị của một nguyên tố có số proton khác nhau. c) Đúng: Khi giải cần đọc đúng dữ kiện, dùng hệ đơn vị thống nhất và kiểm tra điều kiện áp dụng. d) Sai: Có thể bỏ qua dữ kiện, đơn vị và ý nghĩa vật lí của kết quả. Kiến thức cốt lõi: Các đồng vị có cùng số proton nhưng khác số neutron; nguyên tử khối trung bình là trung bình có trọng số theo tỉ lệ đồng vị
\item (Sai) Các đồng vị của một nguyên tố có số proton khác nhau.
\item (Đúng) Khi giải cần đọc đúng dữ kiện, dùng hệ đơn vị thống nhất và kiểm tra điều kiện áp dụng.
\item (Sai) Có thể bỏ qua dữ kiện, đơn vị và ý nghĩa vật lí của kết quả.
\end{itemchoice}
}
\end{ex}

% ===== Câu 70 =====
% ID: L12C4B21-18-TL
% Mức: VDC
\dangbt{Viết kí hiệu hạt nhân và xác định số proton, neutron, electron}
\begin{bt}
% ID: L12C4B21-18-TL
% Mức: VDC
So sánh nhận định đúng “Với hạt nhân $^{A}_{Z}X$, số proton là $Z$, số neutron là $A-Z$; nguyên tử trung hòa có $Z$ electron.” với nhận định sai “Trong kí hiệu $^{A}_{Z}X$, số neutron bằng $Z$.”; chỉ rõ căn cứ Vật lí.
\loigiai{
Cần nêu được: Với hạt nhân $^{A}_{Z}X$, số proton là $Z$, số neutron là $A-Z$; nguyên tử trung hòa có $Z$ electron. Khi vận dụng phải xác định đúng dữ kiện, điều kiện áp dụng, hệ đơn vị và kiểm tra tính hợp lí của kết quả. Nhận định “Trong kí hiệu $^{A}_{Z}X$, số neutron bằng $Z$.” sai vì trái với kiến thức trên.
}
\end{bt}

% ===== Câu 71 =====
% ID: L12C4B21-18-TLN
% Mức: VDC
\begin{ex}
% ID: L12C4B21-18-TLN
% Mức: VDC
Trong bốn phát biểu sau về viết kí hiệu hạt nhân và xác định số proton, neutron, electron, có bao nhiêu phát biểu đúng? Chỉ ghi một số nguyên. (1) Kết quả cần có ý nghĩa vật lí. (2) Trong kí hiệu $^{A}_{Z}X$, số neutron bằng $Z$. (3) Phải dùng đúng định luật cho tình huống. (4) Với hạt nhân $^{A}_{Z}X$, số proton là $Z$, số neutron là $A-Z$; nguyên tử trung hòa có $Z$ electron.
\shortans{3}
\loigiai{
Đối chiếu từng phát biểu với kiến thức cốt lõi: Với hạt nhân $^{A}_{Z}X$, số proton là $Z$, số neutron là $A-Z$; nguyên tử trung hòa có $Z$ electron. Có 3 phát biểu đúng.
}
\end{ex}

% ===== Câu 72 =====
% ID: L12C4B21-18-TN
% Mức: VD
\dangbt{Nhận biết cấu tạo hạt nhân, nuclon, số khối và điện tích hạt nhân}
\begin{ex}
% ID: L12C4B21-18-TN
% Mức: VD
Cơ sở Vật lí đúng để giải bài là gì về nhận biết cấu tạo hạt nhân, nuclon, số khối và điện tích hạt nhân?
\choice
{Số khối $A$ luôn bằng số neutron $N$.}
{Có thể bỏ qua điều kiện áp dụng và đơn vị trong mọi trường hợp.}
{Kết quả của bài toán không phụ thuộc dữ kiện đã cho.}
{\True Hạt nhân gồm proton và neutron; số khối $A=Z+N$, điện tích hạt nhân bằng $+Ze$.}
\loigiai{
Hạt nhân gồm proton và neutron; số khối $A=Z+N$, điện tích hạt nhân bằng $+Ze$. Vì vậy phương án $D$ đúng; các phương án còn lại trái với kiến thức hoặc điều kiện áp dụng.
}
\end{ex}

% ===== Câu 73 =====
% ID: L12C4B21-19-DS
% Mức: VD
\dangbt{Nhận biết đồng vị và tính nguyên tử khối trung bình}
\begin{ex}
% ID: L12C4B21-19-DS
% Mức: VD
Xét nhận biết đồng vị và tính nguyên tử khối trung bình (tình huống 4). Hãy xác định tính đúng, sai của bốn phát biểu sau.
\choiceTF
{Các đồng vị của một nguyên tố có số proton khác nhau.}
{\True Các đồng vị có cùng số proton nhưng khác số neutron; nguyên tử khối trung bình là trung bình có trọng số theo tỉ lệ đồng vị.}
{Có thể bỏ qua dữ kiện, đơn vị và ý nghĩa vật lí của kết quả.}
{\True Khi giải cần đọc đúng dữ kiện, dùng hệ đơn vị thống nhất và kiểm tra điều kiện áp dụng.}
\loigiai{
\begin{itemchoice}
\item (Sai) Các đồng vị của một nguyên tố có số proton khác nhau. (Vì): Các đồng vị của một nguyên tố có số proton khác nhau. b) Đúng: Các đồng vị có cùng số proton nhưng khác số neutron; nguyên tử khối trung bình là trung bình có trọng số theo tỉ lệ đồng vị. c) Sai: Có thể bỏ qua dữ kiện, đơn vị và ý nghĩa vật lí của kết quả. d) Đúng: Khi giải cần đọc đúng dữ kiện, dùng hệ đơn vị thống nhất và kiểm tra điều kiện áp dụng. Kiến thức cốt lõi: Các đồng vị có cùng số proton nhưng khác số neutron; nguyên tử khối trung bình là trung bình có trọng số theo tỉ lệ đồng vị
\item (Đúng) Các đồng vị có cùng số proton nhưng khác số neutron; nguyên tử khối trung bình là trung bình có trọng số theo tỉ lệ đồng vị.
\item (Sai) Có thể bỏ qua dữ kiện, đơn vị và ý nghĩa vật lí của kết quả.
\item (Đúng) Khi giải cần đọc đúng dữ kiện, dùng hệ đơn vị thống nhất và kiểm tra điều kiện áp dụng.
\end{itemchoice}
}
\end{ex}

% ===== Câu 74 =====
% ID: L12C4B21-19-TL
% Mức: TH
\begin{bt}
% ID: L12C4B21-19-TL
% Mức: TH
Trình bày kiến thức cốt lõi của nhận biết đồng vị và tính nguyên tử khối trung bình và nêu điều kiện áp dụng.
\loigiai{
Cần nêu được: Các đồng vị có cùng số proton nhưng khác số neutron; nguyên tử khối trung bình là trung bình có trọng số theo tỉ lệ đồng vị. Khi vận dụng phải xác định đúng dữ kiện, điều kiện áp dụng, hệ đơn vị và kiểm tra tính hợp lí của kết quả. Nhận định “Các đồng vị của một nguyên tố có số proton khác nhau.” sai vì trái với kiến thức trên.
}
\end{bt}

% ===== Câu 75 =====
% ID: L12C4B21-19-TLN
% Mức: TH
\begin{ex}
% ID: L12C4B21-19-TLN
% Mức: TH
Trong bốn phát biểu sau về nhận biết đồng vị và tính nguyên tử khối trung bình, có bao nhiêu phát biểu đúng? Chỉ ghi một số nguyên. (1) Các đồng vị có cùng số proton nhưng khác số neutron; nguyên tử khối trung bình là trung bình có trọng số theo tỉ lệ đồng vị. (2) Các đồng vị của một nguyên tố có số proton khác nhau. (3) Cần kiểm tra điều kiện áp dụng trước khi tính toán. (4) Có thể bỏ qua đơn vị của kết quả.
\shortans{2}
\loigiai{
Đối chiếu từng phát biểu với kiến thức cốt lõi: Các đồng vị có cùng số proton nhưng khác số neutron; nguyên tử khối trung bình là trung bình có trọng số theo tỉ lệ đồng vị. Có 2 phát biểu đúng.
}
\end{ex}

% ===== Câu 76 =====
% ID: L12C4B21-19-TN
% Mức: VD
\dangbt{Nhận biết cấu tạo hạt nhân, nuclon, số khối và điện tích hạt nhân}
\begin{ex}
% ID: L12C4B21-19-TN
% Mức: VD
Trong một bài thực hành, nhận xét nào đúng về nhận biết cấu tạo hạt nhân, nuclon, số khối và điện tích hạt nhân?
\choice
{\True Hạt nhân gồm proton và neutron; số khối $A=Z+N$, điện tích hạt nhân bằng $+Ze$.}
{Số khối $A$ luôn bằng số neutron $N$.}
{Có thể bỏ qua điều kiện áp dụng và đơn vị trong mọi trường hợp.}
{Kết quả của bài toán không phụ thuộc dữ kiện đã cho.}
\loigiai{
Hạt nhân gồm proton và neutron; số khối $A=Z+N$, điện tích hạt nhân bằng $+Ze$. Vì vậy phương án $A$ đúng; các phương án còn lại trái với kiến thức hoặc điều kiện áp dụng.
}
\end{ex}

% ===== Câu 77 =====
% ID: L12C4B21-20-DS
% Mức: VD
\dangbt{Nhận biết đồng vị và tính nguyên tử khối trung bình}
\begin{ex}
% ID: L12C4B21-20-DS
% Mức: VD
Xét nhận biết đồng vị và tính nguyên tử khối trung bình (tình huống 5). Hãy xác định tính đúng, sai của bốn phát biểu sau.
\choiceTF
{\True Các đồng vị có cùng số proton nhưng khác số neutron; nguyên tử khối trung bình là trung bình có trọng số theo tỉ lệ đồng vị.}
{Có thể bỏ qua dữ kiện, đơn vị và ý nghĩa vật lí của kết quả.}
{Các đồng vị của một nguyên tố có số proton khác nhau.}
{Có thể bỏ qua dữ kiện, đơn vị và ý nghĩa vật lí của kết quả.}
\loigiai{
\begin{itemchoice}
\item (Đúng) Các đồng vị có cùng số proton nhưng khác số neutron; nguyên tử khối trung bình là trung bình có trọng số theo tỉ lệ đồng vị. (Vì): Các đồng vị có cùng số proton nhưng khác số neutron; nguyên tử khối trung bình là trung bình có trọng số theo tỉ lệ đồng vị. b) Sai: Có thể bỏ qua dữ kiện, đơn vị và ý nghĩa vật lí của kết quả. c) Sai: Các đồng vị của một nguyên tố có số proton khác nhau. d) Sai: Có thể bỏ qua dữ kiện, đơn vị và ý nghĩa vật lí của kết quả. Kiến thức cốt lõi: Các đồng vị có cùng số proton nhưng khác số neutron; nguyên tử khối trung bình là trung bình có trọng số theo tỉ lệ đồng vị
\item (Sai) Có thể bỏ qua dữ kiện, đơn vị và ý nghĩa vật lí của kết quả.
\item (Sai) Các đồng vị của một nguyên tố có số proton khác nhau.
\item (Sai) Có thể bỏ qua dữ kiện, đơn vị và ý nghĩa vật lí của kết quả.
\end{itemchoice}
}
\end{ex}

% ===== Câu 78 =====
% ID: L12C4B21-20-TL
% Mức: TH
\begin{bt}
% ID: L12C4B21-20-TL
% Mức: TH
Giải thích vì sao nhận định sau đúng: “Các đồng vị có cùng số proton nhưng khác số neutron; nguyên tử khối trung bình là trung bình có trọng số theo tỉ lệ đồng vị.”
\loigiai{
Cần nêu được: Các đồng vị có cùng số proton nhưng khác số neutron; nguyên tử khối trung bình là trung bình có trọng số theo tỉ lệ đồng vị. Khi vận dụng phải xác định đúng dữ kiện, điều kiện áp dụng, hệ đơn vị và kiểm tra tính hợp lí của kết quả. Nhận định “Các đồng vị của một nguyên tố có số proton khác nhau.” sai vì trái với kiến thức trên.
}
\end{bt}

% ===== Câu 79 =====
% ID: L12C4B21-20-TLN
% Mức: TH
\begin{ex}
% ID: L12C4B21-20-TLN
% Mức: TH
Trong bốn phát biểu sau về nhận biết đồng vị và tính nguyên tử khối trung bình, có bao nhiêu phát biểu đúng? Chỉ ghi một số nguyên. (1) Các đồng vị của một nguyên tố có số proton khác nhau. (2) Các đồng vị có cùng số proton nhưng khác số neutron; nguyên tử khối trung bình là trung bình có trọng số theo tỉ lệ đồng vị. (3) Kết quả phải phù hợp ý nghĩa vật lí. (4) Mọi đại lượng đều có cùng bản chất.
\shortans{2}
\loigiai{
Đối chiếu từng phát biểu với kiến thức cốt lõi: Các đồng vị có cùng số proton nhưng khác số neutron; nguyên tử khối trung bình là trung bình có trọng số theo tỉ lệ đồng vị. Có 2 phát biểu đúng.
}
\end{ex}

% ===== Câu 80 =====
% ID: L12C4B21-20-TN
% Mức: VD
\dangbt{Nhận biết cấu tạo hạt nhân, nuclon, số khối và điện tích hạt nhân}
\begin{ex}
% ID: L12C4B21-20-TN
% Mức: VD
Kết luận đúng khi vận dụng là về nhận biết cấu tạo hạt nhân, nuclon, số khối và điện tích hạt nhân?
\choice
{Số khối $A$ luôn bằng số neutron $N$.}
{\True Hạt nhân gồm proton và neutron; số khối $A=Z+N$, điện tích hạt nhân bằng $+Ze$.}
{Có thể bỏ qua điều kiện áp dụng và đơn vị trong mọi trường hợp.}
{Kết quả của bài toán không phụ thuộc dữ kiện đã cho.}
\loigiai{
Hạt nhân gồm proton và neutron; số khối $A=Z+N$, điện tích hạt nhân bằng $+Ze$. Vì vậy phương án $B$ đúng; các phương án còn lại trái với kiến thức hoặc điều kiện áp dụng.
}
\end{ex}

% ===== Câu 81 =====
% ID: L12C4B21-21-DS
% Mức: TH
\dangbt{Tính kích thước, bán kính và khối lượng riêng của hạt nhân}
\begin{ex}
% ID: L12C4B21-21-DS
% Mức: TH
Xét tính kích thước, bán kính và khối lượng riêng của hạt nhân (tình huống 1). Hãy xác định tính đúng, sai của bốn phát biểu sau.
\choiceTF
{\True Bán kính hạt nhân gần đúng $R=R_0A^{1/3}$ nên thể tích tỉ lệ với $A$ và khối lượng riêng hạt nhân xấp xỉ không đổi.}
{Bán kính hạt nhân tỉ lệ thuận với số khối $A$.}
{\True Khi giải cần đọc đúng dữ kiện, dùng hệ đơn vị thống nhất và kiểm tra điều kiện áp dụng.}
{Có thể bỏ qua dữ kiện, đơn vị và ý nghĩa vật lí của kết quả.}
\loigiai{
\begin{itemchoice}
\item (Đúng) Bán kính hạt nhân gần đúng $R=R_0A^{1/3}$ nên thể tích tỉ lệ với $A$ và khối lượng riêng hạt nhân xấp xỉ không đổi. (Vì): Bán kính hạt nhân gần đúng $R=R_0A^{1/3}$ nên thể tích tỉ lệ với $A$ và khối lượng riêng hạt nhân xấp xỉ không đổi. b) Sai: Bán kính hạt nhân tỉ lệ thuận với số khối $A$. c) Đúng: Khi giải cần đọc đúng dữ kiện, dùng hệ đơn vị thống nhất và kiểm tra điều kiện áp dụng. d) Sai: Có thể bỏ qua dữ kiện, đơn vị và ý nghĩa vật lí của kết quả. Kiến thức cốt lõi: Bán kính hạt nhân gần đúng $R=R_0A^{1/3}$ nên thể tích tỉ lệ với $A$ và khối lượng riêng hạt nhân xấp xỉ không đổi
\item (Sai) Bán kính hạt nhân tỉ lệ thuận với số khối $A$.
\item (Đúng) Khi giải cần đọc đúng dữ kiện, dùng hệ đơn vị thống nhất và kiểm tra điều kiện áp dụng.
\item (Sai) Có thể bỏ qua dữ kiện, đơn vị và ý nghĩa vật lí của kết quả.
\end{itemchoice}
}
\end{ex}

% ===== Câu 82 =====
% ID: L12C4B21-21-TL
% Mức: VD
\dangbt{Nhận biết đồng vị và tính nguyên tử khối trung bình}
\begin{bt}
% ID: L12C4B21-21-TL
% Mức: VD
Phân tích sai lầm trong nhận định: “Các đồng vị của một nguyên tố có số proton khác nhau.”
\loigiai{
Cần nêu được: Các đồng vị có cùng số proton nhưng khác số neutron; nguyên tử khối trung bình là trung bình có trọng số theo tỉ lệ đồng vị. Khi vận dụng phải xác định đúng dữ kiện, điều kiện áp dụng, hệ đơn vị và kiểm tra tính hợp lí của kết quả. Nhận định “Các đồng vị của một nguyên tố có số proton khác nhau.” sai vì trái với kiến thức trên.
}
\end{bt}

% ===== Câu 83 =====
% ID: L12C4B21-21-TLN
% Mức: TH
\begin{ex}
% ID: L12C4B21-21-TLN
% Mức: TH
Trong bốn phát biểu sau về nhận biết đồng vị và tính nguyên tử khối trung bình, có bao nhiêu phát biểu đúng? Chỉ ghi một số nguyên. (1) Cần đổi các đại lượng về hệ đơn vị thống nhất. (2) Các đồng vị có cùng số proton nhưng khác số neutron; nguyên tử khối trung bình là trung bình có trọng số theo tỉ lệ đồng vị. (3) Các đồng vị của một nguyên tố có số proton khác nhau. (4) Không cần kiểm tra dữ kiện.
\shortans{2}
\loigiai{
Đối chiếu từng phát biểu với kiến thức cốt lõi: Các đồng vị có cùng số proton nhưng khác số neutron; nguyên tử khối trung bình là trung bình có trọng số theo tỉ lệ đồng vị. Có 2 phát biểu đúng.
}
\end{ex}

% ===== Câu 84 =====
% ID: L12C4B21-21-TN
% Mức: NB
\dangbt{Viết kí hiệu hạt nhân và xác định số proton, neutron, electron}
\begin{ex}
% ID: L12C4B21-21-TN
% Mức: NB
Phát biểu nào sau đây đúng về viết kí hiệu hạt nhân và xác định số proton, neutron, electron?
\choice
{\True Với hạt nhân $^{A}_{Z}X$, số proton là $Z$, số neutron là $A-Z$; nguyên tử trung hòa có $Z$ electron.}
{Trong kí hiệu $^{A}_{Z}X$, số neutron bằng $Z$.}
{Có thể bỏ qua điều kiện áp dụng và đơn vị trong mọi trường hợp.}
{Kết quả của bài toán không phụ thuộc dữ kiện đã cho.}
\loigiai{
Với hạt nhân $^{A}_{Z}X$, số proton là $Z$, số neutron là $A-Z$; nguyên tử trung hòa có $Z$ electron. Vì vậy phương án $A$ đúng; các phương án còn lại trái với kiến thức hoặc điều kiện áp dụng.
}
\end{ex}

% ===== Câu 85 =====
% ID: L12C4B21-22-DS
% Mức: TH
\dangbt{Tính kích thước, bán kính và khối lượng riêng của hạt nhân}
\begin{ex}
% ID: L12C4B21-22-DS
% Mức: TH
Xét tính kích thước, bán kính và khối lượng riêng của hạt nhân (tình huống 2). Hãy xác định tính đúng, sai của bốn phát biểu sau.
\choiceTF
{Bán kính hạt nhân tỉ lệ thuận với số khối $A$.}
{\True Bán kính hạt nhân gần đúng $R=R_0A^{1/3}$ nên thể tích tỉ lệ với $A$ và khối lượng riêng hạt nhân xấp xỉ không đổi.}
{Có thể bỏ qua dữ kiện, đơn vị và ý nghĩa vật lí của kết quả.}
{\True Khi giải cần đọc đúng dữ kiện, dùng hệ đơn vị thống nhất và kiểm tra điều kiện áp dụng.}
\loigiai{
\begin{itemchoice}
\item (Sai) Bán kính hạt nhân tỉ lệ thuận với số khối $A$. (Vì): Bán kính hạt nhân tỉ lệ thuận với số khối $A$. b) Đúng: Bán kính hạt nhân gần đúng $R=R_0A^{1/3}$ nên thể tích tỉ lệ với $A$ và khối lượng riêng hạt nhân xấp xỉ không đổi. c) Sai: Có thể bỏ qua dữ kiện, đơn vị và ý nghĩa vật lí của kết quả. d) Đúng: Khi giải cần đọc đúng dữ kiện, dùng hệ đơn vị thống nhất và kiểm tra điều kiện áp dụng. Kiến thức cốt lõi: Bán kính hạt nhân gần đúng $R=R_0A^{1/3}$ nên thể tích tỉ lệ với $A$ và khối lượng riêng hạt nhân xấp xỉ không đổi
\item (Đúng) Bán kính hạt nhân gần đúng $R=R_0A^{1/3}$ nên thể tích tỉ lệ với $A$ và khối lượng riêng hạt nhân xấp xỉ không đổi.
\item (Sai) Có thể bỏ qua dữ kiện, đơn vị và ý nghĩa vật lí của kết quả.
\item (Đúng) Khi giải cần đọc đúng dữ kiện, dùng hệ đơn vị thống nhất và kiểm tra điều kiện áp dụng.
\end{itemchoice}
}
\end{ex}

% ===== Câu 86 =====
% ID: L12C4B21-22-TL
% Mức: VD
\dangbt{Nhận biết đồng vị và tính nguyên tử khối trung bình}
\begin{bt}
% ID: L12C4B21-22-TL
% Mức: VD
Đề xuất các bước giải một bài toán thuộc dạng nhận biết đồng vị và tính nguyên tử khối trung bình.
\loigiai{
Cần nêu được: Các đồng vị có cùng số proton nhưng khác số neutron; nguyên tử khối trung bình là trung bình có trọng số theo tỉ lệ đồng vị. Khi vận dụng phải xác định đúng dữ kiện, điều kiện áp dụng, hệ đơn vị và kiểm tra tính hợp lí của kết quả. Nhận định “Các đồng vị của một nguyên tố có số proton khác nhau.” sai vì trái với kiến thức trên.
}
\end{bt}

% ===== Câu 87 =====
% ID: L12C4B21-22-TLN
% Mức: VD
\begin{ex}
% ID: L12C4B21-22-TLN
% Mức: VD
Trong bốn phát biểu sau về nhận biết đồng vị và tính nguyên tử khối trung bình, có bao nhiêu phát biểu đúng? Chỉ ghi một số nguyên. (1) Các đồng vị có cùng số proton nhưng khác số neutron; nguyên tử khối trung bình là trung bình có trọng số theo tỉ lệ đồng vị. (2) Cần đối chiếu kết quả với điều kiện bài toán. (3) Các đồng vị của một nguyên tố có số proton khác nhau. (4) Mọi trường hợp đều cho cùng một kết quả.
\shortans{2}
\loigiai{
Đối chiếu từng phát biểu với kiến thức cốt lõi: Các đồng vị có cùng số proton nhưng khác số neutron; nguyên tử khối trung bình là trung bình có trọng số theo tỉ lệ đồng vị. Có 2 phát biểu đúng.
}
\end{ex}

% ===== Câu 88 =====
% ID: L12C4B21-22-TN
% Mức: NB
\dangbt{Viết kí hiệu hạt nhân và xác định số proton, neutron, electron}
\begin{ex}
% ID: L12C4B21-22-TN
% Mức: NB
Nhận định nào phù hợp với kiến thức Vật lí về viết kí hiệu hạt nhân và xác định số proton, neutron, electron?
\choice
{Trong kí hiệu $^{A}_{Z}X$, số neutron bằng $Z$.}
{\True Với hạt nhân $^{A}_{Z}X$, số proton là $Z$, số neutron là $A-Z$; nguyên tử trung hòa có $Z$ electron.}
{Có thể bỏ qua điều kiện áp dụng và đơn vị trong mọi trường hợp.}
{Kết quả của bài toán không phụ thuộc dữ kiện đã cho.}
\loigiai{
Với hạt nhân $^{A}_{Z}X$, số proton là $Z$, số neutron là $A-Z$; nguyên tử trung hòa có $Z$ electron. Vì vậy phương án $B$ đúng; các phương án còn lại trái với kiến thức hoặc điều kiện áp dụng.
}
\end{ex}

% ===== Câu 89 =====
% ID: L12C4B21-23-DS
% Mức: VD
\dangbt{Tính kích thước, bán kính và khối lượng riêng của hạt nhân}
\begin{ex}
% ID: L12C4B21-23-DS
% Mức: VD
Xét tính kích thước, bán kính và khối lượng riêng của hạt nhân (tình huống 3). Hãy xác định tính đúng, sai của bốn phát biểu sau.
\choiceTF
{\True Bán kính hạt nhân gần đúng $R=R_0A^{1/3}$ nên thể tích tỉ lệ với $A$ và khối lượng riêng hạt nhân xấp xỉ không đổi.}
{Bán kính hạt nhân tỉ lệ thuận với số khối $A$.}
{\True Khi giải cần đọc đúng dữ kiện, dùng hệ đơn vị thống nhất và kiểm tra điều kiện áp dụng.}
{Có thể bỏ qua dữ kiện, đơn vị và ý nghĩa vật lí của kết quả.}
\loigiai{
\begin{itemchoice}
\item (Đúng) Bán kính hạt nhân gần đúng $R=R_0A^{1/3}$ nên thể tích tỉ lệ với $A$ và khối lượng riêng hạt nhân xấp xỉ không đổi. (Vì): Bán kính hạt nhân gần đúng $R=R_0A^{1/3}$ nên thể tích tỉ lệ với $A$ và khối lượng riêng hạt nhân xấp xỉ không đổi. b) Sai: Bán kính hạt nhân tỉ lệ thuận với số khối $A$. c) Đúng: Khi giải cần đọc đúng dữ kiện, dùng hệ đơn vị thống nhất và kiểm tra điều kiện áp dụng. d) Sai: Có thể bỏ qua dữ kiện, đơn vị và ý nghĩa vật lí của kết quả. Kiến thức cốt lõi: Bán kính hạt nhân gần đúng $R=R_0A^{1/3}$ nên thể tích tỉ lệ với $A$ và khối lượng riêng hạt nhân xấp xỉ không đổi
\item (Sai) Bán kính hạt nhân tỉ lệ thuận với số khối $A$.
\item (Đúng) Khi giải cần đọc đúng dữ kiện, dùng hệ đơn vị thống nhất và kiểm tra điều kiện áp dụng.
\item (Sai) Có thể bỏ qua dữ kiện, đơn vị và ý nghĩa vật lí của kết quả.
\end{itemchoice}
}
\end{ex}

% ===== Câu 90 =====
% ID: L12C4B21-23-TL
% Mức: VD
\dangbt{Nhận biết đồng vị và tính nguyên tử khối trung bình}
\begin{bt}
% ID: L12C4B21-23-TL
% Mức: VD
Nêu một tình huống thực tiễn liên quan đến nhận biết đồng vị và tính nguyên tử khối trung bình và giải thích bằng kiến thức Vật lí.
\loigiai{
Cần nêu được: Các đồng vị có cùng số proton nhưng khác số neutron; nguyên tử khối trung bình là trung bình có trọng số theo tỉ lệ đồng vị. Khi vận dụng phải xác định đúng dữ kiện, điều kiện áp dụng, hệ đơn vị và kiểm tra tính hợp lí của kết quả. Nhận định “Các đồng vị của một nguyên tố có số proton khác nhau.” sai vì trái với kiến thức trên.
}
\end{bt}

% ===== Câu 91 =====
% ID: L12C4B21-23-TLN
% Mức: VD
\begin{ex}
% ID: L12C4B21-23-TLN
% Mức: VD
Trong bốn phát biểu sau về nhận biết đồng vị và tính nguyên tử khối trung bình, có bao nhiêu phát biểu đúng? Chỉ ghi một số nguyên. (1) Các đồng vị của một nguyên tố có số proton khác nhau. (2) Cần đọc đúng dữ kiện và giả thiết. (3) Các đồng vị có cùng số proton nhưng khác số neutron; nguyên tử khối trung bình là trung bình có trọng số theo tỉ lệ đồng vị. (4) Có thể thay số trước khi chọn công thức.
\shortans{2}
\loigiai{
Đối chiếu từng phát biểu với kiến thức cốt lõi: Các đồng vị có cùng số proton nhưng khác số neutron; nguyên tử khối trung bình là trung bình có trọng số theo tỉ lệ đồng vị. Có 2 phát biểu đúng.
}
\end{ex}

% ===== Câu 92 =====
% ID: L12C4B21-23-TN
% Mức: NB
\dangbt{Viết kí hiệu hạt nhân và xác định số proton, neutron, electron}
\begin{ex}
% ID: L12C4B21-23-TN
% Mức: NB
Chọn kết luận chính xác về viết kí hiệu hạt nhân và xác định số proton, neutron, electron?
\choice
{Trong kí hiệu $^{A}_{Z}X$, số neutron bằng $Z$.}
{Có thể bỏ qua điều kiện áp dụng và đơn vị trong mọi trường hợp.}
{\True Với hạt nhân $^{A}_{Z}X$, số proton là $Z$, số neutron là $A-Z$; nguyên tử trung hòa có $Z$ electron.}
{Kết quả của bài toán không phụ thuộc dữ kiện đã cho.}
\loigiai{
Với hạt nhân $^{A}_{Z}X$, số proton là $Z$, số neutron là $A-Z$; nguyên tử trung hòa có $Z$ electron. Vì vậy phương án $C$ đúng; các phương án còn lại trái với kiến thức hoặc điều kiện áp dụng.
}
\end{ex}

% ===== Câu 93 =====
% ID: L12C4B21-24-DS
% Mức: VD
\dangbt{Tính kích thước, bán kính và khối lượng riêng của hạt nhân}
\begin{ex}
% ID: L12C4B21-24-DS
% Mức: VD
Xét tính kích thước, bán kính và khối lượng riêng của hạt nhân (tình huống 4). Hãy xác định tính đúng, sai của bốn phát biểu sau.
\choiceTF
{Bán kính hạt nhân tỉ lệ thuận với số khối $A$.}
{\True Bán kính hạt nhân gần đúng $R=R_0A^{1/3}$ nên thể tích tỉ lệ với $A$ và khối lượng riêng hạt nhân xấp xỉ không đổi.}
{Có thể bỏ qua dữ kiện, đơn vị và ý nghĩa vật lí của kết quả.}
{\True Khi giải cần đọc đúng dữ kiện, dùng hệ đơn vị thống nhất và kiểm tra điều kiện áp dụng.}
\loigiai{
\begin{itemchoice}
\item (Sai) Bán kính hạt nhân tỉ lệ thuận với số khối $A$. (Vì): Bán kính hạt nhân tỉ lệ thuận với số khối $A$. b) Đúng: Bán kính hạt nhân gần đúng $R=R_0A^{1/3}$ nên thể tích tỉ lệ với $A$ và khối lượng riêng hạt nhân xấp xỉ không đổi. c) Sai: Có thể bỏ qua dữ kiện, đơn vị và ý nghĩa vật lí của kết quả. d) Đúng: Khi giải cần đọc đúng dữ kiện, dùng hệ đơn vị thống nhất và kiểm tra điều kiện áp dụng. Kiến thức cốt lõi: Bán kính hạt nhân gần đúng $R=R_0A^{1/3}$ nên thể tích tỉ lệ với $A$ và khối lượng riêng hạt nhân xấp xỉ không đổi
\item (Đúng) Bán kính hạt nhân gần đúng $R=R_0A^{1/3}$ nên thể tích tỉ lệ với $A$ và khối lượng riêng hạt nhân xấp xỉ không đổi.
\item (Sai) Có thể bỏ qua dữ kiện, đơn vị và ý nghĩa vật lí của kết quả.
\item (Đúng) Khi giải cần đọc đúng dữ kiện, dùng hệ đơn vị thống nhất và kiểm tra điều kiện áp dụng.
\end{itemchoice}
}
\end{ex}

% ===== Câu 94 =====
% ID: L12C4B21-24-TL
% Mức: VDC
\dangbt{Nhận biết đồng vị và tính nguyên tử khối trung bình}
\begin{bt}
% ID: L12C4B21-24-TL
% Mức: VDC
So sánh nhận định đúng “Các đồng vị có cùng số proton nhưng khác số neutron; nguyên tử khối trung bình là trung bình có trọng số theo tỉ lệ đồng vị.” với nhận định sai “Các đồng vị của một nguyên tố có số proton khác nhau.”; chỉ rõ căn cứ Vật lí.
\loigiai{
Cần nêu được: Các đồng vị có cùng số proton nhưng khác số neutron; nguyên tử khối trung bình là trung bình có trọng số theo tỉ lệ đồng vị. Khi vận dụng phải xác định đúng dữ kiện, điều kiện áp dụng, hệ đơn vị và kiểm tra tính hợp lí của kết quả. Nhận định “Các đồng vị của một nguyên tố có số proton khác nhau.” sai vì trái với kiến thức trên.
}
\end{bt}

% ===== Câu 95 =====
% ID: L12C4B21-24-TLN
% Mức: VDC
\begin{ex}
% ID: L12C4B21-24-TLN
% Mức: VDC
Trong bốn phát biểu sau về nhận biết đồng vị và tính nguyên tử khối trung bình, có bao nhiêu phát biểu đúng? Chỉ ghi một số nguyên. (1) Kết quả cần có ý nghĩa vật lí. (2) Các đồng vị của một nguyên tố có số proton khác nhau. (3) Phải dùng đúng định luật cho tình huống. (4) Các đồng vị có cùng số proton nhưng khác số neutron; nguyên tử khối trung bình là trung bình có trọng số theo tỉ lệ đồng vị.
\shortans{3}
\loigiai{
Đối chiếu từng phát biểu với kiến thức cốt lõi: Các đồng vị có cùng số proton nhưng khác số neutron; nguyên tử khối trung bình là trung bình có trọng số theo tỉ lệ đồng vị. Có 3 phát biểu đúng.
}
\end{ex}

% ===== Câu 96 =====
% ID: L12C4B21-24-TN
% Mức: TH
\dangbt{Viết kí hiệu hạt nhân và xác định số proton, neutron, electron}
\begin{ex}
% ID: L12C4B21-24-TN
% Mức: TH
Nội dung nào mô tả đúng nhất về viết kí hiệu hạt nhân và xác định số proton, neutron, electron?
\choice
{Trong kí hiệu $^{A}_{Z}X$, số neutron bằng $Z$.}
{Có thể bỏ qua điều kiện áp dụng và đơn vị trong mọi trường hợp.}
{Kết quả của bài toán không phụ thuộc dữ kiện đã cho.}
{\True Với hạt nhân $^{A}_{Z}X$, số proton là $Z$, số neutron là $A-Z$; nguyên tử trung hòa có $Z$ electron.}
\loigiai{
Với hạt nhân $^{A}_{Z}X$, số proton là $Z$, số neutron là $A-Z$; nguyên tử trung hòa có $Z$ electron. Vì vậy phương án $D$ đúng; các phương án còn lại trái với kiến thức hoặc điều kiện áp dụng.
}
\end{ex}

% ===== Câu 97 =====
% ID: L12C4B21-25-DS
% Mức: VD
\dangbt{Tính kích thước, bán kính và khối lượng riêng của hạt nhân}
\begin{ex}
% ID: L12C4B21-25-DS
% Mức: VD
Xét tính kích thước, bán kính và khối lượng riêng của hạt nhân (tình huống 5). Hãy xác định tính đúng, sai của bốn phát biểu sau.
\choiceTF
{\True Bán kính hạt nhân gần đúng $R=R_0A^{1/3}$ nên thể tích tỉ lệ với $A$ và khối lượng riêng hạt nhân xấp xỉ không đổi.}
{Có thể bỏ qua dữ kiện, đơn vị và ý nghĩa vật lí của kết quả.}
{Bán kính hạt nhân tỉ lệ thuận với số khối $A$.}
{Có thể bỏ qua dữ kiện, đơn vị và ý nghĩa vật lí của kết quả.}
\loigiai{
\begin{itemchoice}
\item (Đúng) Bán kính hạt nhân gần đúng $R=R_0A^{1/3}$ nên thể tích tỉ lệ với $A$ và khối lượng riêng hạt nhân xấp xỉ không đổi. (Vì): Bán kính hạt nhân gần đúng $R=R_0A^{1/3}$ nên thể tích tỉ lệ với $A$ và khối lượng riêng hạt nhân xấp xỉ không đổi. b) Sai: Có thể bỏ qua dữ kiện, đơn vị và ý nghĩa vật lí của kết quả. c) Sai: Bán kính hạt nhân tỉ lệ thuận với số khối $A$. d) Sai: Có thể bỏ qua dữ kiện, đơn vị và ý nghĩa vật lí của kết quả. Kiến thức cốt lõi: Bán kính hạt nhân gần đúng $R=R_0A^{1/3}$ nên thể tích tỉ lệ với $A$ và khối lượng riêng hạt nhân xấp xỉ không đổi
\item (Sai) Có thể bỏ qua dữ kiện, đơn vị và ý nghĩa vật lí của kết quả.
\item (Sai) Bán kính hạt nhân tỉ lệ thuận với số khối $A$.
\item (Sai) Có thể bỏ qua dữ kiện, đơn vị và ý nghĩa vật lí của kết quả.
\end{itemchoice}
}
\end{ex}

% ===== Câu 98 =====
% ID: L12C4B21-25-TL
% Mức: TH
\begin{bt}
% ID: L12C4B21-25-TL
% Mức: TH
Trình bày kiến thức cốt lõi của tính kích thước, bán kính và khối lượng riêng của hạt nhân và nêu điều kiện áp dụng.
\loigiai{
Cần nêu được: Bán kính hạt nhân gần đúng $R=R_0A^{1/3}$ nên thể tích tỉ lệ với $A$ và khối lượng riêng hạt nhân xấp xỉ không đổi. Khi vận dụng phải xác định đúng dữ kiện, điều kiện áp dụng, hệ đơn vị và kiểm tra tính hợp lí của kết quả. Nhận định “Bán kính hạt nhân tỉ lệ thuận với số khối $A$.” sai vì trái với kiến thức trên.
}
\end{bt}

% ===== Câu 99 =====
% ID: L12C4B21-25-TLN
% Mức: TH
\begin{ex}
% ID: L12C4B21-25-TLN
% Mức: TH
Trong bốn phát biểu sau về tính kích thước, bán kính và khối lượng riêng của hạt nhân, có bao nhiêu phát biểu đúng? Chỉ ghi một số nguyên. (1) Bán kính hạt nhân gần đúng $R=R_0A^{1/3}$ nên thể tích tỉ lệ với $A$ và khối lượng riêng hạt nhân xấp xỉ không đổi. (2) Bán kính hạt nhân tỉ lệ thuận với số khối $A$. (3) Cần kiểm tra điều kiện áp dụng trước khi tính toán. (4) Có thể bỏ qua đơn vị của kết quả.
\shortans{2}
\loigiai{
Đối chiếu từng phát biểu với kiến thức cốt lõi: Bán kính hạt nhân gần đúng $R=R_0A^{1/3}$ nên thể tích tỉ lệ với $A$ và khối lượng riêng hạt nhân xấp xỉ không đổi. Có 2 phát biểu đúng.
}
\end{ex}

% ===== Câu 100 =====
% ID: L12C4B21-25-TN
% Mức: TH
\dangbt{Viết kí hiệu hạt nhân và xác định số proton, neutron, electron}
\begin{ex}
% ID: L12C4B21-25-TN
% Mức: TH
Khi phân tích, phát biểu nào đúng về viết kí hiệu hạt nhân và xác định số proton, neutron, electron?
\choice
{\True Với hạt nhân $^{A}_{Z}X$, số proton là $Z$, số neutron là $A-Z$; nguyên tử trung hòa có $Z$ electron.}
{Trong kí hiệu $^{A}_{Z}X$, số neutron bằng $Z$.}
{Có thể bỏ qua điều kiện áp dụng và đơn vị trong mọi trường hợp.}
{Kết quả của bài toán không phụ thuộc dữ kiện đã cho.}
\loigiai{
Với hạt nhân $^{A}_{Z}X$, số proton là $Z$, số neutron là $A-Z$; nguyên tử trung hòa có $Z$ electron. Vì vậy phương án $A$ đúng; các phương án còn lại trái với kiến thức hoặc điều kiện áp dụng.
}
\end{ex}

% ===== Câu 101 =====
% ID: L12C4B21-26-DS
% Mức: TH
\dangbt{Vận dụng hệ thức khối lượng – năng lượng và độ hụt khối}
\begin{ex}
% ID: L12C4B21-26-DS
% Mức: TH
Xét vận dụng hệ thức khối lượng – năng lượng và độ hụt khối (tình huống 1). Hãy xác định tính đúng, sai của bốn phát biểu sau.
\choiceTF
{\True Khối lượng và năng lượng liên hệ bởi $E=mc^2$; độ hụt khối là phần chênh giữa tổng khối lượng nuclon riêng rẽ và khối lượng hạt nhân.}
{Độ hụt khối của hạt nhân bền luôn âm.}
{\True Khi giải cần đọc đúng dữ kiện, dùng hệ đơn vị thống nhất và kiểm tra điều kiện áp dụng.}
{Có thể bỏ qua dữ kiện, đơn vị và ý nghĩa vật lí của kết quả.}
\loigiai{
\begin{itemchoice}
\item (Đúng) Khối lượng và năng lượng liên hệ bởi $E=mc^2$; độ hụt khối là phần chênh giữa tổng khối lượng nuclon riêng rẽ và khối lượng hạt nhân. (Vì): Khối lượng và năng lượng liên hệ bởi $E=mc^2$; độ hụt khối là phần chênh giữa tổng khối lượng nuclon riêng rẽ và khối lượng hạt nhân. b) Sai: Độ hụt khối của hạt nhân bền luôn âm. c) Đúng: Khi giải cần đọc đúng dữ kiện, dùng hệ đơn vị thống nhất và kiểm tra điều kiện áp dụng. d) Sai: Có thể bỏ qua dữ kiện, đơn vị và ý nghĩa vật lí của kết quả. Kiến thức cốt lõi: Khối lượng và năng lượng liên hệ bởi $E=mc^2$; độ hụt khối là phần chênh giữa tổng khối lượng nuclon riêng rẽ và khối lượng hạt nhân
\item (Sai) Độ hụt khối của hạt nhân bền luôn âm.
\item (Đúng) Khi giải cần đọc đúng dữ kiện, dùng hệ đơn vị thống nhất và kiểm tra điều kiện áp dụng.
\item (Sai) Có thể bỏ qua dữ kiện, đơn vị và ý nghĩa vật lí của kết quả.
\end{itemchoice}
}
\end{ex}

% ===== Câu 102 =====
% ID: L12C4B21-26-TL
% Mức: TH
\dangbt{Tính kích thước, bán kính và khối lượng riêng của hạt nhân}
\begin{bt}
% ID: L12C4B21-26-TL
% Mức: TH
Giải thích vì sao nhận định sau đúng: “Bán kính hạt nhân gần đúng $R=R_0A^{1/3}$ nên thể tích tỉ lệ với $A$ và khối lượng riêng hạt nhân xấp xỉ không đổi.”
\loigiai{
Cần nêu được: Bán kính hạt nhân gần đúng $R=R_0A^{1/3}$ nên thể tích tỉ lệ với $A$ và khối lượng riêng hạt nhân xấp xỉ không đổi. Khi vận dụng phải xác định đúng dữ kiện, điều kiện áp dụng, hệ đơn vị và kiểm tra tính hợp lí của kết quả. Nhận định “Bán kính hạt nhân tỉ lệ thuận với số khối $A$.” sai vì trái với kiến thức trên.
}
\end{bt}

% ===== Câu 103 =====
% ID: L12C4B21-26-TLN
% Mức: TH
\begin{ex}
% ID: L12C4B21-26-TLN
% Mức: TH
Trong bốn phát biểu sau về tính kích thước, bán kính và khối lượng riêng của hạt nhân, có bao nhiêu phát biểu đúng? Chỉ ghi một số nguyên. (1) Bán kính hạt nhân tỉ lệ thuận với số khối $A$. (2) Bán kính hạt nhân gần đúng $R=R_0A^{1/3}$ nên thể tích tỉ lệ với $A$ và khối lượng riêng hạt nhân xấp xỉ không đổi. (3) Kết quả phải phù hợp ý nghĩa vật lí. (4) Mọi đại lượng đều có cùng bản chất.
\shortans{2}
\loigiai{
Đối chiếu từng phát biểu với kiến thức cốt lõi: Bán kính hạt nhân gần đúng $R=R_0A^{1/3}$ nên thể tích tỉ lệ với $A$ và khối lượng riêng hạt nhân xấp xỉ không đổi. Có 2 phát biểu đúng.
}
\end{ex}

% ===== Câu 104 =====
% ID: L12C4B21-26-TN
% Mức: TH
\dangbt{Viết kí hiệu hạt nhân và xác định số proton, neutron, electron}
\begin{ex}
% ID: L12C4B21-26-TN
% Mức: TH
Kết luận nào của học sinh là đúng về viết kí hiệu hạt nhân và xác định số proton, neutron, electron?
\choice
{Trong kí hiệu $^{A}_{Z}X$, số neutron bằng $Z$.}
{\True Với hạt nhân $^{A}_{Z}X$, số proton là $Z$, số neutron là $A-Z$; nguyên tử trung hòa có $Z$ electron.}
{Có thể bỏ qua điều kiện áp dụng và đơn vị trong mọi trường hợp.}
{Kết quả của bài toán không phụ thuộc dữ kiện đã cho.}
\loigiai{
Với hạt nhân $^{A}_{Z}X$, số proton là $Z$, số neutron là $A-Z$; nguyên tử trung hòa có $Z$ electron. Vì vậy phương án $B$ đúng; các phương án còn lại trái với kiến thức hoặc điều kiện áp dụng.
}
\end{ex}

% ===== Câu 105 =====
% ID: L12C4B21-27-DS
% Mức: TH
\dangbt{Vận dụng hệ thức khối lượng – năng lượng và độ hụt khối}
\begin{ex}
% ID: L12C4B21-27-DS
% Mức: TH
Xét vận dụng hệ thức khối lượng – năng lượng và độ hụt khối (tình huống 2). Hãy xác định tính đúng, sai của bốn phát biểu sau.
\choiceTF
{Độ hụt khối của hạt nhân bền luôn âm.}
{\True Khối lượng và năng lượng liên hệ bởi $E=mc^2$; độ hụt khối là phần chênh giữa tổng khối lượng nuclon riêng rẽ và khối lượng hạt nhân.}
{Có thể bỏ qua dữ kiện, đơn vị và ý nghĩa vật lí của kết quả.}
{\True Khi giải cần đọc đúng dữ kiện, dùng hệ đơn vị thống nhất và kiểm tra điều kiện áp dụng.}
\loigiai{
\begin{itemchoice}
\item (Sai) Độ hụt khối của hạt nhân bền luôn âm. (Vì): Độ hụt khối của hạt nhân bền luôn âm. b) Đúng: Khối lượng và năng lượng liên hệ bởi $E=mc^2$; độ hụt khối là phần chênh giữa tổng khối lượng nuclon riêng rẽ và khối lượng hạt nhân. c) Sai: Có thể bỏ qua dữ kiện, đơn vị và ý nghĩa vật lí của kết quả. d) Đúng: Khi giải cần đọc đúng dữ kiện, dùng hệ đơn vị thống nhất và kiểm tra điều kiện áp dụng. Kiến thức cốt lõi: Khối lượng và năng lượng liên hệ bởi $E=mc^2$; độ hụt khối là phần chênh giữa tổng khối lượng nuclon riêng rẽ và khối lượng hạt nhân
\item (Đúng) Khối lượng và năng lượng liên hệ bởi $E=mc^2$; độ hụt khối là phần chênh giữa tổng khối lượng nuclon riêng rẽ và khối lượng hạt nhân.
\item (Sai) Có thể bỏ qua dữ kiện, đơn vị và ý nghĩa vật lí của kết quả.
\item (Đúng) Khi giải cần đọc đúng dữ kiện, dùng hệ đơn vị thống nhất và kiểm tra điều kiện áp dụng.
\end{itemchoice}
}
\end{ex}

% ===== Câu 106 =====
% ID: L12C4B21-27-TL
% Mức: VD
\dangbt{Tính kích thước, bán kính và khối lượng riêng của hạt nhân}
\begin{bt}
% ID: L12C4B21-27-TL
% Mức: VD
Phân tích sai lầm trong nhận định: “Bán kính hạt nhân tỉ lệ thuận với số khối $A$.”
\loigiai{
Cần nêu được: Bán kính hạt nhân gần đúng $R=R_0A^{1/3}$ nên thể tích tỉ lệ với $A$ và khối lượng riêng hạt nhân xấp xỉ không đổi. Khi vận dụng phải xác định đúng dữ kiện, điều kiện áp dụng, hệ đơn vị và kiểm tra tính hợp lí của kết quả. Nhận định “Bán kính hạt nhân tỉ lệ thuận với số khối $A$.” sai vì trái với kiến thức trên.
}
\end{bt}

% ===== Câu 107 =====
% ID: L12C4B21-27-TLN
% Mức: TH
\begin{ex}
% ID: L12C4B21-27-TLN
% Mức: TH
Trong bốn phát biểu sau về tính kích thước, bán kính và khối lượng riêng của hạt nhân, có bao nhiêu phát biểu đúng? Chỉ ghi một số nguyên. (1) Cần đổi các đại lượng về hệ đơn vị thống nhất. (2) Bán kính hạt nhân gần đúng $R=R_0A^{1/3}$ nên thể tích tỉ lệ với $A$ và khối lượng riêng hạt nhân xấp xỉ không đổi. (3) Bán kính hạt nhân tỉ lệ thuận với số khối $A$. (4) Không cần kiểm tra dữ kiện.
\shortans{2}
\loigiai{
Đối chiếu từng phát biểu với kiến thức cốt lõi: Bán kính hạt nhân gần đúng $R=R_0A^{1/3}$ nên thể tích tỉ lệ với $A$ và khối lượng riêng hạt nhân xấp xỉ không đổi. Có 2 phát biểu đúng.
}
\end{ex}

% ===== Câu 108 =====
% ID: L12C4B21-27-TN
% Mức: TH
\dangbt{Viết kí hiệu hạt nhân và xác định số proton, neutron, electron}
\begin{ex}
% ID: L12C4B21-27-TN
% Mức: TH
Chọn phát biểu không trái với kiến thức về viết kí hiệu hạt nhân và xác định số proton, neutron, electron?
\choice
{Trong kí hiệu $^{A}_{Z}X$, số neutron bằng $Z$.}
{Có thể bỏ qua điều kiện áp dụng và đơn vị trong mọi trường hợp.}
{\True Với hạt nhân $^{A}_{Z}X$, số proton là $Z$, số neutron là $A-Z$; nguyên tử trung hòa có $Z$ electron.}
{Kết quả của bài toán không phụ thuộc dữ kiện đã cho.}
\loigiai{
Với hạt nhân $^{A}_{Z}X$, số proton là $Z$, số neutron là $A-Z$; nguyên tử trung hòa có $Z$ electron. Vì vậy phương án $C$ đúng; các phương án còn lại trái với kiến thức hoặc điều kiện áp dụng.
}
\end{ex}

% ===== Câu 109 =====
% ID: L12C4B21-28-DS
% Mức: VD
\dangbt{Vận dụng hệ thức khối lượng – năng lượng và độ hụt khối}
\begin{ex}
% ID: L12C4B21-28-DS
% Mức: VD
Xét vận dụng hệ thức khối lượng – năng lượng và độ hụt khối (tình huống 3). Hãy xác định tính đúng, sai của bốn phát biểu sau.
\choiceTF
{\True Khối lượng và năng lượng liên hệ bởi $E=mc^2$; độ hụt khối là phần chênh giữa tổng khối lượng nuclon riêng rẽ và khối lượng hạt nhân.}
{Độ hụt khối của hạt nhân bền luôn âm.}
{\True Khi giải cần đọc đúng dữ kiện, dùng hệ đơn vị thống nhất và kiểm tra điều kiện áp dụng.}
{Có thể bỏ qua dữ kiện, đơn vị và ý nghĩa vật lí của kết quả.}
\loigiai{
\begin{itemchoice}
\item (Đúng) Khối lượng và năng lượng liên hệ bởi $E=mc^2$; độ hụt khối là phần chênh giữa tổng khối lượng nuclon riêng rẽ và khối lượng hạt nhân. (Vì): Khối lượng và năng lượng liên hệ bởi $E=mc^2$; độ hụt khối là phần chênh giữa tổng khối lượng nuclon riêng rẽ và khối lượng hạt nhân. b) Sai: Độ hụt khối của hạt nhân bền luôn âm. c) Đúng: Khi giải cần đọc đúng dữ kiện, dùng hệ đơn vị thống nhất và kiểm tra điều kiện áp dụng. d) Sai: Có thể bỏ qua dữ kiện, đơn vị và ý nghĩa vật lí của kết quả. Kiến thức cốt lõi: Khối lượng và năng lượng liên hệ bởi $E=mc^2$; độ hụt khối là phần chênh giữa tổng khối lượng nuclon riêng rẽ và khối lượng hạt nhân
\item (Sai) Độ hụt khối của hạt nhân bền luôn âm.
\item (Đúng) Khi giải cần đọc đúng dữ kiện, dùng hệ đơn vị thống nhất và kiểm tra điều kiện áp dụng.
\item (Sai) Có thể bỏ qua dữ kiện, đơn vị và ý nghĩa vật lí của kết quả.
\end{itemchoice}
}
\end{ex}

% ===== Câu 110 =====
% ID: L12C4B21-28-TL
% Mức: VD
\dangbt{Tính kích thước, bán kính và khối lượng riêng của hạt nhân}
\begin{bt}
% ID: L12C4B21-28-TL
% Mức: VD
Đề xuất các bước giải một bài toán thuộc dạng tính kích thước, bán kính và khối lượng riêng của hạt nhân.
\loigiai{
Cần nêu được: Bán kính hạt nhân gần đúng $R=R_0A^{1/3}$ nên thể tích tỉ lệ với $A$ và khối lượng riêng hạt nhân xấp xỉ không đổi. Khi vận dụng phải xác định đúng dữ kiện, điều kiện áp dụng, hệ đơn vị và kiểm tra tính hợp lí của kết quả. Nhận định “Bán kính hạt nhân tỉ lệ thuận với số khối $A$.” sai vì trái với kiến thức trên.
}
\end{bt}

% ===== Câu 111 =====
% ID: L12C4B21-28-TLN
% Mức: VD
\begin{ex}
% ID: L12C4B21-28-TLN
% Mức: VD
Trong bốn phát biểu sau về tính kích thước, bán kính và khối lượng riêng của hạt nhân, có bao nhiêu phát biểu đúng? Chỉ ghi một số nguyên. (1) Bán kính hạt nhân gần đúng $R=R_0A^{1/3}$ nên thể tích tỉ lệ với $A$ và khối lượng riêng hạt nhân xấp xỉ không đổi. (2) Cần đối chiếu kết quả với điều kiện bài toán. (3) Bán kính hạt nhân tỉ lệ thuận với số khối $A$. (4) Mọi trường hợp đều cho cùng một kết quả.
\shortans{2}
\loigiai{
Đối chiếu từng phát biểu với kiến thức cốt lõi: Bán kính hạt nhân gần đúng $R=R_0A^{1/3}$ nên thể tích tỉ lệ với $A$ và khối lượng riêng hạt nhân xấp xỉ không đổi. Có 2 phát biểu đúng.
}
\end{ex}

% ===== Câu 112 =====
% ID: L12C4B21-28-TN
% Mức: VD
\dangbt{Viết kí hiệu hạt nhân và xác định số proton, neutron, electron}
\begin{ex}
% ID: L12C4B21-28-TN
% Mức: VD
Cơ sở Vật lí đúng để giải bài là gì về viết kí hiệu hạt nhân và xác định số proton, neutron, electron?
\choice
{Trong kí hiệu $^{A}_{Z}X$, số neutron bằng $Z$.}
{Có thể bỏ qua điều kiện áp dụng và đơn vị trong mọi trường hợp.}
{Kết quả của bài toán không phụ thuộc dữ kiện đã cho.}
{\True Với hạt nhân $^{A}_{Z}X$, số proton là $Z$, số neutron là $A-Z$; nguyên tử trung hòa có $Z$ electron.}
\loigiai{
Với hạt nhân $^{A}_{Z}X$, số proton là $Z$, số neutron là $A-Z$; nguyên tử trung hòa có $Z$ electron. Vì vậy phương án $D$ đúng; các phương án còn lại trái với kiến thức hoặc điều kiện áp dụng.
}
\end{ex}

% ===== Câu 113 =====
% ID: L12C4B21-29-DS
% Mức: VD
\dangbt{Vận dụng hệ thức khối lượng – năng lượng và độ hụt khối}
\begin{ex}
% ID: L12C4B21-29-DS
% Mức: VD
Xét vận dụng hệ thức khối lượng – năng lượng và độ hụt khối (tình huống 4). Hãy xác định tính đúng, sai của bốn phát biểu sau.
\choiceTF
{Độ hụt khối của hạt nhân bền luôn âm.}
{\True Khối lượng và năng lượng liên hệ bởi $E=mc^2$; độ hụt khối là phần chênh giữa tổng khối lượng nuclon riêng rẽ và khối lượng hạt nhân.}
{Có thể bỏ qua dữ kiện, đơn vị và ý nghĩa vật lí của kết quả.}
{\True Khi giải cần đọc đúng dữ kiện, dùng hệ đơn vị thống nhất và kiểm tra điều kiện áp dụng.}
\loigiai{
\begin{itemchoice}
\item (Sai) Độ hụt khối của hạt nhân bền luôn âm. (Vì): Độ hụt khối của hạt nhân bền luôn âm. b) Đúng: Khối lượng và năng lượng liên hệ bởi $E=mc^2$; độ hụt khối là phần chênh giữa tổng khối lượng nuclon riêng rẽ và khối lượng hạt nhân. c) Sai: Có thể bỏ qua dữ kiện, đơn vị và ý nghĩa vật lí của kết quả. d) Đúng: Khi giải cần đọc đúng dữ kiện, dùng hệ đơn vị thống nhất và kiểm tra điều kiện áp dụng. Kiến thức cốt lõi: Khối lượng và năng lượng liên hệ bởi $E=mc^2$; độ hụt khối là phần chênh giữa tổng khối lượng nuclon riêng rẽ và khối lượng hạt nhân
\item (Đúng) Khối lượng và năng lượng liên hệ bởi $E=mc^2$; độ hụt khối là phần chênh giữa tổng khối lượng nuclon riêng rẽ và khối lượng hạt nhân.
\item (Sai) Có thể bỏ qua dữ kiện, đơn vị và ý nghĩa vật lí của kết quả.
\item (Đúng) Khi giải cần đọc đúng dữ kiện, dùng hệ đơn vị thống nhất và kiểm tra điều kiện áp dụng.
\end{itemchoice}
}
\end{ex}

% ===== Câu 114 =====
% ID: L12C4B21-29-TL
% Mức: VD
\dangbt{Tính kích thước, bán kính và khối lượng riêng của hạt nhân}
\begin{bt}
% ID: L12C4B21-29-TL
% Mức: VD
Nêu một tình huống thực tiễn liên quan đến tính kích thước, bán kính và khối lượng riêng của hạt nhân và giải thích bằng kiến thức Vật lí.
\loigiai{
Cần nêu được: Bán kính hạt nhân gần đúng $R=R_0A^{1/3}$ nên thể tích tỉ lệ với $A$ và khối lượng riêng hạt nhân xấp xỉ không đổi. Khi vận dụng phải xác định đúng dữ kiện, điều kiện áp dụng, hệ đơn vị và kiểm tra tính hợp lí của kết quả. Nhận định “Bán kính hạt nhân tỉ lệ thuận với số khối $A$.” sai vì trái với kiến thức trên.
}
\end{bt}

% ===== Câu 115 =====
% ID: L12C4B21-29-TLN
% Mức: VD
\begin{ex}
% ID: L12C4B21-29-TLN
% Mức: VD
Trong bốn phát biểu sau về tính kích thước, bán kính và khối lượng riêng của hạt nhân, có bao nhiêu phát biểu đúng? Chỉ ghi một số nguyên. (1) Bán kính hạt nhân tỉ lệ thuận với số khối $A$. (2) Cần đọc đúng dữ kiện và giả thiết. (3) Bán kính hạt nhân gần đúng $R=R_0A^{1/3}$ nên thể tích tỉ lệ với $A$ và khối lượng riêng hạt nhân xấp xỉ không đổi. (4) Có thể thay số trước khi chọn công thức.
\shortans{2}
\loigiai{
Đối chiếu từng phát biểu với kiến thức cốt lõi: Bán kính hạt nhân gần đúng $R=R_0A^{1/3}$ nên thể tích tỉ lệ với $A$ và khối lượng riêng hạt nhân xấp xỉ không đổi. Có 2 phát biểu đúng.
}
\end{ex}

% ===== Câu 116 =====
% ID: L12C4B21-29-TN
% Mức: VD
\dangbt{Viết kí hiệu hạt nhân và xác định số proton, neutron, electron}
\begin{ex}
% ID: L12C4B21-29-TN
% Mức: VD
Trong một bài thực hành, nhận xét nào đúng về viết kí hiệu hạt nhân và xác định số proton, neutron, electron?
\choice
{\True Với hạt nhân $^{A}_{Z}X$, số proton là $Z$, số neutron là $A-Z$; nguyên tử trung hòa có $Z$ electron.}
{Trong kí hiệu $^{A}_{Z}X$, số neutron bằng $Z$.}
{Có thể bỏ qua điều kiện áp dụng và đơn vị trong mọi trường hợp.}
{Kết quả của bài toán không phụ thuộc dữ kiện đã cho.}
\loigiai{
Với hạt nhân $^{A}_{Z}X$, số proton là $Z$, số neutron là $A-Z$; nguyên tử trung hòa có $Z$ electron. Vì vậy phương án $A$ đúng; các phương án còn lại trái với kiến thức hoặc điều kiện áp dụng.
}
\end{ex}

% ===== Câu 117 =====
% ID: L12C4B21-30-DS
% Mức: VD
\dangbt{Vận dụng hệ thức khối lượng – năng lượng và độ hụt khối}
\begin{ex}
% ID: L12C4B21-30-DS
% Mức: VD
Xét vận dụng hệ thức khối lượng – năng lượng và độ hụt khối (tình huống 5). Hãy xác định tính đúng, sai của bốn phát biểu sau.
\choiceTF
{\True Khối lượng và năng lượng liên hệ bởi $E=mc^2$; độ hụt khối là phần chênh giữa tổng khối lượng nuclon riêng rẽ và khối lượng hạt nhân.}
{Có thể bỏ qua dữ kiện, đơn vị và ý nghĩa vật lí của kết quả.}
{Độ hụt khối của hạt nhân bền luôn âm.}
{Có thể bỏ qua dữ kiện, đơn vị và ý nghĩa vật lí của kết quả.}
\loigiai{
\begin{itemchoice}
\item (Đúng) Khối lượng và năng lượng liên hệ bởi $E=mc^2$; độ hụt khối là phần chênh giữa tổng khối lượng nuclon riêng rẽ và khối lượng hạt nhân. (Vì): Khối lượng và năng lượng liên hệ bởi $E=mc^2$; độ hụt khối là phần chênh giữa tổng khối lượng nuclon riêng rẽ và khối lượng hạt nhân. b) Sai: Có thể bỏ qua dữ kiện, đơn vị và ý nghĩa vật lí của kết quả. c) Sai: Độ hụt khối của hạt nhân bền luôn âm. d) Sai: Có thể bỏ qua dữ kiện, đơn vị và ý nghĩa vật lí của kết quả. Kiến thức cốt lõi: Khối lượng và năng lượng liên hệ bởi $E=mc^2$; độ hụt khối là phần chênh giữa tổng khối lượng nuclon riêng rẽ và khối lượng hạt nhân
\item (Sai) Có thể bỏ qua dữ kiện, đơn vị và ý nghĩa vật lí của kết quả.
\item (Sai) Độ hụt khối của hạt nhân bền luôn âm.
\item (Sai) Có thể bỏ qua dữ kiện, đơn vị và ý nghĩa vật lí của kết quả.
\end{itemchoice}
}
\end{ex}

% ===== Câu 118 =====
% ID: L12C4B21-30-TL
% Mức: VDC
\dangbt{Tính kích thước, bán kính và khối lượng riêng của hạt nhân}
\begin{bt}
% ID: L12C4B21-30-TL
% Mức: VDC
So sánh nhận định đúng “Bán kính hạt nhân gần đúng $R=R_0A^{1/3}$ nên thể tích tỉ lệ với $A$ và khối lượng riêng hạt nhân xấp xỉ không đổi.” với nhận định sai “Bán kính hạt nhân tỉ lệ thuận với số khối $A$.”; chỉ rõ căn cứ Vật lí.
\loigiai{
Cần nêu được: Bán kính hạt nhân gần đúng $R=R_0A^{1/3}$ nên thể tích tỉ lệ với $A$ và khối lượng riêng hạt nhân xấp xỉ không đổi. Khi vận dụng phải xác định đúng dữ kiện, điều kiện áp dụng, hệ đơn vị và kiểm tra tính hợp lí của kết quả. Nhận định “Bán kính hạt nhân tỉ lệ thuận với số khối $A$.” sai vì trái với kiến thức trên.
}
\end{bt}

% ===== Câu 119 =====
% ID: L12C4B21-30-TLN
% Mức: VDC
\begin{ex}
% ID: L12C4B21-30-TLN
% Mức: VDC
Trong bốn phát biểu sau về tính kích thước, bán kính và khối lượng riêng của hạt nhân, có bao nhiêu phát biểu đúng? Chỉ ghi một số nguyên. (1) Kết quả cần có ý nghĩa vật lí. (2) Bán kính hạt nhân tỉ lệ thuận với số khối $A$. (3) Phải dùng đúng định luật cho tình huống. (4) Bán kính hạt nhân gần đúng $R=R_0A^{1/3}$ nên thể tích tỉ lệ với $A$ và khối lượng riêng hạt nhân xấp xỉ không đổi.
\shortans{3}
\loigiai{
Đối chiếu từng phát biểu với kiến thức cốt lõi: Bán kính hạt nhân gần đúng $R=R_0A^{1/3}$ nên thể tích tỉ lệ với $A$ và khối lượng riêng hạt nhân xấp xỉ không đổi. Có 3 phát biểu đúng.
}
\end{ex}

% ===== Câu 120 =====
% ID: L12C4B21-30-TN
% Mức: VD
\dangbt{Viết kí hiệu hạt nhân và xác định số proton, neutron, electron}
\begin{ex}
% ID: L12C4B21-30-TN
% Mức: VD
Kết luận đúng khi vận dụng là về viết kí hiệu hạt nhân và xác định số proton, neutron, electron?
\choice
{Trong kí hiệu $^{A}_{Z}X$, số neutron bằng $Z$.}
{\True Với hạt nhân $^{A}_{Z}X$, số proton là $Z$, số neutron là $A-Z$; nguyên tử trung hòa có $Z$ electron.}
{Có thể bỏ qua điều kiện áp dụng và đơn vị trong mọi trường hợp.}
{Kết quả của bài toán không phụ thuộc dữ kiện đã cho.}
\loigiai{
Với hạt nhân $^{A}_{Z}X$, số proton là $Z$, số neutron là $A-Z$; nguyên tử trung hòa có $Z$ electron. Vì vậy phương án $B$ đúng; các phương án còn lại trái với kiến thức hoặc điều kiện áp dụng.
}
\end{ex}

% ===== Câu 121 =====
% ID: L12C4B21-31-DS
% Mức: TH
\dangbt{Bài toán tổng hợp về cấu trúc và đặc trưng của hạt nhân}
\begin{ex}
% ID: L12C4B21-31-DS
% Mức: TH
Xét bài toán tổng hợp về cấu trúc và đặc trưng của hạt nhân (tình huống 1). Hãy xác định tính đúng, sai của bốn phát biểu sau.
\choiceTF
{\True Khi giải bài tổng hợp phải xác định đúng $A,Z,N$, đổi đơn vị khối lượng và dùng nhất quán $E=mc^2$ hoặc $1\,u\,c^2=931,5\,MeV$.}
{Có thể đồng nhất số khối với khối lượng hạt nhân tính theo đơn vị $u$ trong mọi phép tính chính xác.}
{\True Khi giải cần đọc đúng dữ kiện, dùng hệ đơn vị thống nhất và kiểm tra điều kiện áp dụng.}
{Có thể bỏ qua dữ kiện, đơn vị và ý nghĩa vật lí của kết quả.}
\loigiai{
\begin{itemchoice}
\item (Đúng) Khi giải bài tổng hợp phải xác định đúng $A,Z,N$, đổi đơn vị khối lượng và dùng nhất quán $E=mc^2$ hoặc $1\,u\,c^2=931,5\,MeV$. (Vì): Khi giải bài tổng hợp phải xác định đúng $A,Z,N$, đổi đơn vị khối lượng và dùng nhất quán $E=mc^2$ hoặc $1\,u\,c^2=931,5\,MeV$. b) Sai: Có thể đồng nhất số khối với khối lượng hạt nhân tính theo đơn vị $u$ trong mọi phép tính chính xác. c) Đúng: Khi giải cần đọc đúng dữ kiện, dùng hệ đơn vị thống nhất và kiểm tra điều kiện áp dụng. d) Sai: Có thể bỏ qua dữ kiện, đơn vị và ý nghĩa vật lí của kết quả. Kiến thức cốt lõi: Khi giải bài tổng hợp phải xác định đúng $A,Z,N$, đổi đơn vị khối lượng và dùng nhất quán $E=mc^2$ hoặc $1\,u\,c^2=931,5\,MeV$
\item (Sai) Có thể đồng nhất số khối với khối lượng hạt nhân tính theo đơn vị $u$ trong mọi phép tính chính xác.
\item (Đúng) Khi giải cần đọc đúng dữ kiện, dùng hệ đơn vị thống nhất và kiểm tra điều kiện áp dụng.
\item (Sai) Có thể bỏ qua dữ kiện, đơn vị và ý nghĩa vật lí của kết quả.
\end{itemchoice}
}
\end{ex}

% ===== Câu 122 =====
% ID: L12C4B21-31-TL
% Mức: TH
\dangbt{Vận dụng hệ thức khối lượng – năng lượng và độ hụt khối}
\begin{bt}
% ID: L12C4B21-31-TL
% Mức: TH
Trình bày kiến thức cốt lõi của vận dụng hệ thức khối lượng – năng lượng và độ hụt khối và nêu điều kiện áp dụng.
\loigiai{
Cần nêu được: Khối lượng và năng lượng liên hệ bởi $E=mc^2$; độ hụt khối là phần chênh giữa tổng khối lượng nuclon riêng rẽ và khối lượng hạt nhân. Khi vận dụng phải xác định đúng dữ kiện, điều kiện áp dụng, hệ đơn vị và kiểm tra tính hợp lí của kết quả. Nhận định “Độ hụt khối của hạt nhân bền luôn âm.” sai vì trái với kiến thức trên.
}
\end{bt}

% ===== Câu 123 =====
% ID: L12C4B21-31-TLN
% Mức: TH
\begin{ex}
% ID: L12C4B21-31-TLN
% Mức: TH
Trong bốn phát biểu sau về vận dụng hệ thức khối lượng – năng lượng và độ hụt khối, có bao nhiêu phát biểu đúng? Chỉ ghi một số nguyên. (1) Khối lượng và năng lượng liên hệ bởi $E=mc^2$; độ hụt khối là phần chênh giữa tổng khối lượng nuclon riêng rẽ và khối lượng hạt nhân. (2) Độ hụt khối của hạt nhân bền luôn âm. (3) Cần kiểm tra điều kiện áp dụng trước khi tính toán. (4) Có thể bỏ qua đơn vị của kết quả.
\shortans{2}
\loigiai{
Đối chiếu từng phát biểu với kiến thức cốt lõi: Khối lượng và năng lượng liên hệ bởi $E=mc^2$; độ hụt khối là phần chênh giữa tổng khối lượng nuclon riêng rẽ và khối lượng hạt nhân. Có 2 phát biểu đúng.
}
\end{ex}

% ===== Câu 124 =====
% ID: L12C4B21-31-TN
% Mức: NB
\dangbt{Nhận biết đồng vị và tính nguyên tử khối trung bình}
\begin{ex}
% ID: L12C4B21-31-TN
% Mức: NB
Phát biểu nào sau đây đúng về nhận biết đồng vị và tính nguyên tử khối trung bình?
\choice
{\True Các đồng vị có cùng số proton nhưng khác số neutron; nguyên tử khối trung bình là trung bình có trọng số theo tỉ lệ đồng vị.}
{Các đồng vị của một nguyên tố có số proton khác nhau.}
{Có thể bỏ qua điều kiện áp dụng và đơn vị trong mọi trường hợp.}
{Kết quả của bài toán không phụ thuộc dữ kiện đã cho.}
\loigiai{
Các đồng vị có cùng số proton nhưng khác số neutron; nguyên tử khối trung bình là trung bình có trọng số theo tỉ lệ đồng vị. Vì vậy phương án $A$ đúng; các phương án còn lại trái với kiến thức hoặc điều kiện áp dụng.
}
\end{ex}

% ===== Câu 125 =====
% ID: L12C4B21-32-DS
% Mức: TH
\dangbt{Bài toán tổng hợp về cấu trúc và đặc trưng của hạt nhân}
\begin{ex}
% ID: L12C4B21-32-DS
% Mức: TH
Xét bài toán tổng hợp về cấu trúc và đặc trưng của hạt nhân (tình huống 2). Hãy xác định tính đúng, sai của bốn phát biểu sau.
\choiceTF
{Có thể đồng nhất số khối với khối lượng hạt nhân tính theo đơn vị $u$ trong mọi phép tính chính xác.}
{\True Khi giải bài tổng hợp phải xác định đúng $A,Z,N$, đổi đơn vị khối lượng và dùng nhất quán $E=mc^2$ hoặc $1\,u\,c^2=931,5\,MeV$.}
{Có thể bỏ qua dữ kiện, đơn vị và ý nghĩa vật lí của kết quả.}
{\True Khi giải cần đọc đúng dữ kiện, dùng hệ đơn vị thống nhất và kiểm tra điều kiện áp dụng.}
\loigiai{
\begin{itemchoice}
\item (Sai) Có thể đồng nhất số khối với khối lượng hạt nhân tính theo đơn vị $u$ trong mọi phép tính chính xác. (Vì): Có thể đồng nhất số khối với khối lượng hạt nhân tính theo đơn vị $u$ trong mọi phép tính chính xác. b) Đúng: Khi giải bài tổng hợp phải xác định đúng $A,Z,N$, đổi đơn vị khối lượng và dùng nhất quán $E=mc^2$ hoặc $1\,u\,c^2=931,5\,MeV$. c) Sai: Có thể bỏ qua dữ kiện, đơn vị và ý nghĩa vật lí của kết quả. d) Đúng: Khi giải cần đọc đúng dữ kiện, dùng hệ đơn vị thống nhất và kiểm tra điều kiện áp dụng. Kiến thức cốt lõi: Khi giải bài tổng hợp phải xác định đúng $A,Z,N$, đổi đơn vị khối lượng và dùng nhất quán $E=mc^2$ hoặc $1\,u\,c^2=931,5\,MeV$
\item (Đúng) Khi giải bài tổng hợp phải xác định đúng $A,Z,N$, đổi đơn vị khối lượng và dùng nhất quán $E=mc^2$ hoặc $1\,u\,c^2=931,5\,MeV$.
\item (Sai) Có thể bỏ qua dữ kiện, đơn vị và ý nghĩa vật lí của kết quả.
\item (Đúng) Khi giải cần đọc đúng dữ kiện, dùng hệ đơn vị thống nhất và kiểm tra điều kiện áp dụng.
\end{itemchoice}
}
\end{ex}

% ===== Câu 126 =====
% ID: L12C4B21-32-TL
% Mức: TH
\dangbt{Vận dụng hệ thức khối lượng – năng lượng và độ hụt khối}
\begin{bt}
% ID: L12C4B21-32-TL
% Mức: TH
Giải thích vì sao nhận định sau đúng: “Khối lượng và năng lượng liên hệ bởi $E=mc^2$; độ hụt khối là phần chênh giữa tổng khối lượng nuclon riêng rẽ và khối lượng hạt nhân.”
\loigiai{
Cần nêu được: Khối lượng và năng lượng liên hệ bởi $E=mc^2$; độ hụt khối là phần chênh giữa tổng khối lượng nuclon riêng rẽ và khối lượng hạt nhân. Khi vận dụng phải xác định đúng dữ kiện, điều kiện áp dụng, hệ đơn vị và kiểm tra tính hợp lí của kết quả. Nhận định “Độ hụt khối của hạt nhân bền luôn âm.” sai vì trái với kiến thức trên.
}
\end{bt}

% ===== Câu 127 =====
% ID: L12C4B21-32-TLN
% Mức: TH
\begin{ex}
% ID: L12C4B21-32-TLN
% Mức: TH
Trong bốn phát biểu sau về vận dụng hệ thức khối lượng – năng lượng và độ hụt khối, có bao nhiêu phát biểu đúng? Chỉ ghi một số nguyên. (1) Độ hụt khối của hạt nhân bền luôn âm. (2) Khối lượng và năng lượng liên hệ bởi $E=mc^2$; độ hụt khối là phần chênh giữa tổng khối lượng nuclon riêng rẽ và khối lượng hạt nhân. (3) Kết quả phải phù hợp ý nghĩa vật lí. (4) Mọi đại lượng đều có cùng bản chất.
\shortans{2}
\loigiai{
Đối chiếu từng phát biểu với kiến thức cốt lõi: Khối lượng và năng lượng liên hệ bởi $E=mc^2$; độ hụt khối là phần chênh giữa tổng khối lượng nuclon riêng rẽ và khối lượng hạt nhân. Có 2 phát biểu đúng.
}
\end{ex}

% ===== Câu 128 =====
% ID: L12C4B21-32-TN
% Mức: NB
\dangbt{Nhận biết đồng vị và tính nguyên tử khối trung bình}
\begin{ex}
% ID: L12C4B21-32-TN
% Mức: NB
Nhận định nào phù hợp với kiến thức Vật lí về nhận biết đồng vị và tính nguyên tử khối trung bình?
\choice
{Các đồng vị của một nguyên tố có số proton khác nhau.}
{\True Các đồng vị có cùng số proton nhưng khác số neutron; nguyên tử khối trung bình là trung bình có trọng số theo tỉ lệ đồng vị.}
{Có thể bỏ qua điều kiện áp dụng và đơn vị trong mọi trường hợp.}
{Kết quả của bài toán không phụ thuộc dữ kiện đã cho.}
\loigiai{
Các đồng vị có cùng số proton nhưng khác số neutron; nguyên tử khối trung bình là trung bình có trọng số theo tỉ lệ đồng vị. Vì vậy phương án $B$ đúng; các phương án còn lại trái với kiến thức hoặc điều kiện áp dụng.
}
\end{ex}

% ===== Câu 129 =====
% ID: L12C4B21-33-DS
% Mức: VD
\dangbt{Bài toán tổng hợp về cấu trúc và đặc trưng của hạt nhân}
\begin{ex}
% ID: L12C4B21-33-DS
% Mức: VD
Xét bài toán tổng hợp về cấu trúc và đặc trưng của hạt nhân (tình huống 3). Hãy xác định tính đúng, sai của bốn phát biểu sau.
\choiceTF
{\True Khi giải bài tổng hợp phải xác định đúng $A,Z,N$, đổi đơn vị khối lượng và dùng nhất quán $E=mc^2$ hoặc $1\,u\,c^2=931,5\,MeV$.}
{Có thể đồng nhất số khối với khối lượng hạt nhân tính theo đơn vị $u$ trong mọi phép tính chính xác.}
{\True Khi giải cần đọc đúng dữ kiện, dùng hệ đơn vị thống nhất và kiểm tra điều kiện áp dụng.}
{Có thể bỏ qua dữ kiện, đơn vị và ý nghĩa vật lí của kết quả.}
\loigiai{
\begin{itemchoice}
\item (Đúng) Khi giải bài tổng hợp phải xác định đúng $A,Z,N$, đổi đơn vị khối lượng và dùng nhất quán $E=mc^2$ hoặc $1\,u\,c^2=931,5\,MeV$. (Vì): Khi giải bài tổng hợp phải xác định đúng $A,Z,N$, đổi đơn vị khối lượng và dùng nhất quán $E=mc^2$ hoặc $1\,u\,c^2=931,5\,MeV$. b) Sai: Có thể đồng nhất số khối với khối lượng hạt nhân tính theo đơn vị $u$ trong mọi phép tính chính xác. c) Đúng: Khi giải cần đọc đúng dữ kiện, dùng hệ đơn vị thống nhất và kiểm tra điều kiện áp dụng. d) Sai: Có thể bỏ qua dữ kiện, đơn vị và ý nghĩa vật lí của kết quả. Kiến thức cốt lõi: Khi giải bài tổng hợp phải xác định đúng $A,Z,N$, đổi đơn vị khối lượng và dùng nhất quán $E=mc^2$ hoặc $1\,u\,c^2=931,5\,MeV$
\item (Sai) Có thể đồng nhất số khối với khối lượng hạt nhân tính theo đơn vị $u$ trong mọi phép tính chính xác.
\item (Đúng) Khi giải cần đọc đúng dữ kiện, dùng hệ đơn vị thống nhất và kiểm tra điều kiện áp dụng.
\item (Sai) Có thể bỏ qua dữ kiện, đơn vị và ý nghĩa vật lí của kết quả.
\end{itemchoice}
}
\end{ex}

% ===== Câu 130 =====
% ID: L12C4B21-33-TL
% Mức: VD
\dangbt{Vận dụng hệ thức khối lượng – năng lượng và độ hụt khối}
\begin{bt}
% ID: L12C4B21-33-TL
% Mức: VD
Phân tích sai lầm trong nhận định: “Độ hụt khối của hạt nhân bền luôn âm.”
\loigiai{
Cần nêu được: Khối lượng và năng lượng liên hệ bởi $E=mc^2$; độ hụt khối là phần chênh giữa tổng khối lượng nuclon riêng rẽ và khối lượng hạt nhân. Khi vận dụng phải xác định đúng dữ kiện, điều kiện áp dụng, hệ đơn vị và kiểm tra tính hợp lí của kết quả. Nhận định “Độ hụt khối của hạt nhân bền luôn âm.” sai vì trái với kiến thức trên.
}
\end{bt}

% ===== Câu 131 =====
% ID: L12C4B21-33-TLN
% Mức: TH
\begin{ex}
% ID: L12C4B21-33-TLN
% Mức: TH
Trong bốn phát biểu sau về vận dụng hệ thức khối lượng – năng lượng và độ hụt khối, có bao nhiêu phát biểu đúng? Chỉ ghi một số nguyên. (1) Cần đổi các đại lượng về hệ đơn vị thống nhất. (2) Khối lượng và năng lượng liên hệ bởi $E=mc^2$; độ hụt khối là phần chênh giữa tổng khối lượng nuclon riêng rẽ và khối lượng hạt nhân. (3) Độ hụt khối của hạt nhân bền luôn âm. (4) Không cần kiểm tra dữ kiện.
\shortans{2}
\loigiai{
Đối chiếu từng phát biểu với kiến thức cốt lõi: Khối lượng và năng lượng liên hệ bởi $E=mc^2$; độ hụt khối là phần chênh giữa tổng khối lượng nuclon riêng rẽ và khối lượng hạt nhân. Có 2 phát biểu đúng.
}
\end{ex}

% ===== Câu 132 =====
% ID: L12C4B21-33-TN
% Mức: NB
\dangbt{Nhận biết đồng vị và tính nguyên tử khối trung bình}
\begin{ex}
% ID: L12C4B21-33-TN
% Mức: NB
Chọn kết luận chính xác về nhận biết đồng vị và tính nguyên tử khối trung bình?
\choice
{Các đồng vị của một nguyên tố có số proton khác nhau.}
{Có thể bỏ qua điều kiện áp dụng và đơn vị trong mọi trường hợp.}
{\True Các đồng vị có cùng số proton nhưng khác số neutron; nguyên tử khối trung bình là trung bình có trọng số theo tỉ lệ đồng vị.}
{Kết quả của bài toán không phụ thuộc dữ kiện đã cho.}
\loigiai{
Các đồng vị có cùng số proton nhưng khác số neutron; nguyên tử khối trung bình là trung bình có trọng số theo tỉ lệ đồng vị. Vì vậy phương án $C$ đúng; các phương án còn lại trái với kiến thức hoặc điều kiện áp dụng.
}
\end{ex}

% ===== Câu 133 =====
% ID: L12C4B21-34-DS
% Mức: VD
\dangbt{Bài toán tổng hợp về cấu trúc và đặc trưng của hạt nhân}
\begin{ex}
% ID: L12C4B21-34-DS
% Mức: VD
Xét bài toán tổng hợp về cấu trúc và đặc trưng của hạt nhân (tình huống 4). Hãy xác định tính đúng, sai của bốn phát biểu sau.
\choiceTF
{Có thể đồng nhất số khối với khối lượng hạt nhân tính theo đơn vị $u$ trong mọi phép tính chính xác.}
{\True Khi giải bài tổng hợp phải xác định đúng $A,Z,N$, đổi đơn vị khối lượng và dùng nhất quán $E=mc^2$ hoặc $1\,u\,c^2=931,5\,MeV$.}
{Có thể bỏ qua dữ kiện, đơn vị và ý nghĩa vật lí của kết quả.}
{\True Khi giải cần đọc đúng dữ kiện, dùng hệ đơn vị thống nhất và kiểm tra điều kiện áp dụng.}
\loigiai{
\begin{itemchoice}
\item (Sai) Có thể đồng nhất số khối với khối lượng hạt nhân tính theo đơn vị $u$ trong mọi phép tính chính xác. (Vì): Có thể đồng nhất số khối với khối lượng hạt nhân tính theo đơn vị $u$ trong mọi phép tính chính xác. b) Đúng: Khi giải bài tổng hợp phải xác định đúng $A,Z,N$, đổi đơn vị khối lượng và dùng nhất quán $E=mc^2$ hoặc $1\,u\,c^2=931,5\,MeV$. c) Sai: Có thể bỏ qua dữ kiện, đơn vị và ý nghĩa vật lí của kết quả. d) Đúng: Khi giải cần đọc đúng dữ kiện, dùng hệ đơn vị thống nhất và kiểm tra điều kiện áp dụng. Kiến thức cốt lõi: Khi giải bài tổng hợp phải xác định đúng $A,Z,N$, đổi đơn vị khối lượng và dùng nhất quán $E=mc^2$ hoặc $1\,u\,c^2=931,5\,MeV$
\item (Đúng) Khi giải bài tổng hợp phải xác định đúng $A,Z,N$, đổi đơn vị khối lượng và dùng nhất quán $E=mc^2$ hoặc $1\,u\,c^2=931,5\,MeV$.
\item (Sai) Có thể bỏ qua dữ kiện, đơn vị và ý nghĩa vật lí của kết quả.
\item (Đúng) Khi giải cần đọc đúng dữ kiện, dùng hệ đơn vị thống nhất và kiểm tra điều kiện áp dụng.
\end{itemchoice}
}
\end{ex}

% ===== Câu 134 =====
% ID: L12C4B21-34-TL
% Mức: VD
\dangbt{Vận dụng hệ thức khối lượng – năng lượng và độ hụt khối}
\begin{bt}
% ID: L12C4B21-34-TL
% Mức: VD
Đề xuất các bước giải một bài toán thuộc dạng vận dụng hệ thức khối lượng – năng lượng và độ hụt khối.
\loigiai{
Cần nêu được: Khối lượng và năng lượng liên hệ bởi $E=mc^2$; độ hụt khối là phần chênh giữa tổng khối lượng nuclon riêng rẽ và khối lượng hạt nhân. Khi vận dụng phải xác định đúng dữ kiện, điều kiện áp dụng, hệ đơn vị và kiểm tra tính hợp lí của kết quả. Nhận định “Độ hụt khối của hạt nhân bền luôn âm.” sai vì trái với kiến thức trên.
}
\end{bt}

% ===== Câu 135 =====
% ID: L12C4B21-34-TLN
% Mức: VD
\begin{ex}
% ID: L12C4B21-34-TLN
% Mức: VD
Trong bốn phát biểu sau về vận dụng hệ thức khối lượng – năng lượng và độ hụt khối, có bao nhiêu phát biểu đúng? Chỉ ghi một số nguyên. (1) Khối lượng và năng lượng liên hệ bởi $E=mc^2$; độ hụt khối là phần chênh giữa tổng khối lượng nuclon riêng rẽ và khối lượng hạt nhân. (2) Cần đối chiếu kết quả với điều kiện bài toán. (3) Độ hụt khối của hạt nhân bền luôn âm. (4) Mọi trường hợp đều cho cùng một kết quả.
\shortans{2}
\loigiai{
Đối chiếu từng phát biểu với kiến thức cốt lõi: Khối lượng và năng lượng liên hệ bởi $E=mc^2$; độ hụt khối là phần chênh giữa tổng khối lượng nuclon riêng rẽ và khối lượng hạt nhân. Có 2 phát biểu đúng.
}
\end{ex}

% ===== Câu 136 =====
% ID: L12C4B21-34-TN
% Mức: TH
\dangbt{Nhận biết đồng vị và tính nguyên tử khối trung bình}
\begin{ex}
% ID: L12C4B21-34-TN
% Mức: TH
Nội dung nào mô tả đúng nhất về nhận biết đồng vị và tính nguyên tử khối trung bình?
\choice
{Các đồng vị của một nguyên tố có số proton khác nhau.}
{Có thể bỏ qua điều kiện áp dụng và đơn vị trong mọi trường hợp.}
{Kết quả của bài toán không phụ thuộc dữ kiện đã cho.}
{\True Các đồng vị có cùng số proton nhưng khác số neutron; nguyên tử khối trung bình là trung bình có trọng số theo tỉ lệ đồng vị.}
\loigiai{
Các đồng vị có cùng số proton nhưng khác số neutron; nguyên tử khối trung bình là trung bình có trọng số theo tỉ lệ đồng vị. Vì vậy phương án $D$ đúng; các phương án còn lại trái với kiến thức hoặc điều kiện áp dụng.
}
\end{ex}

% ===== Câu 137 =====
% ID: L12C4B21-35-DS
% Mức: VD
\dangbt{Bài toán tổng hợp về cấu trúc và đặc trưng của hạt nhân}
\begin{ex}
% ID: L12C4B21-35-DS
% Mức: VD
Xét bài toán tổng hợp về cấu trúc và đặc trưng của hạt nhân (tình huống 5). Hãy xác định tính đúng, sai của bốn phát biểu sau.
\choiceTF
{\True Khi giải bài tổng hợp phải xác định đúng $A,Z,N$, đổi đơn vị khối lượng và dùng nhất quán $E=mc^2$ hoặc $1\,u\,c^2=931,5\,MeV$.}
{Có thể bỏ qua dữ kiện, đơn vị và ý nghĩa vật lí của kết quả.}
{Có thể đồng nhất số khối với khối lượng hạt nhân tính theo đơn vị $u$ trong mọi phép tính chính xác.}
{Có thể bỏ qua dữ kiện, đơn vị và ý nghĩa vật lí của kết quả.}
\loigiai{
\begin{itemchoice}
\item (Đúng) Khi giải bài tổng hợp phải xác định đúng $A,Z,N$, đổi đơn vị khối lượng và dùng nhất quán $E=mc^2$ hoặc $1\,u\,c^2=931,5\,MeV$. (Vì): Khi giải bài tổng hợp phải xác định đúng $A,Z,N$, đổi đơn vị khối lượng và dùng nhất quán $E=mc^2$ hoặc $1\,u\,c^2=931,5\,MeV$. b) Sai: Có thể bỏ qua dữ kiện, đơn vị và ý nghĩa vật lí của kết quả. c) Sai: Có thể đồng nhất số khối với khối lượng hạt nhân tính theo đơn vị $u$ trong mọi phép tính chính xác. d) Sai: Có thể bỏ qua dữ kiện, đơn vị và ý nghĩa vật lí của kết quả. Kiến thức cốt lõi: Khi giải bài tổng hợp phải xác định đúng $A,Z,N$, đổi đơn vị khối lượng và dùng nhất quán $E=mc^2$ hoặc $1\,u\,c^2=931,5\,MeV$
\item (Sai) Có thể bỏ qua dữ kiện, đơn vị và ý nghĩa vật lí của kết quả.
\item (Sai) Có thể đồng nhất số khối với khối lượng hạt nhân tính theo đơn vị $u$ trong mọi phép tính chính xác.
\item (Sai) Có thể bỏ qua dữ kiện, đơn vị và ý nghĩa vật lí của kết quả.
\end{itemchoice}
}
\end{ex}

% ===== Câu 138 =====
% ID: L12C4B21-35-TL
% Mức: VD
\dangbt{Vận dụng hệ thức khối lượng – năng lượng và độ hụt khối}
\begin{bt}
% ID: L12C4B21-35-TL
% Mức: VD
Nêu một tình huống thực tiễn liên quan đến vận dụng hệ thức khối lượng – năng lượng và độ hụt khối và giải thích bằng kiến thức Vật lí.
\loigiai{
Cần nêu được: Khối lượng và năng lượng liên hệ bởi $E=mc^2$; độ hụt khối là phần chênh giữa tổng khối lượng nuclon riêng rẽ và khối lượng hạt nhân. Khi vận dụng phải xác định đúng dữ kiện, điều kiện áp dụng, hệ đơn vị và kiểm tra tính hợp lí của kết quả. Nhận định “Độ hụt khối của hạt nhân bền luôn âm.” sai vì trái với kiến thức trên.
}
\end{bt}

% ===== Câu 139 =====
% ID: L12C4B21-35-TLN
% Mức: VD
\begin{ex}
% ID: L12C4B21-35-TLN
% Mức: VD
Trong bốn phát biểu sau về vận dụng hệ thức khối lượng – năng lượng và độ hụt khối, có bao nhiêu phát biểu đúng? Chỉ ghi một số nguyên. (1) Độ hụt khối của hạt nhân bền luôn âm. (2) Cần đọc đúng dữ kiện và giả thiết. (3) Khối lượng và năng lượng liên hệ bởi $E=mc^2$; độ hụt khối là phần chênh giữa tổng khối lượng nuclon riêng rẽ và khối lượng hạt nhân. (4) Có thể thay số trước khi chọn công thức.
\shortans{2}
\loigiai{
Đối chiếu từng phát biểu với kiến thức cốt lõi: Khối lượng và năng lượng liên hệ bởi $E=mc^2$; độ hụt khối là phần chênh giữa tổng khối lượng nuclon riêng rẽ và khối lượng hạt nhân. Có 2 phát biểu đúng.
}
\end{ex}

% ===== Câu 140 =====
% ID: L12C4B21-35-TN
% Mức: TH
\dangbt{Nhận biết đồng vị và tính nguyên tử khối trung bình}
\begin{ex}
% ID: L12C4B21-35-TN
% Mức: TH
Khi phân tích, phát biểu nào đúng về nhận biết đồng vị và tính nguyên tử khối trung bình?
\choice
{\True Các đồng vị có cùng số proton nhưng khác số neutron; nguyên tử khối trung bình là trung bình có trọng số theo tỉ lệ đồng vị.}
{Các đồng vị của một nguyên tố có số proton khác nhau.}
{Có thể bỏ qua điều kiện áp dụng và đơn vị trong mọi trường hợp.}
{Kết quả của bài toán không phụ thuộc dữ kiện đã cho.}
\loigiai{
Các đồng vị có cùng số proton nhưng khác số neutron; nguyên tử khối trung bình là trung bình có trọng số theo tỉ lệ đồng vị. Vì vậy phương án $A$ đúng; các phương án còn lại trái với kiến thức hoặc điều kiện áp dụng.
}
\end{ex}

% ===== Câu 141 =====
% ID: L12C4B21-36-TL
% Mức: VDC
\dangbt{Vận dụng hệ thức khối lượng – năng lượng và độ hụt khối}
\begin{bt}
% ID: L12C4B21-36-TL
% Mức: VDC
So sánh nhận định đúng “Khối lượng và năng lượng liên hệ bởi $E=mc^2$; độ hụt khối là phần chênh giữa tổng khối lượng nuclon riêng rẽ và khối lượng hạt nhân.” với nhận định sai “Độ hụt khối của hạt nhân bền luôn âm.”; chỉ rõ căn cứ Vật lí.
\loigiai{
Cần nêu được: Khối lượng và năng lượng liên hệ bởi $E=mc^2$; độ hụt khối là phần chênh giữa tổng khối lượng nuclon riêng rẽ và khối lượng hạt nhân. Khi vận dụng phải xác định đúng dữ kiện, điều kiện áp dụng, hệ đơn vị và kiểm tra tính hợp lí của kết quả. Nhận định “Độ hụt khối của hạt nhân bền luôn âm.” sai vì trái với kiến thức trên.
}
\end{bt}

% ===== Câu 142 =====
% ID: L12C4B21-36-TLN
% Mức: VDC
\begin{ex}
% ID: L12C4B21-36-TLN
% Mức: VDC
Trong bốn phát biểu sau về vận dụng hệ thức khối lượng – năng lượng và độ hụt khối, có bao nhiêu phát biểu đúng? Chỉ ghi một số nguyên. (1) Kết quả cần có ý nghĩa vật lí. (2) Độ hụt khối của hạt nhân bền luôn âm. (3) Phải dùng đúng định luật cho tình huống. (4) Khối lượng và năng lượng liên hệ bởi $E=mc^2$; độ hụt khối là phần chênh giữa tổng khối lượng nuclon riêng rẽ và khối lượng hạt nhân.
\shortans{3}
\loigiai{
Đối chiếu từng phát biểu với kiến thức cốt lõi: Khối lượng và năng lượng liên hệ bởi $E=mc^2$; độ hụt khối là phần chênh giữa tổng khối lượng nuclon riêng rẽ và khối lượng hạt nhân. Có 3 phát biểu đúng.
}
\end{ex}

% ===== Câu 143 =====
% ID: L12C4B21-36-TN
% Mức: TH
\dangbt{Nhận biết đồng vị và tính nguyên tử khối trung bình}
\begin{ex}
% ID: L12C4B21-36-TN
% Mức: TH
Kết luận nào của học sinh là đúng về nhận biết đồng vị và tính nguyên tử khối trung bình?
\choice
{Các đồng vị của một nguyên tố có số proton khác nhau.}
{\True Các đồng vị có cùng số proton nhưng khác số neutron; nguyên tử khối trung bình là trung bình có trọng số theo tỉ lệ đồng vị.}
{Có thể bỏ qua điều kiện áp dụng và đơn vị trong mọi trường hợp.}
{Kết quả của bài toán không phụ thuộc dữ kiện đã cho.}
\loigiai{
Các đồng vị có cùng số proton nhưng khác số neutron; nguyên tử khối trung bình là trung bình có trọng số theo tỉ lệ đồng vị. Vì vậy phương án $B$ đúng; các phương án còn lại trái với kiến thức hoặc điều kiện áp dụng.
}
\end{ex}

% ===== Câu 144 =====
% ID: L12C4B21-37-TL
% Mức: TH
\dangbt{Bài toán tổng hợp về cấu trúc và đặc trưng của hạt nhân}
\begin{bt}
% ID: L12C4B21-37-TL
% Mức: TH
Trình bày kiến thức cốt lõi của bài toán tổng hợp về cấu trúc và đặc trưng của hạt nhân và nêu điều kiện áp dụng.
\loigiai{
Cần nêu được: Khi giải bài tổng hợp phải xác định đúng $A,Z,N$, đổi đơn vị khối lượng và dùng nhất quán $E=mc^2$ hoặc $1\,u\,c^2=931,5\,MeV$. Khi vận dụng phải xác định đúng dữ kiện, điều kiện áp dụng, hệ đơn vị và kiểm tra tính hợp lí của kết quả. Nhận định “Có thể đồng nhất số khối với khối lượng hạt nhân tính theo đơn vị $u$ trong mọi phép tính chính xác.” sai vì trái với kiến thức trên.
}
\end{bt}

% ===== Câu 145 =====
% ID: L12C4B21-37-TLN
% Mức: TH
\begin{ex}
% ID: L12C4B21-37-TLN
% Mức: TH
Trong bốn phát biểu sau về bài toán tổng hợp về cấu trúc và đặc trưng của hạt nhân, có bao nhiêu phát biểu đúng? Chỉ ghi một số nguyên. (1) Khi giải bài tổng hợp phải xác định đúng $A,Z,N$, đổi đơn vị khối lượng và dùng nhất quán $E=mc^2$ hoặc $1\,u\,c^2=931,5\,MeV$. (2) Có thể đồng nhất số khối với khối lượng hạt nhân tính theo đơn vị $u$ trong mọi phép tính chính xác. (3) Cần kiểm tra điều kiện áp dụng trước khi tính toán. (4) Có thể bỏ qua đơn vị của kết quả.
\shortans{2}
\loigiai{
Đối chiếu từng phát biểu với kiến thức cốt lõi: Khi giải bài tổng hợp phải xác định đúng $A,Z,N$, đổi đơn vị khối lượng và dùng nhất quán $E=mc^2$ hoặc $1\,u\,c^2=931,5\,MeV$. Có 2 phát biểu đúng.
}
\end{ex}

% ===== Câu 146 =====
% ID: L12C4B21-37-TN
% Mức: TH
\dangbt{Nhận biết đồng vị và tính nguyên tử khối trung bình}
\begin{ex}
% ID: L12C4B21-37-TN
% Mức: TH
Chọn phát biểu không trái với kiến thức về nhận biết đồng vị và tính nguyên tử khối trung bình?
\choice
{Các đồng vị của một nguyên tố có số proton khác nhau.}
{Có thể bỏ qua điều kiện áp dụng và đơn vị trong mọi trường hợp.}
{\True Các đồng vị có cùng số proton nhưng khác số neutron; nguyên tử khối trung bình là trung bình có trọng số theo tỉ lệ đồng vị.}
{Kết quả của bài toán không phụ thuộc dữ kiện đã cho.}
\loigiai{
Các đồng vị có cùng số proton nhưng khác số neutron; nguyên tử khối trung bình là trung bình có trọng số theo tỉ lệ đồng vị. Vì vậy phương án $C$ đúng; các phương án còn lại trái với kiến thức hoặc điều kiện áp dụng.
}
\end{ex}

% ===== Câu 147 =====
% ID: L12C4B21-38-TL
% Mức: TH
\dangbt{Bài toán tổng hợp về cấu trúc và đặc trưng của hạt nhân}
\begin{bt}
% ID: L12C4B21-38-TL
% Mức: TH
Giải thích vì sao nhận định sau đúng: “Khi giải bài tổng hợp phải xác định đúng $A,Z,N$, đổi đơn vị khối lượng và dùng nhất quán $E=mc^2$ hoặc $1\,u\,c^2=931,5\,MeV$.”
\loigiai{
Cần nêu được: Khi giải bài tổng hợp phải xác định đúng $A,Z,N$, đổi đơn vị khối lượng và dùng nhất quán $E=mc^2$ hoặc $1\,u\,c^2=931,5\,MeV$. Khi vận dụng phải xác định đúng dữ kiện, điều kiện áp dụng, hệ đơn vị và kiểm tra tính hợp lí của kết quả. Nhận định “Có thể đồng nhất số khối với khối lượng hạt nhân tính theo đơn vị $u$ trong mọi phép tính chính xác.” sai vì trái với kiến thức trên.
}
\end{bt}

% ===== Câu 148 =====
% ID: L12C4B21-38-TLN
% Mức: TH
\begin{ex}
% ID: L12C4B21-38-TLN
% Mức: TH
Trong bốn phát biểu sau về bài toán tổng hợp về cấu trúc và đặc trưng của hạt nhân, có bao nhiêu phát biểu đúng? Chỉ ghi một số nguyên. (1) Có thể đồng nhất số khối với khối lượng hạt nhân tính theo đơn vị $u$ trong mọi phép tính chính xác. (2) Khi giải bài tổng hợp phải xác định đúng $A,Z,N$, đổi đơn vị khối lượng và dùng nhất quán $E=mc^2$ hoặc $1\,u\,c^2=931,5\,MeV$. (3) Kết quả phải phù hợp ý nghĩa vật lí. (4) Mọi đại lượng đều có cùng bản chất.
\shortans{2}
\loigiai{
Đối chiếu từng phát biểu với kiến thức cốt lõi: Khi giải bài tổng hợp phải xác định đúng $A,Z,N$, đổi đơn vị khối lượng và dùng nhất quán $E=mc^2$ hoặc $1\,u\,c^2=931,5\,MeV$. Có 2 phát biểu đúng.
}
\end{ex}

% ===== Câu 149 =====
% ID: L12C4B21-38-TN
% Mức: VD
\dangbt{Nhận biết đồng vị và tính nguyên tử khối trung bình}
\begin{ex}
% ID: L12C4B21-38-TN
% Mức: VD
Cơ sở Vật lí đúng để giải bài là gì về nhận biết đồng vị và tính nguyên tử khối trung bình?
\choice
{Các đồng vị của một nguyên tố có số proton khác nhau.}
{Có thể bỏ qua điều kiện áp dụng và đơn vị trong mọi trường hợp.}
{Kết quả của bài toán không phụ thuộc dữ kiện đã cho.}
{\True Các đồng vị có cùng số proton nhưng khác số neutron; nguyên tử khối trung bình là trung bình có trọng số theo tỉ lệ đồng vị.}
\loigiai{
Các đồng vị có cùng số proton nhưng khác số neutron; nguyên tử khối trung bình là trung bình có trọng số theo tỉ lệ đồng vị. Vì vậy phương án $D$ đúng; các phương án còn lại trái với kiến thức hoặc điều kiện áp dụng.
}
\end{ex}

% ===== Câu 150 =====
% ID: L12C4B21-39-TL
% Mức: VD
\dangbt{Bài toán tổng hợp về cấu trúc và đặc trưng của hạt nhân}
\begin{bt}
% ID: L12C4B21-39-TL
% Mức: VD
Phân tích sai lầm trong nhận định: “Có thể đồng nhất số khối với khối lượng hạt nhân tính theo đơn vị $u$ trong mọi phép tính chính xác.”
\loigiai{
Cần nêu được: Khi giải bài tổng hợp phải xác định đúng $A,Z,N$, đổi đơn vị khối lượng và dùng nhất quán $E=mc^2$ hoặc $1\,u\,c^2=931,5\,MeV$. Khi vận dụng phải xác định đúng dữ kiện, điều kiện áp dụng, hệ đơn vị và kiểm tra tính hợp lí của kết quả. Nhận định “Có thể đồng nhất số khối với khối lượng hạt nhân tính theo đơn vị $u$ trong mọi phép tính chính xác.” sai vì trái với kiến thức trên.
}
\end{bt}

% ===== Câu 151 =====
% ID: L12C4B21-39-TLN
% Mức: TH
\begin{ex}
% ID: L12C4B21-39-TLN
% Mức: TH
Trong bốn phát biểu sau về bài toán tổng hợp về cấu trúc và đặc trưng của hạt nhân, có bao nhiêu phát biểu đúng? Chỉ ghi một số nguyên. (1) Cần đổi các đại lượng về hệ đơn vị thống nhất. (2) Khi giải bài tổng hợp phải xác định đúng $A,Z,N$, đổi đơn vị khối lượng và dùng nhất quán $E=mc^2$ hoặc $1\,u\,c^2=931,5\,MeV$. (3) Có thể đồng nhất số khối với khối lượng hạt nhân tính theo đơn vị $u$ trong mọi phép tính chính xác. (4) Không cần kiểm tra dữ kiện.
\shortans{2}
\loigiai{
Đối chiếu từng phát biểu với kiến thức cốt lõi: Khi giải bài tổng hợp phải xác định đúng $A,Z,N$, đổi đơn vị khối lượng và dùng nhất quán $E=mc^2$ hoặc $1\,u\,c^2=931,5\,MeV$. Có 2 phát biểu đúng.
}
\end{ex}

% ===== Câu 152 =====
% ID: L12C4B21-39-TN
% Mức: VD
\dangbt{Nhận biết đồng vị và tính nguyên tử khối trung bình}
\begin{ex}
% ID: L12C4B21-39-TN
% Mức: VD
Trong một bài thực hành, nhận xét nào đúng về nhận biết đồng vị và tính nguyên tử khối trung bình?
\choice
{\True Các đồng vị có cùng số proton nhưng khác số neutron; nguyên tử khối trung bình là trung bình có trọng số theo tỉ lệ đồng vị.}
{Các đồng vị của một nguyên tố có số proton khác nhau.}
{Có thể bỏ qua điều kiện áp dụng và đơn vị trong mọi trường hợp.}
{Kết quả của bài toán không phụ thuộc dữ kiện đã cho.}
\loigiai{
Các đồng vị có cùng số proton nhưng khác số neutron; nguyên tử khối trung bình là trung bình có trọng số theo tỉ lệ đồng vị. Vì vậy phương án $A$ đúng; các phương án còn lại trái với kiến thức hoặc điều kiện áp dụng.
}
\end{ex}

% ===== Câu 153 =====
% ID: L12C4B21-40-TL
% Mức: VD
\dangbt{Bài toán tổng hợp về cấu trúc và đặc trưng của hạt nhân}
\begin{bt}
% ID: L12C4B21-40-TL
% Mức: VD
Đề xuất các bước giải một bài toán thuộc dạng bài toán tổng hợp về cấu trúc và đặc trưng của hạt nhân.
\loigiai{
Cần nêu được: Khi giải bài tổng hợp phải xác định đúng $A,Z,N$, đổi đơn vị khối lượng và dùng nhất quán $E=mc^2$ hoặc $1\,u\,c^2=931,5\,MeV$. Khi vận dụng phải xác định đúng dữ kiện, điều kiện áp dụng, hệ đơn vị và kiểm tra tính hợp lí của kết quả. Nhận định “Có thể đồng nhất số khối với khối lượng hạt nhân tính theo đơn vị $u$ trong mọi phép tính chính xác.” sai vì trái với kiến thức trên.
}
\end{bt}

% ===== Câu 154 =====
% ID: L12C4B21-40-TLN
% Mức: VD
\begin{ex}
% ID: L12C4B21-40-TLN
% Mức: VD
Trong bốn phát biểu sau về bài toán tổng hợp về cấu trúc và đặc trưng của hạt nhân, có bao nhiêu phát biểu đúng? Chỉ ghi một số nguyên. (1) Khi giải bài tổng hợp phải xác định đúng $A,Z,N$, đổi đơn vị khối lượng và dùng nhất quán $E=mc^2$ hoặc $1\,u\,c^2=931,5\,MeV$. (2) Cần đối chiếu kết quả với điều kiện bài toán. (3) Có thể đồng nhất số khối với khối lượng hạt nhân tính theo đơn vị $u$ trong mọi phép tính chính xác. (4) Mọi trường hợp đều cho cùng một kết quả.
\shortans{2}
\loigiai{
Đối chiếu từng phát biểu với kiến thức cốt lõi: Khi giải bài tổng hợp phải xác định đúng $A,Z,N$, đổi đơn vị khối lượng và dùng nhất quán $E=mc^2$ hoặc $1\,u\,c^2=931,5\,MeV$. Có 2 phát biểu đúng.
}
\end{ex}

% ===== Câu 155 =====
% ID: L12C4B21-40-TN
% Mức: VD
\dangbt{Nhận biết đồng vị và tính nguyên tử khối trung bình}
\begin{ex}
% ID: L12C4B21-40-TN
% Mức: VD
Kết luận đúng khi vận dụng là về nhận biết đồng vị và tính nguyên tử khối trung bình?
\choice
{Các đồng vị của một nguyên tố có số proton khác nhau.}
{\True Các đồng vị có cùng số proton nhưng khác số neutron; nguyên tử khối trung bình là trung bình có trọng số theo tỉ lệ đồng vị.}
{Có thể bỏ qua điều kiện áp dụng và đơn vị trong mọi trường hợp.}
{Kết quả của bài toán không phụ thuộc dữ kiện đã cho.}
\loigiai{
Các đồng vị có cùng số proton nhưng khác số neutron; nguyên tử khối trung bình là trung bình có trọng số theo tỉ lệ đồng vị. Vì vậy phương án $B$ đúng; các phương án còn lại trái với kiến thức hoặc điều kiện áp dụng.
}
\end{ex}

% ===== Câu 156 =====
% ID: L12C4B21-41-TL
% Mức: VD
\dangbt{Bài toán tổng hợp về cấu trúc và đặc trưng của hạt nhân}
\begin{bt}
% ID: L12C4B21-41-TL
% Mức: VD
Nêu một tình huống thực tiễn liên quan đến bài toán tổng hợp về cấu trúc và đặc trưng của hạt nhân và giải thích bằng kiến thức Vật lí.
\loigiai{
Cần nêu được: Khi giải bài tổng hợp phải xác định đúng $A,Z,N$, đổi đơn vị khối lượng và dùng nhất quán $E=mc^2$ hoặc $1\,u\,c^2=931,5\,MeV$. Khi vận dụng phải xác định đúng dữ kiện, điều kiện áp dụng, hệ đơn vị và kiểm tra tính hợp lí của kết quả. Nhận định “Có thể đồng nhất số khối với khối lượng hạt nhân tính theo đơn vị $u$ trong mọi phép tính chính xác.” sai vì trái với kiến thức trên.
}
\end{bt}

% ===== Câu 157 =====
% ID: L12C4B21-41-TLN
% Mức: VD
\begin{ex}
% ID: L12C4B21-41-TLN
% Mức: VD
Trong bốn phát biểu sau về bài toán tổng hợp về cấu trúc và đặc trưng của hạt nhân, có bao nhiêu phát biểu đúng? Chỉ ghi một số nguyên. (1) Có thể đồng nhất số khối với khối lượng hạt nhân tính theo đơn vị $u$ trong mọi phép tính chính xác. (2) Cần đọc đúng dữ kiện và giả thiết. (3) Khi giải bài tổng hợp phải xác định đúng $A,Z,N$, đổi đơn vị khối lượng và dùng nhất quán $E=mc^2$ hoặc $1\,u\,c^2=931,5\,MeV$. (4) Có thể thay số trước khi chọn công thức.
\shortans{2}
\loigiai{
Đối chiếu từng phát biểu với kiến thức cốt lõi: Khi giải bài tổng hợp phải xác định đúng $A,Z,N$, đổi đơn vị khối lượng và dùng nhất quán $E=mc^2$ hoặc $1\,u\,c^2=931,5\,MeV$. Có 2 phát biểu đúng.
}
\end{ex}

% ===== Câu 158 =====
% ID: L12C4B21-41-TN
% Mức: NB
\dangbt{Tính kích thước, bán kính và khối lượng riêng của hạt nhân}
\begin{ex}
% ID: L12C4B21-41-TN
% Mức: NB
Phát biểu nào sau đây đúng về tính kích thước, bán kính và khối lượng riêng của hạt nhân?
\choice
{\True Bán kính hạt nhân gần đúng $R=R_0A^{1/3}$ nên thể tích tỉ lệ với $A$ và khối lượng riêng hạt nhân xấp xỉ không đổi.}
{Bán kính hạt nhân tỉ lệ thuận với số khối $A$.}
{Có thể bỏ qua điều kiện áp dụng và đơn vị trong mọi trường hợp.}
{Kết quả của bài toán không phụ thuộc dữ kiện đã cho.}
\loigiai{
Bán kính hạt nhân gần đúng $R=R_0A^{1/3}$ nên thể tích tỉ lệ với $A$ và khối lượng riêng hạt nhân xấp xỉ không đổi. Vì vậy phương án $A$ đúng; các phương án còn lại trái với kiến thức hoặc điều kiện áp dụng.
}
\end{ex}

% ===== Câu 159 =====
% ID: L12C4B21-42-TL
% Mức: VDC
\dangbt{Bài toán tổng hợp về cấu trúc và đặc trưng của hạt nhân}
\begin{bt}
% ID: L12C4B21-42-TL
% Mức: VDC
So sánh nhận định đúng “Khi giải bài tổng hợp phải xác định đúng $A,Z,N$, đổi đơn vị khối lượng và dùng nhất quán $E=mc^2$ hoặc $1\,u\,c^2=931,5\,MeV$.” với nhận định sai “Có thể đồng nhất số khối với khối lượng hạt nhân tính theo đơn vị $u$ trong mọi phép tính chính xác.”; chỉ rõ căn cứ Vật lí.
\loigiai{
Cần nêu được: Khi giải bài tổng hợp phải xác định đúng $A,Z,N$, đổi đơn vị khối lượng và dùng nhất quán $E=mc^2$ hoặc $1\,u\,c^2=931,5\,MeV$. Khi vận dụng phải xác định đúng dữ kiện, điều kiện áp dụng, hệ đơn vị và kiểm tra tính hợp lí của kết quả. Nhận định “Có thể đồng nhất số khối với khối lượng hạt nhân tính theo đơn vị $u$ trong mọi phép tính chính xác.” sai vì trái với kiến thức trên.
}
\end{bt}

% ===== Câu 160 =====
% ID: L12C4B21-42-TLN
% Mức: VDC
\begin{ex}
% ID: L12C4B21-42-TLN
% Mức: VDC
Trong bốn phát biểu sau về bài toán tổng hợp về cấu trúc và đặc trưng của hạt nhân, có bao nhiêu phát biểu đúng? Chỉ ghi một số nguyên. (1) Kết quả cần có ý nghĩa vật lí. (2) Có thể đồng nhất số khối với khối lượng hạt nhân tính theo đơn vị $u$ trong mọi phép tính chính xác. (3) Phải dùng đúng định luật cho tình huống. (4) Khi giải bài tổng hợp phải xác định đúng $A,Z,N$, đổi đơn vị khối lượng và dùng nhất quán $E=mc^2$ hoặc $1\,u\,c^2=931,5\,MeV$.
\shortans{3}
\loigiai{
Đối chiếu từng phát biểu với kiến thức cốt lõi: Khi giải bài tổng hợp phải xác định đúng $A,Z,N$, đổi đơn vị khối lượng và dùng nhất quán $E=mc^2$ hoặc $1\,u\,c^2=931,5\,MeV$. Có 3 phát biểu đúng.
}
\end{ex}

% ===== Câu 161 =====
% ID: L12C4B21-42-TN
% Mức: NB
\dangbt{Tính kích thước, bán kính và khối lượng riêng của hạt nhân}
\begin{ex}
% ID: L12C4B21-42-TN
% Mức: NB
Nhận định nào phù hợp với kiến thức Vật lí về tính kích thước, bán kính và khối lượng riêng của hạt nhân?
\choice
{Bán kính hạt nhân tỉ lệ thuận với số khối $A$.}
{\True Bán kính hạt nhân gần đúng $R=R_0A^{1/3}$ nên thể tích tỉ lệ với $A$ và khối lượng riêng hạt nhân xấp xỉ không đổi.}
{Có thể bỏ qua điều kiện áp dụng và đơn vị trong mọi trường hợp.}
{Kết quả của bài toán không phụ thuộc dữ kiện đã cho.}
\loigiai{
Bán kính hạt nhân gần đúng $R=R_0A^{1/3}$ nên thể tích tỉ lệ với $A$ và khối lượng riêng hạt nhân xấp xỉ không đổi. Vì vậy phương án $B$ đúng; các phương án còn lại trái với kiến thức hoặc điều kiện áp dụng.
}
\end{ex}

% ===== Câu 162 =====
% ID: L12C4B21-43-TN
% Mức: NB
\begin{ex}
% ID: L12C4B21-43-TN
% Mức: NB
Chọn kết luận chính xác về tính kích thước, bán kính và khối lượng riêng của hạt nhân?
\choice
{Bán kính hạt nhân tỉ lệ thuận với số khối $A$.}
{Có thể bỏ qua điều kiện áp dụng và đơn vị trong mọi trường hợp.}
{\True Bán kính hạt nhân gần đúng $R=R_0A^{1/3}$ nên thể tích tỉ lệ với $A$ và khối lượng riêng hạt nhân xấp xỉ không đổi.}
{Kết quả của bài toán không phụ thuộc dữ kiện đã cho.}
\loigiai{
Bán kính hạt nhân gần đúng $R=R_0A^{1/3}$ nên thể tích tỉ lệ với $A$ và khối lượng riêng hạt nhân xấp xỉ không đổi. Vì vậy phương án $C$ đúng; các phương án còn lại trái với kiến thức hoặc điều kiện áp dụng.
}
\end{ex}

% ===== Câu 163 =====
% ID: L12C4B21-44-TN
% Mức: TH
\begin{ex}
% ID: L12C4B21-44-TN
% Mức: TH
Nội dung nào mô tả đúng nhất về tính kích thước, bán kính và khối lượng riêng của hạt nhân?
\choice
{Bán kính hạt nhân tỉ lệ thuận với số khối $A$.}
{Có thể bỏ qua điều kiện áp dụng và đơn vị trong mọi trường hợp.}
{Kết quả của bài toán không phụ thuộc dữ kiện đã cho.}
{\True Bán kính hạt nhân gần đúng $R=R_0A^{1/3}$ nên thể tích tỉ lệ với $A$ và khối lượng riêng hạt nhân xấp xỉ không đổi.}
\loigiai{
Bán kính hạt nhân gần đúng $R=R_0A^{1/3}$ nên thể tích tỉ lệ với $A$ và khối lượng riêng hạt nhân xấp xỉ không đổi. Vì vậy phương án $D$ đúng; các phương án còn lại trái với kiến thức hoặc điều kiện áp dụng.
}
\end{ex}

% ===== Câu 164 =====
% ID: L12C4B21-45-TN
% Mức: TH
\begin{ex}
% ID: L12C4B21-45-TN
% Mức: TH
Khi phân tích, phát biểu nào đúng về tính kích thước, bán kính và khối lượng riêng của hạt nhân?
\choice
{\True Bán kính hạt nhân gần đúng $R=R_0A^{1/3}$ nên thể tích tỉ lệ với $A$ và khối lượng riêng hạt nhân xấp xỉ không đổi.}
{Bán kính hạt nhân tỉ lệ thuận với số khối $A$.}
{Có thể bỏ qua điều kiện áp dụng và đơn vị trong mọi trường hợp.}
{Kết quả của bài toán không phụ thuộc dữ kiện đã cho.}
\loigiai{
Bán kính hạt nhân gần đúng $R=R_0A^{1/3}$ nên thể tích tỉ lệ với $A$ và khối lượng riêng hạt nhân xấp xỉ không đổi. Vì vậy phương án $A$ đúng; các phương án còn lại trái với kiến thức hoặc điều kiện áp dụng.
}
\end{ex}

% ===== Câu 165 =====
% ID: L12C4B21-46-TN
% Mức: TH
\begin{ex}
% ID: L12C4B21-46-TN
% Mức: TH
Kết luận nào của học sinh là đúng về tính kích thước, bán kính và khối lượng riêng của hạt nhân?
\choice
{Bán kính hạt nhân tỉ lệ thuận với số khối $A$.}
{\True Bán kính hạt nhân gần đúng $R=R_0A^{1/3}$ nên thể tích tỉ lệ với $A$ và khối lượng riêng hạt nhân xấp xỉ không đổi.}
{Có thể bỏ qua điều kiện áp dụng và đơn vị trong mọi trường hợp.}
{Kết quả của bài toán không phụ thuộc dữ kiện đã cho.}
\loigiai{
Bán kính hạt nhân gần đúng $R=R_0A^{1/3}$ nên thể tích tỉ lệ với $A$ và khối lượng riêng hạt nhân xấp xỉ không đổi. Vì vậy phương án $B$ đúng; các phương án còn lại trái với kiến thức hoặc điều kiện áp dụng.
}
\end{ex}

% ===== Câu 166 =====
% ID: L12C4B21-47-TN
% Mức: TH
\begin{ex}
% ID: L12C4B21-47-TN
% Mức: TH
Chọn phát biểu không trái với kiến thức về tính kích thước, bán kính và khối lượng riêng của hạt nhân?
\choice
{Bán kính hạt nhân tỉ lệ thuận với số khối $A$.}
{Có thể bỏ qua điều kiện áp dụng và đơn vị trong mọi trường hợp.}
{\True Bán kính hạt nhân gần đúng $R=R_0A^{1/3}$ nên thể tích tỉ lệ với $A$ và khối lượng riêng hạt nhân xấp xỉ không đổi.}
{Kết quả của bài toán không phụ thuộc dữ kiện đã cho.}
\loigiai{
Bán kính hạt nhân gần đúng $R=R_0A^{1/3}$ nên thể tích tỉ lệ với $A$ và khối lượng riêng hạt nhân xấp xỉ không đổi. Vì vậy phương án $C$ đúng; các phương án còn lại trái với kiến thức hoặc điều kiện áp dụng.
}
\end{ex}

% ===== Câu 167 =====
% ID: L12C4B21-48-TN
% Mức: VD
\begin{ex}
% ID: L12C4B21-48-TN
% Mức: VD
Cơ sở Vật lí đúng để giải bài là gì về tính kích thước, bán kính và khối lượng riêng của hạt nhân?
\choice
{Bán kính hạt nhân tỉ lệ thuận với số khối $A$.}
{Có thể bỏ qua điều kiện áp dụng và đơn vị trong mọi trường hợp.}
{Kết quả của bài toán không phụ thuộc dữ kiện đã cho.}
{\True Bán kính hạt nhân gần đúng $R=R_0A^{1/3}$ nên thể tích tỉ lệ với $A$ và khối lượng riêng hạt nhân xấp xỉ không đổi.}
\loigiai{
Bán kính hạt nhân gần đúng $R=R_0A^{1/3}$ nên thể tích tỉ lệ với $A$ và khối lượng riêng hạt nhân xấp xỉ không đổi. Vì vậy phương án $D$ đúng; các phương án còn lại trái với kiến thức hoặc điều kiện áp dụng.
}
\end{ex}

% ===== Câu 168 =====
% ID: L12C4B21-49-TN
% Mức: VD
\begin{ex}
% ID: L12C4B21-49-TN
% Mức: VD
Trong một bài thực hành, nhận xét nào đúng về tính kích thước, bán kính và khối lượng riêng của hạt nhân?
\choice
{\True Bán kính hạt nhân gần đúng $R=R_0A^{1/3}$ nên thể tích tỉ lệ với $A$ và khối lượng riêng hạt nhân xấp xỉ không đổi.}
{Bán kính hạt nhân tỉ lệ thuận với số khối $A$.}
{Có thể bỏ qua điều kiện áp dụng và đơn vị trong mọi trường hợp.}
{Kết quả của bài toán không phụ thuộc dữ kiện đã cho.}
\loigiai{
Bán kính hạt nhân gần đúng $R=R_0A^{1/3}$ nên thể tích tỉ lệ với $A$ và khối lượng riêng hạt nhân xấp xỉ không đổi. Vì vậy phương án $A$ đúng; các phương án còn lại trái với kiến thức hoặc điều kiện áp dụng.
}
\end{ex}

% ===== Câu 169 =====
% ID: L12C4B21-50-TN
% Mức: VD
\begin{ex}
% ID: L12C4B21-50-TN
% Mức: VD
Kết luận đúng khi vận dụng là về tính kích thước, bán kính và khối lượng riêng của hạt nhân?
\choice
{Bán kính hạt nhân tỉ lệ thuận với số khối $A$.}
{\True Bán kính hạt nhân gần đúng $R=R_0A^{1/3}$ nên thể tích tỉ lệ với $A$ và khối lượng riêng hạt nhân xấp xỉ không đổi.}
{Có thể bỏ qua điều kiện áp dụng và đơn vị trong mọi trường hợp.}
{Kết quả của bài toán không phụ thuộc dữ kiện đã cho.}
\loigiai{
Bán kính hạt nhân gần đúng $R=R_0A^{1/3}$ nên thể tích tỉ lệ với $A$ và khối lượng riêng hạt nhân xấp xỉ không đổi. Vì vậy phương án $B$ đúng; các phương án còn lại trái với kiến thức hoặc điều kiện áp dụng.
}
\end{ex}

% ===== Câu 170 =====
% ID: L12C4B21-51-TN
% Mức: NB
\dangbt{Vận dụng hệ thức khối lượng – năng lượng và độ hụt khối}
\begin{ex}
% ID: L12C4B21-51-TN
% Mức: NB
Phát biểu nào sau đây đúng về vận dụng hệ thức khối lượng – năng lượng và độ hụt khối?
\choice
{\True Khối lượng và năng lượng liên hệ bởi $E=mc^2$; độ hụt khối là phần chênh giữa tổng khối lượng nuclon riêng rẽ và khối lượng hạt nhân.}
{Độ hụt khối của hạt nhân bền luôn âm.}
{Có thể bỏ qua điều kiện áp dụng và đơn vị trong mọi trường hợp.}
{Kết quả của bài toán không phụ thuộc dữ kiện đã cho.}
\loigiai{
Khối lượng và năng lượng liên hệ bởi $E=mc^2$; độ hụt khối là phần chênh giữa tổng khối lượng nuclon riêng rẽ và khối lượng hạt nhân. Vì vậy phương án $A$ đúng; các phương án còn lại trái với kiến thức hoặc điều kiện áp dụng.
}
\end{ex}

% ===== Câu 171 =====
% ID: L12C4B21-52-TN
% Mức: NB
\begin{ex}
% ID: L12C4B21-52-TN
% Mức: NB
Nhận định nào phù hợp với kiến thức Vật lí về vận dụng hệ thức khối lượng – năng lượng và độ hụt khối?
\choice
{Độ hụt khối của hạt nhân bền luôn âm.}
{\True Khối lượng và năng lượng liên hệ bởi $E=mc^2$; độ hụt khối là phần chênh giữa tổng khối lượng nuclon riêng rẽ và khối lượng hạt nhân.}
{Có thể bỏ qua điều kiện áp dụng và đơn vị trong mọi trường hợp.}
{Kết quả của bài toán không phụ thuộc dữ kiện đã cho.}
\loigiai{
Khối lượng và năng lượng liên hệ bởi $E=mc^2$; độ hụt khối là phần chênh giữa tổng khối lượng nuclon riêng rẽ và khối lượng hạt nhân. Vì vậy phương án $B$ đúng; các phương án còn lại trái với kiến thức hoặc điều kiện áp dụng.
}
\end{ex}

% ===== Câu 172 =====
% ID: L12C4B21-53-TN
% Mức: NB
\begin{ex}
% ID: L12C4B21-53-TN
% Mức: NB
Chọn kết luận chính xác về vận dụng hệ thức khối lượng – năng lượng và độ hụt khối?
\choice
{Độ hụt khối của hạt nhân bền luôn âm.}
{Có thể bỏ qua điều kiện áp dụng và đơn vị trong mọi trường hợp.}
{\True Khối lượng và năng lượng liên hệ bởi $E=mc^2$; độ hụt khối là phần chênh giữa tổng khối lượng nuclon riêng rẽ và khối lượng hạt nhân.}
{Kết quả của bài toán không phụ thuộc dữ kiện đã cho.}
\loigiai{
Khối lượng và năng lượng liên hệ bởi $E=mc^2$; độ hụt khối là phần chênh giữa tổng khối lượng nuclon riêng rẽ và khối lượng hạt nhân. Vì vậy phương án $C$ đúng; các phương án còn lại trái với kiến thức hoặc điều kiện áp dụng.
}
\end{ex}

% ===== Câu 173 =====
% ID: L12C4B21-54-TN
% Mức: TH
\begin{ex}
% ID: L12C4B21-54-TN
% Mức: TH
Nội dung nào mô tả đúng nhất về vận dụng hệ thức khối lượng – năng lượng và độ hụt khối?
\choice
{Độ hụt khối của hạt nhân bền luôn âm.}
{Có thể bỏ qua điều kiện áp dụng và đơn vị trong mọi trường hợp.}
{Kết quả của bài toán không phụ thuộc dữ kiện đã cho.}
{\True Khối lượng và năng lượng liên hệ bởi $E=mc^2$; độ hụt khối là phần chênh giữa tổng khối lượng nuclon riêng rẽ và khối lượng hạt nhân.}
\loigiai{
Khối lượng và năng lượng liên hệ bởi $E=mc^2$; độ hụt khối là phần chênh giữa tổng khối lượng nuclon riêng rẽ và khối lượng hạt nhân. Vì vậy phương án $D$ đúng; các phương án còn lại trái với kiến thức hoặc điều kiện áp dụng.
}
\end{ex}

% ===== Câu 174 =====
% ID: L12C4B21-55-TN
% Mức: TH
\begin{ex}
% ID: L12C4B21-55-TN
% Mức: TH
Khi phân tích, phát biểu nào đúng về vận dụng hệ thức khối lượng – năng lượng và độ hụt khối?
\choice
{\True Khối lượng và năng lượng liên hệ bởi $E=mc^2$; độ hụt khối là phần chênh giữa tổng khối lượng nuclon riêng rẽ và khối lượng hạt nhân.}
{Độ hụt khối của hạt nhân bền luôn âm.}
{Có thể bỏ qua điều kiện áp dụng và đơn vị trong mọi trường hợp.}
{Kết quả của bài toán không phụ thuộc dữ kiện đã cho.}
\loigiai{
Khối lượng và năng lượng liên hệ bởi $E=mc^2$; độ hụt khối là phần chênh giữa tổng khối lượng nuclon riêng rẽ và khối lượng hạt nhân. Vì vậy phương án $A$ đúng; các phương án còn lại trái với kiến thức hoặc điều kiện áp dụng.
}
\end{ex}

% ===== Câu 175 =====
% ID: L12C4B21-56-TN
% Mức: TH
\begin{ex}
% ID: L12C4B21-56-TN
% Mức: TH
Kết luận nào của học sinh là đúng về vận dụng hệ thức khối lượng – năng lượng và độ hụt khối?
\choice
{Độ hụt khối của hạt nhân bền luôn âm.}
{\True Khối lượng và năng lượng liên hệ bởi $E=mc^2$; độ hụt khối là phần chênh giữa tổng khối lượng nuclon riêng rẽ và khối lượng hạt nhân.}
{Có thể bỏ qua điều kiện áp dụng và đơn vị trong mọi trường hợp.}
{Kết quả của bài toán không phụ thuộc dữ kiện đã cho.}
\loigiai{
Khối lượng và năng lượng liên hệ bởi $E=mc^2$; độ hụt khối là phần chênh giữa tổng khối lượng nuclon riêng rẽ và khối lượng hạt nhân. Vì vậy phương án $B$ đúng; các phương án còn lại trái với kiến thức hoặc điều kiện áp dụng.
}
\end{ex}

% ===== Câu 176 =====
% ID: L12C4B21-57-TN
% Mức: TH
\begin{ex}
% ID: L12C4B21-57-TN
% Mức: TH
Chọn phát biểu không trái với kiến thức về vận dụng hệ thức khối lượng – năng lượng và độ hụt khối?
\choice
{Độ hụt khối của hạt nhân bền luôn âm.}
{Có thể bỏ qua điều kiện áp dụng và đơn vị trong mọi trường hợp.}
{\True Khối lượng và năng lượng liên hệ bởi $E=mc^2$; độ hụt khối là phần chênh giữa tổng khối lượng nuclon riêng rẽ và khối lượng hạt nhân.}
{Kết quả của bài toán không phụ thuộc dữ kiện đã cho.}
\loigiai{
Khối lượng và năng lượng liên hệ bởi $E=mc^2$; độ hụt khối là phần chênh giữa tổng khối lượng nuclon riêng rẽ và khối lượng hạt nhân. Vì vậy phương án $C$ đúng; các phương án còn lại trái với kiến thức hoặc điều kiện áp dụng.
}
\end{ex}

% ===== Câu 177 =====
% ID: L12C4B21-58-TN
% Mức: VD
\begin{ex}
% ID: L12C4B21-58-TN
% Mức: VD
Cơ sở Vật lí đúng để giải bài là gì về vận dụng hệ thức khối lượng – năng lượng và độ hụt khối?
\choice
{Độ hụt khối của hạt nhân bền luôn âm.}
{Có thể bỏ qua điều kiện áp dụng và đơn vị trong mọi trường hợp.}
{Kết quả của bài toán không phụ thuộc dữ kiện đã cho.}
{\True Khối lượng và năng lượng liên hệ bởi $E=mc^2$; độ hụt khối là phần chênh giữa tổng khối lượng nuclon riêng rẽ và khối lượng hạt nhân.}
\loigiai{
Khối lượng và năng lượng liên hệ bởi $E=mc^2$; độ hụt khối là phần chênh giữa tổng khối lượng nuclon riêng rẽ và khối lượng hạt nhân. Vì vậy phương án $D$ đúng; các phương án còn lại trái với kiến thức hoặc điều kiện áp dụng.
}
\end{ex}

% ===== Câu 178 =====
% ID: L12C4B21-59-TN
% Mức: VD
\begin{ex}
% ID: L12C4B21-59-TN
% Mức: VD
Trong một bài thực hành, nhận xét nào đúng về vận dụng hệ thức khối lượng – năng lượng và độ hụt khối?
\choice
{\True Khối lượng và năng lượng liên hệ bởi $E=mc^2$; độ hụt khối là phần chênh giữa tổng khối lượng nuclon riêng rẽ và khối lượng hạt nhân.}
{Độ hụt khối của hạt nhân bền luôn âm.}
{Có thể bỏ qua điều kiện áp dụng và đơn vị trong mọi trường hợp.}
{Kết quả của bài toán không phụ thuộc dữ kiện đã cho.}
\loigiai{
Khối lượng và năng lượng liên hệ bởi $E=mc^2$; độ hụt khối là phần chênh giữa tổng khối lượng nuclon riêng rẽ và khối lượng hạt nhân. Vì vậy phương án $A$ đúng; các phương án còn lại trái với kiến thức hoặc điều kiện áp dụng.
}
\end{ex}

% ===== Câu 179 =====
% ID: L12C4B21-60-TN
% Mức: VD
\begin{ex}
% ID: L12C4B21-60-TN
% Mức: VD
Kết luận đúng khi vận dụng là về vận dụng hệ thức khối lượng – năng lượng và độ hụt khối?
\choice
{Độ hụt khối của hạt nhân bền luôn âm.}
{\True Khối lượng và năng lượng liên hệ bởi $E=mc^2$; độ hụt khối là phần chênh giữa tổng khối lượng nuclon riêng rẽ và khối lượng hạt nhân.}
{Có thể bỏ qua điều kiện áp dụng và đơn vị trong mọi trường hợp.}
{Kết quả của bài toán không phụ thuộc dữ kiện đã cho.}
\loigiai{
Khối lượng và năng lượng liên hệ bởi $E=mc^2$; độ hụt khối là phần chênh giữa tổng khối lượng nuclon riêng rẽ và khối lượng hạt nhân. Vì vậy phương án $B$ đúng; các phương án còn lại trái với kiến thức hoặc điều kiện áp dụng.
}
\end{ex}

% ===== Câu 180 =====
% ID: L12C4B21-61-TN
% Mức: NB
\dangbt{Bài toán tổng hợp về cấu trúc và đặc trưng của hạt nhân}
\begin{ex}
% ID: L12C4B21-61-TN
% Mức: NB
Phát biểu nào sau đây đúng về bài toán tổng hợp về cấu trúc và đặc trưng của hạt nhân?
\choice
{\True Khi giải bài tổng hợp phải xác định đúng $A,Z,N$, đổi đơn vị khối lượng và dùng nhất quán $E=mc^2$ hoặc $1\,u\,c^2=931,5\,MeV$.}
{Có thể đồng nhất số khối với khối lượng hạt nhân tính theo đơn vị $u$ trong mọi phép tính chính xác.}
{Có thể bỏ qua điều kiện áp dụng và đơn vị trong mọi trường hợp.}
{Kết quả của bài toán không phụ thuộc dữ kiện đã cho.}
\loigiai{
Khi giải bài tổng hợp phải xác định đúng $A,Z,N$, đổi đơn vị khối lượng và dùng nhất quán $E=mc^2$ hoặc $1\,u\,c^2=931,5\,MeV$. Vì vậy phương án $A$ đúng; các phương án còn lại trái với kiến thức hoặc điều kiện áp dụng.
}
\end{ex}

% ===== Câu 181 =====
% ID: L12C4B21-62-TN
% Mức: NB
\begin{ex}
% ID: L12C4B21-62-TN
% Mức: NB
Nhận định nào phù hợp với kiến thức Vật lí về bài toán tổng hợp về cấu trúc và đặc trưng của hạt nhân?
\choice
{Có thể đồng nhất số khối với khối lượng hạt nhân tính theo đơn vị $u$ trong mọi phép tính chính xác.}
{\True Khi giải bài tổng hợp phải xác định đúng $A,Z,N$, đổi đơn vị khối lượng và dùng nhất quán $E=mc^2$ hoặc $1\,u\,c^2=931,5\,MeV$.}
{Có thể bỏ qua điều kiện áp dụng và đơn vị trong mọi trường hợp.}
{Kết quả của bài toán không phụ thuộc dữ kiện đã cho.}
\loigiai{
Khi giải bài tổng hợp phải xác định đúng $A,Z,N$, đổi đơn vị khối lượng và dùng nhất quán $E=mc^2$ hoặc $1\,u\,c^2=931,5\,MeV$. Vì vậy phương án $B$ đúng; các phương án còn lại trái với kiến thức hoặc điều kiện áp dụng.
}
\end{ex}

% ===== Câu 182 =====
% ID: L12C4B21-63-TN
% Mức: NB
\begin{ex}
% ID: L12C4B21-63-TN
% Mức: NB
Chọn kết luận chính xác về bài toán tổng hợp về cấu trúc và đặc trưng của hạt nhân?
\choice
{Có thể đồng nhất số khối với khối lượng hạt nhân tính theo đơn vị $u$ trong mọi phép tính chính xác.}
{Có thể bỏ qua điều kiện áp dụng và đơn vị trong mọi trường hợp.}
{\True Khi giải bài tổng hợp phải xác định đúng $A,Z,N$, đổi đơn vị khối lượng và dùng nhất quán $E=mc^2$ hoặc $1\,u\,c^2=931,5\,MeV$.}
{Kết quả của bài toán không phụ thuộc dữ kiện đã cho.}
\loigiai{
Khi giải bài tổng hợp phải xác định đúng $A,Z,N$, đổi đơn vị khối lượng và dùng nhất quán $E=mc^2$ hoặc $1\,u\,c^2=931,5\,MeV$. Vì vậy phương án $C$ đúng; các phương án còn lại trái với kiến thức hoặc điều kiện áp dụng.
}
\end{ex}

% ===== Câu 183 =====
% ID: L12C4B21-64-TN
% Mức: TH
\begin{ex}
% ID: L12C4B21-64-TN
% Mức: TH
Nội dung nào mô tả đúng nhất về bài toán tổng hợp về cấu trúc và đặc trưng của hạt nhân?
\choice
{Có thể đồng nhất số khối với khối lượng hạt nhân tính theo đơn vị $u$ trong mọi phép tính chính xác.}
{Có thể bỏ qua điều kiện áp dụng và đơn vị trong mọi trường hợp.}
{Kết quả của bài toán không phụ thuộc dữ kiện đã cho.}
{\True Khi giải bài tổng hợp phải xác định đúng $A,Z,N$, đổi đơn vị khối lượng và dùng nhất quán $E=mc^2$ hoặc $1\,u\,c^2=931,5\,MeV$.}
\loigiai{
Khi giải bài tổng hợp phải xác định đúng $A,Z,N$, đổi đơn vị khối lượng và dùng nhất quán $E=mc^2$ hoặc $1\,u\,c^2=931,5\,MeV$. Vì vậy phương án $D$ đúng; các phương án còn lại trái với kiến thức hoặc điều kiện áp dụng.
}
\end{ex}

% ===== Câu 184 =====
% ID: L12C4B21-65-TN
% Mức: TH
\begin{ex}
% ID: L12C4B21-65-TN
% Mức: TH
Khi phân tích, phát biểu nào đúng về bài toán tổng hợp về cấu trúc và đặc trưng của hạt nhân?
\choice
{\True Khi giải bài tổng hợp phải xác định đúng $A,Z,N$, đổi đơn vị khối lượng và dùng nhất quán $E=mc^2$ hoặc $1\,u\,c^2=931,5\,MeV$.}
{Có thể đồng nhất số khối với khối lượng hạt nhân tính theo đơn vị $u$ trong mọi phép tính chính xác.}
{Có thể bỏ qua điều kiện áp dụng và đơn vị trong mọi trường hợp.}
{Kết quả của bài toán không phụ thuộc dữ kiện đã cho.}
\loigiai{
Khi giải bài tổng hợp phải xác định đúng $A,Z,N$, đổi đơn vị khối lượng và dùng nhất quán $E=mc^2$ hoặc $1\,u\,c^2=931,5\,MeV$. Vì vậy phương án $A$ đúng; các phương án còn lại trái với kiến thức hoặc điều kiện áp dụng.
}
\end{ex}

% ===== Câu 185 =====
% ID: L12C4B21-66-TN
% Mức: TH
\begin{ex}
% ID: L12C4B21-66-TN
% Mức: TH
Kết luận nào của học sinh là đúng về bài toán tổng hợp về cấu trúc và đặc trưng của hạt nhân?
\choice
{Có thể đồng nhất số khối với khối lượng hạt nhân tính theo đơn vị $u$ trong mọi phép tính chính xác.}
{\True Khi giải bài tổng hợp phải xác định đúng $A,Z,N$, đổi đơn vị khối lượng và dùng nhất quán $E=mc^2$ hoặc $1\,u\,c^2=931,5\,MeV$.}
{Có thể bỏ qua điều kiện áp dụng và đơn vị trong mọi trường hợp.}
{Kết quả của bài toán không phụ thuộc dữ kiện đã cho.}
\loigiai{
Khi giải bài tổng hợp phải xác định đúng $A,Z,N$, đổi đơn vị khối lượng và dùng nhất quán $E=mc^2$ hoặc $1\,u\,c^2=931,5\,MeV$. Vì vậy phương án $B$ đúng; các phương án còn lại trái với kiến thức hoặc điều kiện áp dụng.
}
\end{ex}

% ===== Câu 186 =====
% ID: L12C4B21-67-TN
% Mức: TH
\begin{ex}
% ID: L12C4B21-67-TN
% Mức: TH
Chọn phát biểu không trái với kiến thức về bài toán tổng hợp về cấu trúc và đặc trưng của hạt nhân?
\choice
{Có thể đồng nhất số khối với khối lượng hạt nhân tính theo đơn vị $u$ trong mọi phép tính chính xác.}
{Có thể bỏ qua điều kiện áp dụng và đơn vị trong mọi trường hợp.}
{\True Khi giải bài tổng hợp phải xác định đúng $A,Z,N$, đổi đơn vị khối lượng và dùng nhất quán $E=mc^2$ hoặc $1\,u\,c^2=931,5\,MeV$.}
{Kết quả của bài toán không phụ thuộc dữ kiện đã cho.}
\loigiai{
Khi giải bài tổng hợp phải xác định đúng $A,Z,N$, đổi đơn vị khối lượng và dùng nhất quán $E=mc^2$ hoặc $1\,u\,c^2=931,5\,MeV$. Vì vậy phương án $C$ đúng; các phương án còn lại trái với kiến thức hoặc điều kiện áp dụng.
}
\end{ex}

% ===== Câu 187 =====
% ID: L12C4B21-68-TN
% Mức: VD
\begin{ex}
% ID: L12C4B21-68-TN
% Mức: VD
Cơ sở Vật lí đúng để giải bài là gì về bài toán tổng hợp về cấu trúc và đặc trưng của hạt nhân?
\choice
{Có thể đồng nhất số khối với khối lượng hạt nhân tính theo đơn vị $u$ trong mọi phép tính chính xác.}
{Có thể bỏ qua điều kiện áp dụng và đơn vị trong mọi trường hợp.}
{Kết quả của bài toán không phụ thuộc dữ kiện đã cho.}
{\True Khi giải bài tổng hợp phải xác định đúng $A,Z,N$, đổi đơn vị khối lượng và dùng nhất quán $E=mc^2$ hoặc $1\,u\,c^2=931,5\,MeV$.}
\loigiai{
Khi giải bài tổng hợp phải xác định đúng $A,Z,N$, đổi đơn vị khối lượng và dùng nhất quán $E=mc^2$ hoặc $1\,u\,c^2=931,5\,MeV$. Vì vậy phương án $D$ đúng; các phương án còn lại trái với kiến thức hoặc điều kiện áp dụng.
}
\end{ex}

% ===== Câu 188 =====
% ID: L12C4B21-69-TN
% Mức: VD
\begin{ex}
% ID: L12C4B21-69-TN
% Mức: VD
Trong một bài thực hành, nhận xét nào đúng về bài toán tổng hợp về cấu trúc và đặc trưng của hạt nhân?
\choice
{\True Khi giải bài tổng hợp phải xác định đúng $A,Z,N$, đổi đơn vị khối lượng và dùng nhất quán $E=mc^2$ hoặc $1\,u\,c^2=931,5\,MeV$.}
{Có thể đồng nhất số khối với khối lượng hạt nhân tính theo đơn vị $u$ trong mọi phép tính chính xác.}
{Có thể bỏ qua điều kiện áp dụng và đơn vị trong mọi trường hợp.}
{Kết quả của bài toán không phụ thuộc dữ kiện đã cho.}
\loigiai{
Khi giải bài tổng hợp phải xác định đúng $A,Z,N$, đổi đơn vị khối lượng và dùng nhất quán $E=mc^2$ hoặc $1\,u\,c^2=931,5\,MeV$. Vì vậy phương án $A$ đúng; các phương án còn lại trái với kiến thức hoặc điều kiện áp dụng.
}
\end{ex}

% ===== Câu 189 =====
% ID: L12C4B21-70-TN
% Mức: VD
\begin{ex}
% ID: L12C4B21-70-TN
% Mức: VD
Kết luận đúng khi vận dụng là về bài toán tổng hợp về cấu trúc và đặc trưng của hạt nhân?
\choice
{Có thể đồng nhất số khối với khối lượng hạt nhân tính theo đơn vị $u$ trong mọi phép tính chính xác.}
{\True Khi giải bài tổng hợp phải xác định đúng $A,Z,N$, đổi đơn vị khối lượng và dùng nhất quán $E=mc^2$ hoặc $1\,u\,c^2=931,5\,MeV$.}
{Có thể bỏ qua điều kiện áp dụng và đơn vị trong mọi trường hợp.}
{Kết quả của bài toán không phụ thuộc dữ kiện đã cho.}
\loigiai{
Khi giải bài tổng hợp phải xác định đúng $A,Z,N$, đổi đơn vị khối lượng và dùng nhất quán $E=mc^2$ hoặc $1\,u\,c^2=931,5\,MeV$. Vì vậy phương án $B$ đúng; các phương án còn lại trái với kiến thức hoặc điều kiện áp dụng.
}
\end{ex}
